\begin{savequote}
{\tiny Ein Mann ging von Jerusalem nach Jericho hinab und wurde von
Räubern überfallen. Sie plünderten ihn aus und schlugen ihn nieder; dann gingen sie weg und ließen ihn 
halbtot liegen. Zufällig kam ein Priester denselben Weg
herab; er sah ihn und ging weiter. Auch ein Levit kam zu
der  Stelle; er sah ihn und ging weiter. Dann kam ein Mann
aus  Samarien, der auf der Reise war. Als er ihn sah, hatte
er  Mitleid, ging zu ihm hin, goss Öl und Wein auf seine
Wunden  und verband sie. 

Dann hob er ihn auf sein Reittier,
brachte  ihn zu einer Herberge und sorgte für ihn. Am
anderen Morgen  holte er zwei Denare hervor, gab sie dem
Wirt und sagte:  Sorge für ihn, und wenn du mehr für ihn
brauchst, werde  ich es dir bezahlen, wenn ich wiederkomme.
Was meinst du:  Wer von diesen dreien hat sich als der
Nächste dessen erwiesen,  der von den Räubern überfallen
wurde?  

Der Gesetzeslehrer antwortete: Der, der
barmherzig\index{Barmherzigkeit} an  ihm gehandelt hat. Da
sagte Jesus zu ihm: Dann geh und handle genauso!

Lk 10,25-37}
% \qauthor{{\tiny Lk 10,25-37}}
\end{savequote}

\chapter{[MED] Klassische Medizinische Herausforderungen ---}
\label{chp:MED}
\gAasRsN{RS VI}

\emph{Maintainer: Sebastian Gabriel}
\fxnote[author=GABS]{test}

\minitoc


% \begin{verbatim}
% 
% Changelog:
% 
% 2009-03-14: In LaTeX überführt
% 
% 2009-02-25: Start Changelog. Bis inkl. bakt. Meningitis mit Korr. von
% Koch verglichen und korr.
% \end{verbatim}




\clearpage
\section{Herz-Kreislaufstörungen}

\dictum[Goethe]{Blut ist ein ganz besond'rer Saft.}

\label{sec:herz-kreislauf-stoerungen}
\label{topic:herz-kreislauf-stoerungen}

% \paragraph{Überblick über die wichtigsten
% Herz-Kreislauf-Notfälle}
% 
% \begin{itemize}
%     \item Herz-/Kreislaufstillstand
%     \item Herzinsuffizienz
%     \begin{itemize}
%         \item Linksherz-Versagen
%         \item Rechtsherz-Versagen
%     \end{itemize}
%     \item (Akutes) Koronarsyndrom
%     \begin{itemize}
%         \item Angina pectoris
%         \item Myokardinfarkt (MCI )
%     \end{itemize}
%     \item Herzrhythmusstörungen
%     \item Störungen des Bludrucks
%     \begin{itemize}
%         \item Hyp\gEs{o}tonie
%         \item Hyp\gEs{er}tonie
%         \begin{itemize}
%             \item Hypertensive Krise
%             \item Hypertensiver Notfall
%         \end{itemize}
%     \end{itemize}
% \end{itemize}


\subsection{Herz-Kreislauf-Stillstand}

Der Herz-Kreislauf-Stillstand wird
unter
\prettyref{topic:kreislaufstillstand}
und \prettyref{topic:reanimation-eh} besprochen.



\subsection{Herzinsuffizienz}
\label{topic:herzinsuffizienz}
\index{Herzinsuffizienz}
\index{Herzinsuffizienz!kompensierte}
\index{Herzinsuffizienz!dekompensierte}
\index{Rechtsherzinsuffizienz}
\index{Linksherzinsuffizienz}
\index{Globalinsuffizienz}
\index{Insuffizienz!Herz-}


\gDefinition{Unzureichende Auswurfleistung des Herzens.}

Leistungseinschränkung des Herzmuskels führt zu einer verminderten
Auswurfleistung, viele Ursachen sind möglich. Manche Ursachen treten
\gE{akut} auf, manche sind eher \gE{chronischer} Natur.



\gLayout{0}{Einteilung}
{
\subparagraph{Herzteil} 
Je nachdem welcher Herzteil betroffen ist unterteilt man in:

\begin{compactdesc}
	\item[\gT{Rechtsherzinsuffizienz}] Unzureichende Leistung des \gE{rechten}
	Herzens,
	\item[\gT{Linksherzinsuffizienz}]
	 Unzureichende Leistung des \gE{linken} Herzens, und
	\item[\gT{Globalinsuffizienz}] Unzureichende Leistung beider Herzteile.  
\end{compactdesc}

\subparagraph{Kompensation}
Abhängig davon, wie gut der Körper mit der
Einschränkung zurecht kommt unterteilt man:

\begin{compactdesc}
    \item[\gT{Kompensierte Herzinsuffizienz}] Der Körper hat noch
    genug Reserven, kann gegensteuern und kommt mit der
    eingeschränkten Leistung aus. 
    \item[\gT{\gE{De}kompensierte Herzinsuffizienz}] Der Körper kommt aufgrund
    des Fortschreitens der Erkrankung oder aufgrund eines erhöhten
    Bedarfs nicht mehr mit der Leistung aus. Es treten deutliche
    Symptome auf.
\end{compactdesc}
}{
\begin{itemize}
    \item Herzteil
    \begin{itemize}
        \item Rechtsherzinsuffizienz
        \item Linksherzinsuffizienz
        \item Globalinsuffizienz
    \end{itemize}
    \item Kompensation
    \begin{itemize}
        \item Kompensierte Herzinsuffizienz
        \item \gE{De}kompensierte Herzinsuffizienz
    \end{itemize}
\end{itemize}
}{}{}{}




\gLayout{0}{Beurteilung des Schweregrades}
{
Es ist wichtig den Schweregrad einer Herzinsuffizienz zu beurteilen. Sie
muss nicht dramatisch sein, für viele Patienten ist sie ein
chronischer Zustand mit dem sie gut leben können, da sich
der Körper mit der Zeit an die abnehmende Leistungsfähigkeit
gewöhnt.

Problematisch wird es jedoch, wenn der Patient plötzlich
eine Symptomatik entwickelt und die Kompensationsmechanismen (köpereigene, Medikamente, {\dots})
versagen. Man spricht dann von der \gEs{dekompensierten}
Herzinsuffizienz. Man muss die Symptome generell \gE{anhand ihres
Verlaufes beurteilen}: Bekannte, seit längerer Zeit bestehende
Symptome sind weniger wichtig als plötzlich auftretende
Symptome. Plötzlich auftretende Symptome sind oft \gEs{Alarmzeichen}. 
}{
\begin{itemize}
    \item Chronisch: oft Gewöhnungseffekte, Kompensation.
    \item Kompensation versagt: Plötzlich neu auftretende
    Symptome, Beschwerden, Alarmzeichen.
\end{itemize}
}{}{}{}


\gLayout{0}{Alarmzeichen}
{
Weitere klassische Alarmzeichen sind Symptome wie
\gEs{Atemnot}, \gE{plötzlich auftretende
Ödeme} oder \gE{akute Schmerzen}. 

Bei einem plötzlichen
Leistungsabfall des \gE{linken} Herzens kommt es zu einem
\gE{Rückstau} des Blutes \gE{in den Lungenkreislauf}. Aufgrund der Stauung
kommt es zum Austritt von Flüssigkeit in das
Zwischengewebe und es entsteht ein
\gEs{(kardiales) Lungenödem}. Ein typisches Zeichen dafür
ist das \gEs{brodelnde Atemgeräusch}\index{Atemgeräusch!brodelndes - bei
Herzinsuffizienz}. Fällt die Leistung des \gE{rechten}
Herzens, so kann man Beinödeme und/oder gestaute Halsvenen
beobachten. 
}
{
\begin{itemize}
    \item \gEs{Atemnot}
    \item  \gE{Akute Schmerzen}
    \item Stauung:
    \begin{itemize}
    \item \gEs{Lungenödem} mit \gE{"`brodelndem
    Atemgeräusch"'} (\gE{li.} Herz)
    \item \gEs{Beinödeme}, gestaute Halsvenen (\gE{re.}
    Herz)
    \end{itemize}
    \end{itemize}
}{}{}{}


Das kardiale Lungenödem wird im unter
\gSee{topic:lungenoedem} 
besprochen.

\gMassnahmenNG{Spezielle Maßnahmen}{Herzinsuffizienz}{}{
    \begin{itemize}
        \item Lagerung: Oberkörper hoch
        \item Vitale Gefährdung beurteilen \par \gTxBeiVit
        \gTxMassVitBes
        \begin{itemize}
            \item \ce{O2} 10\gLMin, bei Bedarf bis zu 15\gLMin
        \end{itemize}
        \item \ce{O2}: je nach Symptomen 6-8\gLMin, bei
        Bedarf bis zu 15\gLMin
        \item Vitalwerte erheben \& Monitoring
        \item Ursachenforschung
    \end{itemize}
}

\subsubsection{Linksherzinsuffizienz}
\index{Herzinsuffizienz!Links-}\index{Insuffizienz!Linksherz-}\index{Linksherzinsuffizienz}

Die Linksherzinsuffizienz entsteht aufgrund mangelnder Auswurfleistung
des linken Herzens. Durch den \gE{Rückstau von Blut in den
Lungenkreislauf} kann es zu einem \gEs{kardialen Lungenödem}
kommen.

\gLayout{0}{Kardiales Lungenödem und Linksherzinsuffizienz}
{\index{Lungenödem!kardiales - bei Herzinsuffizienz} Patienten mit einem 
\gE{kardialem Lungenödem} sind \gE{lebensbedrohlich} krank. Das Herz verbraucht oft die allerletzten Reserven, schon die
Aufregung durch den Transport im Stiegenhaus zum Auto kann das Herz
endgültig in den Herz-Kreislaufstillstand bringen. 
}{
\begin{itemize}
    \item \gCa{Kardiales Lungenödem = lebensgefährlich
    krank!}
    \item Herz verbraucht die allerletzten Reserven!
\end{itemize}
}{}{}{}

Es kommt immer wieder vor, dass solche Patienten im
Stiegenhaus oder im Auto reanimationspflichtig werden. 

\gCave{Kein Transport von Patienten mit Lungenödem ohne Notarzt. Auch wenn
das Spital {\quotedblbase}um{\textquotesingle}s Eck{\textquotedblleft}
ist.}


\gLayout{0}{\gTxSymptome}{
Der Patient klagt über Atemnot (\gEs{Dyspnoe}). Oft
berichten die Patienten, dass sie nur mehr im Sitzen Luft
bekommen, sie können nicht flach liegen (\gZ{Orthopnoe}).

Die Atemfrequenz und Herzfrequenz sind oft erhöht
(\gEs{Tachypnoe},
\gEs{Tachykardie}\index{Hyperventilation!bei
Herzinsuffizienz}), oft besteht ein \gEs{Hustenreiz}.
Mitunter wirken die Patienten unruhig und ängstlich.

Grundsätzlich ist der Blutdruck vermindert
(\gEs{Hypotonie}). Allerdings kann ein erhöhter Blutdruck
bei einer Blutdruckkrise (\gEs{Hypertonie}, \gE{hypertensive
Krise} (\prettyref{topic:hypertensive-krise}), \gE{hypertensiver
Notfall} (\prettyref{topic:hypertensiver-notfall})) der eigentliche
\gE{Auslöser} einer Herzinsuffizenz sein. Der Blutdruck
kann dann natürlich erhöht sein. 

Ein besonders kritisches Symptom ist das \gEs{Lungenödem}.
Wenn es sehr stark ausgeprägt ist, kann man es mit freiem
Ohr hören (\gE{"`brodelndes Atemgeräusch"'}). Diese
Patienten sind als besonders kritisch einzuschätzen.
}{
\begin{itemize}
    \item \gEs{Dyspnoe}
    \begin{itemize}
        \item \gZ{Orthopnoe:} Patient bekommt nur mehr im Sitzen Luft, kann
        nicht flach schlafen, etc.
    \end{itemize}
    \item \gEs{Tachypnoe}\index{Hyperventilation!bei Herzinsuffizienz}
    \item \gEs{Tachykardie}
    \item \gEs{Hustenreiz} 
    \item Unruhe
    \item \gEs{Hypotonie} oder \gEs{Hypertonie} (\gAbk{RR}\,↑↑ kann Auslöser
    sein)
    \item \gEs{Lungenstauung/Lungenödem} (Blut staut sich
    zurück in den Lungenkreislauf)
\end{itemize}
}{}{}{}






        




\subsubsection{Rechtsherzinsuffizienz}
\label{topic:rechtsherzinsuffizienz}

Die Rechtsherzinsuffizienz entsteht aufgrund mangelnder Auswurfleistung
des rechten Herzens.
\index{Rechtsherzinsuffizienz}\index{Herzinsuffizienz!Rechts-}
\index{Insuffizienz!Rechtsherz-}

\gLayout{0}{\gTxSymptome}
{
Durch den Rückstau des Blutes in das venöse Gefäßsystem
kommt es oft zu Ödemen (Beine, \ldots) und
gestauten Halsvenen. 

Die Patienten müssen häufig nachts urinieren. Über
längere Zeit kann sich die Leber vergrößern und ein sog.
"`Wasserbauch"' (\gZ{Aszites}) bilden\gZ{, dabei wird durch den Rückstau
Wasser in den freien Bauchraum abgepresst.} 
}{
\begin{itemize}
	\item \gEs{Bein-}/\gEs{Hautödeme}
    \item \gEs{gestaute Halsvenen} 
	\item Nächtliches Wasserlassen
	\item Lebervergrößerung, evt. mit Aszites
	({\quotedblbase}Wasserbauch{\textquotedblleft})
\end{itemize}
}{}{}{}


% \begin{wrapfigure}{r}
% %	\begin{figure}
% 		%\gCaptionof{figure}{Aszites und Bein"odeme}	
% 		\includegraphics[width=5cm]{./images/aszites-beinoedeme.png}
% %	\end{figure}
% \end{wrapfigure}





\subsection{Koronare Herzkrankheit und Akutes Koronarsyndrom}
\index{Koronare
Herzkrankheit}\index{Herzkrankheit!Koronare}\index{KHK}
\index{Akutes Koronarsyndrom}\index{Koronarsyndrom!Akutes}


\begin{description}
	\item[Die \gT{Koronare Herzkrankheit}
	(\gT{KHK})] ist eine Erkrankung der
	\gEs{Herzkranzgefäße} (Koronargefäße)
	bei der es zu einer Verengung der Gefäße kommt.

	\item[Das \gT{Akute Koronarsyndrom}]\index{Akutes Koronarsyndrom} ist ein
	Sammelbegriff für \gE{Herzinfarkt}, \gE{stabile} und \gE{instabile Angina pectoris}.
	Bei diesen Krankheitsbildern kommt es zu einer akuten
	\gEs{Unterversorgung des Herzmuskels} mit Blut und
	Sauerstoff. Die Ursache sind in erster Linie verengte oder verstopfte
	Herzkranzgefäße. 

\end{description}

\subsubsection[Angina pectoris]{Angina pectoris}
\index{Angina pectoris}

\gDefinition{Plötzliches Engegefühl in der Brust meist mit
Thoraxschmerz infolge einer vorübergehenden
Unter\-ver\-sorgung des Herzmuskels mit sauerstoffreichen Blut.}

\begin{itemize}
	\item Minderdurchblutung des Herzmuskels (Myokard) aufgrund einer oder
	mehrerer Einengung(en) von Herzkranzgefäßen (Koronararterien)
	oder
	\item Erhöhter Sauerstoffbedarf des Herzens bei Anstrengung
\end{itemize}
{\mdseries
${\rightarrow}$ Missverhältnis Sauerstoffangebot -- 
Sauerstoffbedarf des Herzens ${\rightarrow}$ Schmerz!}

Der akute Angina-Pectoris-Anfall gehört zu der Familie des Akuten
Koronarsyndroms. Dazu gehört auch der Herzinfarkt, dem der gleiche
Mechanismus wie der der Angina Pectoris zu Grunde liegt, wobei es beim Infarkt
im Gegensatz zur Angina pectoris zu einem \gE{totalen}
Herzkranzgefäß\gE{verschluss} mit Absterben von Herzmuskelgewebe kommt.
\gA{KOCH: nein; GABS: Doch. }

Bei den meisten Patienten ist die Angina Pectoris schon bekannt, bzw. es
besteht eine {\quotedblbase}\index{Koronare
Herz-Krankheit}\gT{Koronare Herz-Krankheit}
(\index{KHK}\gT{KHK}){\textquotedblleft}. Die KHK ist
eine Erkrankung der Herzkranzgefäße, diese werden durch
Ablagerung von Fetten und Kalk geschädigt und eingeengt, dadurch kann
weniger Blut durchfließen. So kommt es zu einer Mangelversorgung
des Herzmuskels und zu den Symptomen der Angina pectoris.  


\gLayout{0}{\gTxSymptome}
{
Das Leitsymptom ist der \gEs{akute, nicht atemabhängige,
Brustschmerz}. Der Schmerz ist meist hinter dem Brustbein lokalisiert und
\gEs{strahlt häufig aus}, häufig in den Kiefer, (linken) Arm und in den Rücken,
evt. auch in den Oberbauch. Der Schmerz wird als brennend bzw. stechend beschrieben.
Dazu kommt i.d.R. ein \gEs{Druckgefühl} am Brustkorb, \emph{"`als würde jemand
auf dem Brustkorb stehen."'} \gEs{Todesangst und Vernichtungsgefühl} ist
häufig. 

Ein weiteres wichtiges Leitsymptom ist die \gEs{Atemnot}
(\gE{Dyspnoe}).

Ein akuter Angina pectoris-Anfall ist primär \gEs{nicht sicher von einem
Herzinfarkt unterscheidbar}. Die Symptome treten eher bei Belastung auf und
vergehen oft bei Entspannung, das Vorübergehen der Beschwerden darf aber
natürlich nicht abgewartet werden! 

}{
\begin{itemize}
    \item Brustschmerz, brenndend/stechend, Druckgefühl
    \item Oft Ausstrahlung häufig in den Kiefer,
    linken Arm, Rücken, evt. Oberbauch.
    \item Todesangst, Vernichtungsgefühl.
    \item Dyspnoe.
    \item Primär \gEs{nicht sicher von MCI unterscheidbar}.
    \item Symptome eher bei Belastung, vergehen oft bei Entspannung.
\end{itemize}

}{}{}{}

\gParallel{
\includegraphics[width=\linewidth]{./images/noauth/AASdev20090228-img73.png}
}{
\gCaptionof{figure}{\gAuthNotes{Bild fehlt!}
Schmerzausdehnung beim ischämischen Herzschmerz.}
}


% \gMarried
% {
%  \includegraphics[width=\textwidth]{./images/noauth/AASdev20090228-img72.png}
% \gCaptionof{figure}{\gAuthNotes{Bild fehlt!} Patient mit
% einem akuten Koronarsyndrom, vermutlich Angina Pectoris. Bei kaltem Wetter, beim
% Stiegen steigen mit schwerem Gepäck: Ein plötzliches Engegefühl
% im Brustkorb und ein stechender Schmerz im linken Thorax.}
% }{
% \includegraphics[width=\textwidth]{./images/noauth/AASdev20090228-img73.png}
% \gCaptionof{figure}{\gAuthNotes{Bild fehlt!}
% Schmerzausdehnung beim ischämischen Herzschmerz.}
% }

\gMassnahmenNG{Spezielle Maßnahmen}{Angina pectoris}{}{
\begin{itemize}
	\item \gTxMassVitBes
	\begin{itemize}
		\item Lagerung: Oberkörper hoch
		\item \ce{O2} 10\gLMin, bei Bedarf bis zu 15\gLMin
	\end{itemize}
	\item \gEs{striktes Bewegungsverbot}!
	\item Patient darf eigenes Nitro nur  bei RR
	{\textgreater}100 syst. nehmen!
\end{itemize}
\gAuthNotes{\#13: ?? eigenes Nitro nehmen?

Vorschlag: Nur wenn sein Nitro leer war und er es deswegen nicht
genommen hat bei RRsyst{\textgreater}110

Wenn er es schon genommen hat und es nicht geholfen hat
dann nicht <-- KOCH richtig wichtig
GABS: was soll ``richtig wichtig bedeuten? ob es der Pat. richtig genommen hat?

KOCH: Pat. darf eigenes Nitro nur  bei RR {\textgreater}100 syst.
\underline{selbst ein}nehmen!}
}

\subsubsection{Herzinfarkt - Myokardinfarkt (MCI)}
\index{Herzinfarkt}\index{Infarkt!Herz-}\index{Myokardinfarkt}\index{MCI}
\index{Herzkranzgefäße!Verschluss}\index{Koronargefäße!Verschluss}

\gDefinition{Verschluss eines oder mehrerer Herzkranzgefäße mit
anschließendem Absterben des betroffenen Herzmuskelgewebes.}

\gLayout{0}{\gTxSymptome}
{
Die Symptome entsprechen im Wesentlichen denen der Angina pectoris, doch
vergehen diese nicht mehr und es bleiben -- selbst bei optimaler Behandlung --
Schäden am Herzmuskel bestehen. Bei Diabetikern kann es aufgrund der
Nervenschädigungen auch zu einem schmerzlosen Infarkt kommen, man
spricht dabei auch vom \gT{"`stummen
Infarkt"'}\index{Stummer Infarkt}.

Sonst steht normalerweise der \gE{stechende/brennende}, nicht atemabhängige
Brustschmerz zusammen mit \gEs{Atemnot} und \gEs{Todesangst} im Vordergrund. Der
\gEs{Brustschmerz} kann in Richtung Magen, rechter oder linker Hand, Kiefer
oder auch Rücken \gEs{ausstrahlen}. Zusätzlich zu dem
Schmerz verspürt der Patient auch oft ein \gEs{Engegefühl} im
Brustkorb, \emph{"`als würde jemand darauf stehen"'}. Die
Atemnot kann massiv ausfallen. 

Weiters können aufgrund des eintretenden \gEs{kardiogenen
Schocks} die entsprechenden Schockzeichen beobachtet
werden: Blässe, Kaltschweißigkeit, Hypotonie, etc.
(\ref{topic:schock-kardiogener}).
}{
\begin{itemize}
    \item Symptome wie bei Angina Pectoris, vergehen aber nicht
    mehr
    \begin{itemize}
        \item Stechender/brennender \gEs{Brustschmerz} mit Ausstrahlung.
        \item "`Stummer Infarkt"' bei Diabetikern.
        \item Druckgefühl, {\quotedblbase}als würde jemand am
        Brustkorb stehen{\textquotedblleft}.
        \item Todesangst.
        \item \gEs{Dyspnoe.}
        \item Kardiogener Schock.
        %	\item Keine ausreichende Besserung auf Nitro
    \end{itemize}
\end{itemize}
}{}{}{}
  
\gLayout{0}{\gTxKomplikationen}
{
\begin{itemize}
    \item Herzrhythmusstörungen (alle Arten), bis hin
    zum Atem-/Kreislaufstillstand und Reanimation.
    Besonders häufig kommt es als Frühkomplikation zum
    plötzlichen Kammerflimmern.
    \item Kardiogener Schock \index{Kardiogener
    Schock!beim Herzinfarkt}\index{Schock!Kardiogener!beim
    Herzinfarkt} 
\end{itemize}

}
{
\begin{itemize}
    \item \gEs{Lebensgefahr!}
    \item Jede Herzrhythmusstörung als Komplikation möglich.
    \item Kammerflimmern als Frühkomplikation häufig!
\end{itemize}

}{}{}{}


% \begin{center}
%  \gTempPicMissing{Infarkt [Warning: Image ignored]
%  {AASdev20090228-img74.png}}
% \end{center}

\gMassnahmenNG{Spezielle Maßnahmen}{Herzinfarkt}{}{   
    \begin{itemize}
        \item \gTxMassVitBes 
        \begin{itemize}
            \item Lagerung: Oberkörper hoch
            \item \ce{O2} 10\gLMin, bei Bedarf bis zu 15\gLMin
        \end{itemize}
        \item Beengende Kleidung öffnen
        \item Striktes Bewegungsverbot!
    \end{itemize}
    
    %%%%%%%%%%%%%%%%%%%%%%%%%%%%%%%%%%%%%%%%%%%%%%%%%%%%%%%%%%%
% TODO : ERRATUM: Nitroeinnahme MCI Fehlerhafter Text
\begin{anfxerror}[author=GABS]{ERRATUM: Fehlerhafter Text: Nitroeinnahme MCI}
   Nitroeinnahme MCI: Darauf wird nicht eingegangen. Dieser Teil ist fehlerhaft
   oder sehr unvollständig. Die Aussagen sind unklar oder verwirrend. Eine Überarbeitung ist dringend notwendig!
\end{anfxerror}
%%%%%%%%%%%%%%%%%%%%%%%%%%%%%%%%%%%%%%%%%%%%%%%%%%%%%%%%%%%
}


\gLayout{0}{Weitere Therapie}
{

Herzinfarktpatienten werden {\quotedblbase}besonders{\textquotedblleft}
behandelt und normalerweise nicht auf eine beliebige Interne Station
gebracht\cite{Kalla2006}. Meistens stehen dem Notarzt zur
Auswahl:

\begin{itemize}
	\item Lyse-Therapie mittels eines speziellen Medikaments, oder
	\item Einweisung in ein \index{Herzkatheterlabor}\gEs{Herzkatheterlabor}
	({\quotedblbase}PCI{\textquotedblleft},
	{\quotedblbase}\index{PTCA}\gEs{PTCA}{\textquotedblleft}
	\gZ{{}- Perkutane Transluminale Coronar
	Angioplastie})

\end{itemize}

}
{
\begin{itemize}
    \item Lyse
    \item Herzkatheterlabor (PTCA)
\end{itemize}

}{}{}{}


In Wien gibt es einen \gE{Dienstplan} für die Katheterlabore, d.h. 
jedes Katheterlabor hat an bestimmten Tagen der Woche rund um die 
Uhr Bereitschaftsdienst. Eine Voranmeldung und vorangehende Absprache 
zwischen Notarzt und diensthabenden Oberarzt der Abteilung
ist erforderlich. Die PTCA ist in Wien u.a. verfügbar im AKH, Donauspital, und KH Hietzing.

\gCave{{\bfseries Kein Transport ohne Notarzt!}
	
	Auch nicht wenn
	das Spital {\quotedblbase}um{\textquotesingle}s Eck{\textquotedblleft}
	ist.

	Der Patient braucht vermutlich {\bfseries kein beliebiges
	Spital} sondern ein {\bfseries dienst\-bereites
	Herzkatheterlabor}. Der Notarzt muss entscheiden, ob das
	notwendig und sinnvoll ist.

	Auch wenn der Patient ein Einweisungsformular (Spitalszettel) von einem	
	Arzt hat (Ärztefunkdienst, Hausarzt) oder ein Hausarzt
	meint {\quotedblbase}es sei eh nicht dringend{\textquotedblleft} unbedingt
	einen Notarzt beiziehen.
}

\subsection{Herzrhythmusstörungen}
\label{topic:herzrhythmusstoerungen}
\index{Herzrhythmusstörungen}
\index{Rhythmusstörungen!Herz-}



\begin{table}[h]
\begin{gWideParEnv}
\gCaptionof{table}{Übersicht: Wichtige und häufige
Herzrhythmusstörungen.\label{tab:herzrhythmusstoerungen}}

% Index-Einträge für Tabelle ---
\index{Bradykardie} 
\index{Tachykardie}
\index{Kammerflimmern}\index{Flimmern!Kammer-}
\index{Vorhofflimmern}\index{Flimmern!Vorhof-}
\index{Asystolie}
\index{Pulslose Ventrikuläre Tachykardie} 
\index{Tachylardie!Pulslose Ventrikuläre}
% --- Ende


\begin{minipage}{\linewidth}
\begin{tabulary}{\textwidth}{LLL}
\gTabHeading{Rhythmus}
& 
\gTabHeading{Bedeutung}
& 
\gTabHeading{Bemerkung}
\\\midrule\addlinespace
\gTabSub{Bradykardie}

{\scriptsize (HF\,{\textless}\,60/min)}\gMarginBugfix{(beim
Erwachsenen)} 
& 
Situations\-abhängig 
& 
\emph{"`zu langsam"'}
\\\addlinespace
\gTabSub{Tachykardie} 

{\scriptsize (HF\,{\textgreater}\,100\gMarginBugfix{(beim
Erwachsenen)}/min) }
& 
Situations\-abhängig 
& 
\emph{"`zu schnell"'}
\\\addlinespace\midrule \addlinespace
\gTabSub{Kammerflimmern}
& Reanimation  & \emph{"`schockbarer"'} Rhythmus 
\\\addlinespace
\gTabSub{Pulslose Ventrikuläre Tachy\-kardie} 
& Reanimation 
& \emph{"`schockbarer"'} Rhythmus 
\\\addlinespace
\gTabSub{Asystolie} 

{\scriptsize (HF\,=\,0)} 
& 
Reanimation 
& 
Nulllinie, \emph{"`nicht-schockbarer"'} Rhythmus
\\\addlinespace
\gTabSub{Vorhofflimmern}
& regellose Erregung im Vorhof

oft symptomlos 
& 
\gE{Chronische Erkrankung.} Sehr viele Leute haben
Vorhofflimmern und leben ganz gut damit. 

Risikofaktor für:

\gE{Tachykardie-Anfälle}

\gE{Thrombenbildung}, dadurch erhöhtes Lungenembolie-
und Schlaganfallrisiko. Deshalb Vorbeugung mit
gerinnungs\-hemmenden ({\quotedblbase}blut\-ver\-dünnenden{\textquotedblleft}) Medi\-ka\-menten:
z.B. \gE{"`Marcoumar"'} ${\rightarrow}$ Patient ist blutungsgefährdet
\\\bottomrule
\end{tabulary}
\end{minipage}
\end{gWideParEnv}
\end{table}



\subsubsection{Hier zählt die Geschwindigkeit: Tachy- und Bradykardie} 
\index{Bradykardie}
\index{Tachykardie}

\begin{labeling}[~--]{Tachykardie}
    \item[\gT{Bradykardie}]  zu langsamer Herzschlag     
    ({HF \textless} 60\,\nicefrac{}{min} {\tiny (beim Erwachsenen)})
    
    Leistungssportler können aufgrund ihres trainierten Körpers auch sehr
    niedrige Ruhepulsfrequenzen haben.
    \item[\gT{Tachykardie}]  zu schneller
    \marginpar{Tachometer im Auto:
    Geschwindigkeitsanzeiger, meistens zu schnell ;)}
    Herzschlag     
    (HF {\textgreater} 100\,\nicefrac{}{min} {\tiny (beim
    Erwachsenen)})
\end{labeling}



Denke bei der Beurteilung auch immer an die Umstände und den
Erregungszustand des Patienten!



\subsubsection{Tachykarde Attacke}

\gDefinition{Plötzliche Tachykardie ohne erkennbare Ursache, wobei der Patient
(meistens) Symptome wahrnimmt und eine HF\,>\, 160\,\,\nicefrac{}{min}
vorliegt.}


\gLayout{0}{\Info}
{
Tachykarde Attacken sind ein häufiges Notfallbild. Da oft chronische
Erkrankungen die Ursache sind, haben die Patienten diese
Attacken immer wieder und kennen die Symptome bereits (und oft auch die Behandlung). Die
Ursachen können sehr unterschiedlich sein und sind auch vom Fachmann oft nicht
leicht zu ermitteln. 

Eine große Rolle spielt die \gEs{Anamnese}: Oft sind entsprechende
Grunderkrankungen schon seit langem bekannt, oft gibt es bereits entsprechende Arztbriefe oder
Befunde. Ein rasches Hinweisen auf diese Befunde ist für den nachbehandelnden
Arzt wichtig um eine rasche, zielführende Behandlung einzuleiten.
}{
\begin{itemize}
    \item Oft Grunderkrankung bekannt.
    \item Anamnese!
    \item Befunde, Arztbriefe!
\end{itemize}
}{}{}{}

\gLayout{0}{\gTxSymptome}
{Es liegt klarerweise eine \gEs{Tachykardie} vor. Der Puls
ist i.d.R. hoch, meist über 160\,\,\nicefrac{}{min}. Der Patient klagt oft über ein
\gEs{"`Herzklopfen"'} und ist unruhig. Manchmal wird der Anfall von
Angstgefühlen begleitet.

Oft kommen \gEs{Begleitsymtome} dazu, wie zum Beispiel Atemnot oder ein leichter
Brustschmerz. Dies sind erste Zeichen, dass das Herz an
seine Grenzen der Belastbarkeit stößt.
}{
\begin{itemize}
    \item Tachykardie.
    \item Spürbares Herzklopfen.
    \item evt. Atemnot.
    \item evt. Leichter Brustschmerz.
    \item evt. Angstgefühl.
\end{itemize}
}{}{}{}

\gLayout{0}{{\Lightning}Gefahr }
{
Da das \gEs{Herz plötzlich viel mehr Arbeit leisten
muss}, kann es sein, dass es an seine \gEs{Grenzen}
stößt und \gEs{versagt} bzw. die Blutversorgung
des Herzens für diese Arbeit nicht mehr reicht und ein
Angina-pectoris-Anfall oder sogar ein Herzinfarkt\gA{ KOCH: passt mit der
Definition MCI nicht überein, AP wäre besser oder? GABS: AP
dazugenommen, Infarkt aber gut möglich. } entsteht. 
}
{
\begin{itemize}
    \item Herz überfordert:
    \begin{itemize}
        \item Angina pectoris.
        \item Herzinfarkt.
    \end{itemize}
\end{itemize}
}{}{}{}





\gMassnahmenNG{Spezielle Maßnahmen}{Tachykarde Attacke}{}{
    \begin{itemize}
        \item \gTxBeiVit \gTxMassVitBes 
        \begin{itemize}
            \item Lagerung: Oberkörper hoch
            \item \ce{O2} 10\gLMin, bei Bedarf bis zu
            15\gLMin
            \item Notarzt: Eine medikamentöse Therapie ist meistens
            erforderlich oder ratsam.
        \end{itemize}
        \item Beengende Kleidung öffnen
        \item Striktes Bewegungsverbot!
        \item Psychische Betreuung! Oase der Ruhe schaffen. 
        \item Wenn der Patient versorgt ist (d.h. alle obigen Punkte erfüllt
        wurden) soll -- wenn möglich -- ein \gE{EKG}
        (Extremitätenableitung, \gSee{topic:ekg}) abgeleitet werden. Der
        Ausdruck dient v.a. der \gE{Dokumentation} und der Information für das
        nachbehandelnde Personal (Notarzt, Spital, Facharzt, \ldots).
    \end{itemize}
}



\subsubsection{Besondere Rhythmen}


\gLayout{0}{Kammerflimmern}
{\gCa{Reanimationspflichtig}. 
\index{Kammerflimmern}

NICHT verwechseln mit
Vorhofflimmern! Die elektrische Erregung im Herzen ist ungerichtet, als würden alle Autos 
eines Parkplatzes gleichzeitig in alle Richtungen losfahren 
${\rightarrow}$ der Herzmuskel kontrahiert nicht
geordnet, er "`zittert"' nur (wie ein Pudding), es wird
\gEs{keine Auswurfleistung} erreicht!
}
{
\begin{itemize}
        \item \gCa{Reanimation!}
    \item Wirre Erregung, keine sinnvolle Muskelarbeit ${\rightarrow}$ kein
    Auswurf.

\end{itemize}

}{}{}{}

    
\gLayout{0}{Pulslose Ventrikuläre Tachykardie}
{\gCa{Reanimationspflichtig!} 
\index{Pulslose Ventrikuläre Tachykardie}



Vorstufe zum Kammerflimmern. Das Herz kontrahiert
so schnell, dass es sich nicht mehr füllen kann
${\rightarrow}$ \gEs{keine Auswurfleistung}.
}
{
\begin{itemize}
    \item  \gCa{Reanimation!}
    \item Keine Zeit zum Füllen ${\rightarrow}$ kein Auswurf.
\end{itemize}

}




\gLayout{0}{Asystolie}
{\gCa{Reanimationspflichtig}. 
\index{Asystolie}



Im Herz entsteht keine elektrische
Erregung, der Herzmuskel kontrahiert daher nicht und es erfolgt
{keine Auswurfleistung} 
}{
\begin{itemize}
    \item \gCa{Reanimation!}
    \item Keine elektrische Erregung, keine Muskelarbeit, kein Auswurf.
\end{itemize}
}{}{}{}
    
\gLayout{0}{Vorhofflimmern}
{
\index{Vorhofflimmern}

Regellose Erregung im \gE{Vorhof}, die oft symptomlos ist. 
Sehr viele Leute haben Vorhofflimmern und leben ganz gut damit.

Komplikationen:
\begin{itemize}
    \item Tachykardie-Anfälle
    \item Thrombenbildung
    \begin{itemize}
        \item deshalb Therapie mit \gEs{gerinnungshemmenden}
        ({\quotedblbase}blutverdünnenden{\textquotedblleft}) Medikamenten:
        Marcoumar ${\rightarrow}$ Patient ist Blutungsgefährdet
    \end{itemize}
\end{itemize}
}
{
% TODO VIT: Slide fehlt!
\begin{itemize}
    \item Flimmern im \gEs{Vorhof}, Kammern arbeiten normal.
    \item Häufige und oft chronische Erkrankung.
    \item Thrombenbildung, Tachykardieanfälle.
    \item Oft Therapie mit \gE{gerinnngshemmenden Medikamenten}.
\end{itemize}

}{}{}{}


\subsection{Hypotonie}

\gDefinition{Unangemessen zu niedriger Blutdruck.}
 
\gEs{Richtwert:} Beim Erwachsenen RR {\textless} 100 mmHg
syst. Jeder Mensch hat allerdings im gewissen Rahmen einen
"`eigenen Normalwert"', bei Kindern gelten grundsätzlich
andere Werte.

% einvernähmlich gestrichen KOCH 2008-05-18
% 
% \paragraph{Ursachen}
% 
% \begin{itemize}
%     \item Cardial
%     \item Volumenmangel, Exsikkose
%     \item Regulationsstörung (vasovagal)
%     \item Sportler, junge Frauen
%     \item Schock
% \end{itemize}
% 
% \paragraph{\gTxSymptome}
% 
% \begin{itemize}
%     \item Müdigkeit bis Kollaps
%     \item RR-Abfall mit Bewußtseinseintrübung
%     \item {\quotedblbase}schwarzwerden vor Augen{\textquotedblleft}, Ohnmacht
%     \item Sturz = {\quotedblbase}Therapie{\textquotedblleft}
% \end{itemize}


\subsubsection{Kollaps, Synkope}

\gDefinition{passagere Kreislaufregulationsstörung  mit RR-Abfall,
    Blutverteilungsstörung  ohne Volumsverlust}

\paragraph{\gTxSymptome}

\begin{itemize}
    \item RR ${\downarrow}$
    \item Schwindel (Vertigo )
    \item evt. kurze Bewusstlosigkeit, Ohnmacht
    \item Patient sackt zusammen
    \item Bradykardie oder Tachykardie
    \item Schwarzwerden vor Augen
\end{itemize}

Patient wacht auf dem Boden liegend  auf, weiß nicht, was geschehen
ist...

\paragraph{Ursachen} Die Ursachen sind präklinisch oft
nicht sicher ermittelbar, mitunter bleibt die wahre Ursache
für immer ein Rätsel. Daher ist eine Hospitalisierung zur
bestmöglichen Abklärung angeraten. 

Es gibt jedoch auch typische "`Kollaps-Situationen"': Ein
schwüler Tag in einer überfüllten U-Bahn-Garnitur, nicht
gefrühstückt, auf dem Weg zu einem Vorstellungstermin und
seit gestern in der Menstruation. Oder im stickigen
Veranstaltungszelt bei 40\textdegree , wobei der der ganze
Flüssigkeitskonsum aus 5 (entwässernden) Dosen
RedBull\texttrademark\footnote{Achtung! Schleichwerbung!}
bestand. So mancher war auch ob eines Begräbnisses so
ergriffen, dass er sich fast in die Grube neben den Sarg
dazugelegt hätte \ldots 

ausgehend von:

\begin{itemize}
    \item Herz (z.B. Herzrhythmusstörungen)
    \item Gefäßen (z.B. Hitzebedingte Erweiterung
    der Gefäße mit "`Versacken"' des Blutes in den Beinen)
    \item Blut (z.B. Blutarmut (\gT{Anämie}),
    Regelblutung)
    \item Hirn
    \item Hirngefäßen (z.B. TIA,
    \gSee{topic:tia} )
\end{itemize}

\gMassnahmenNG{Spezielle Maßnahmen}{Kollaps/Synkope}{}{
    
    \begin{itemize}
        \item Erheben ob eine akute vitale Bedrohung besteht
        \item Ausschluss von Notfallsdiagnosen soweit
        möglich, die gesamten zur Verfügung stehenden
        Möglichkeiten sind auszuschöpfen.
        \item Im Besonderen:
        \begin{itemize}
            \item Traumacheck (Sturz?)
            \item Neurocheck inkl. Blutzuckermessung
            \item Temperatur
            \item Ausführliche Anamnese

        \end{itemize}
        
        \item Lagerung: Nach Ausschluss von Herzbeschwerden,
        Atemnot und relevanten Verletzungen: Beine hoch. 
        \item Hospitalisierung zur Abklärung
        anstreben\newline
        Abteilung: Innere Medizin
        
    \end{itemize}
    
}

\subsection{Hypertonie}

\index{Hypertonie}\gT{Hypertonie}
(\index{Bluthochdruck}\gT{Bluthochdruck}): Oft
symptomlose \textit{chronische} Erkrankung welche auf Dauer schwere
Schäden am Körper anrichten kann und im Zuge von plötzlichen
Blutdruckkrisen auch akut Probleme machen kann. 

\begin{itemize}
    \item Chronische Hypertonie: Obergrenze derzeit bei 140/90
    (außer bei Kindern)
\end{itemize}

\gCave{Bei \gEs{Schwangeren} sind die Grenzwerte genau zu beachten!
(\gEs{EPH-Gestose}, \prettyref{topic:eph-gestose})!}

\paragraph{Ursachen}

\begin{itemize}
    \item Widerstandserhöhung (durch Engstellung der Gefäße)
    \item vermehrte Herzarbeit
    \item Sympathikusaktivierung, Adrenalin+Cortison (Stresshormone) durch
    Dauerstress,  Bewegungsmangel, genetische Veranlagung
    \item u.a. \ldots  
\end{itemize}

\paragraph{Langzeitfolgen}

Die Hypertonie hat in der Notfallmedizin auch eine wichtige
indirekte Bedeutung: Sie ist Ursache für viele andere
Erkrankungen, die ihrerseits akut lebensbedrohlich werden
können. 

Zwei Wirkungen der Hypertonie sind dabei besonders wichtig:
Die \gEs{Gefäßschädigung} und die \gEs{Belastung des Herzens durch
den hohen Gegendruck}. 

\begin{wrapfigure}{r}{3cm}
\includegraphics[width=\linewidth]{./images/teaser/imgp0481.jpg}
\end{wrapfigure}Bei der Gefäßschädigung kommt es zur Ablagerung durch Kalk
und sonstige Plaques und über längere Zeit zur \gE{Verengung
oder Verstopfung der Gefäße}. Je nach Organ kann das sehr
massive Auswirkungen haben. Im Herz führt es zur \gE{Angina
Pectoris} oder zum \gE{Herzinfarkt}, im Hirn zu \gE{Schlaganfällen}, in
den Nieren zu einer \gE{Niereninsuffizienz}, usw. Selbst große
Gefäße wie die Hauptschlagader (Aorta) können derartig
geschädigt werden dass sie reißen können (Aortenaneurysma) --- eine
für den Patienten sehr lebensgefährliche Komplikation. 

Die Schädigung des Herzmuskels ist die Folge von dem
erhöhten Druck gegen den das Herz pumpen muss. Der Muskel
vergrößert sich um mit der Belastung fertig zu werden, kann
aber dann nicht mehr ausreichend mit sauerstoffreichen
Blut versorgt werden. 

Wir haben also gesehen, dass die Hypertonie ein wichtiger
\gE{Risikofaktor} für viele schwere Notfälle ist. Daher ist es
auch verständlich dass die Behandlung der Hypertonie einen
großen Stellenwert in der medizinischen Versorgung der
Bevölkerung hat. 

\gParallel{
\begin{itemize}
    \item \gEs{Gefäßschädigung}
    \begin{itemize}
        \item Herzkranzgefäße ${\rightarrow}$ Infarkt
        \item Hirngefäße ${\rightarrow}$ Schlaganfall
        \item Nierengefäße ${\rightarrow}$ Niereninsuffizienz ${\rightarrow}$
        Dialyse ${\rightarrow}$ typischer KTW-Patient
        \item Hauptschlagader ${\rightarrow}$ Aneurysma        
    \end{itemize}
\end{itemize}
}{
\begin{itemize}
    \item \gEs{Herzmuskelschädigung} \par  (Herz muss gegen
    größeren Widerstand pumpen)
    \begin{itemize}
        \item krankhafte Vergrößerung des Herzmuskels
    \end{itemize}
\end{itemize}
}



\subsubsection{Hypertensive Krise und Hypertensiver Notfall}
\label{topic:hypertensive-krise}
\label{topic:hypertensiver-notfall}

\gLayout{2}{\gTxBeschreibung}
{
Eine plötzliche, sehr starke Erhöhung des Blutdrucks wird als hypertensive
Krise oder hypertensiver Notfall bezeichnet. Der Unterschied zwischen den
beiden Diagnosen ist das Fehlen bzw. Vorhandensein von typischen Symptomen, s.
\prettyref{tab:hxpertension-krise-notfall}.

}{}{}{}{}

\gLayout{0}{\gTxSymptome}
{
S. \prettyref{tab:hxpertension-krise-notfall}. Kopfschmerzen und Nasenbluten
(\gT{Epistaxis}\index{Epistaxis}) sind häufige Leitsymptome für einen
hypertensiven Notfall. 

}{
S. \prettyref{tab:hxpertension-krise-notfall}.
}{}{}{}

\begin{table}[H]
\gCaptionof{table}{Unterscheidung Hypertensive Krise und Hypertensiver
Notfall.\label{tab:hxpertension-krise-notfall}}

\begin{tabularx}{\textwidth}{XX}
\gTabHeading{Hypertensive Krise}
&
\gTabHeading{Hypertensiver Notfall}
\\\addlinespace\midrule\addlinespace
\gAbk{RR}\,{\textgreater}\,230/130\,mmHg\footnote{Richtwert}

\gEs{Patient fühlt sich gut}
 &
\gAbk{RR}\,{\textgreater}\,230/130\,mmHg\footnote{Richtwert}

\gEs{+ Symptome}
\\\addlinespace
Zufallsbefund?

Oder steht der Blutdruck doch in Zusammenhang mit der
Berufung?

\textit{Patient hat überhaupt keine Beschwerden?} Warum hat
er Sie dann überhaupt gerufen?

Wenn es wirklich nur ein Zufallsbefund ist ({\quotedblbase}Gesundheitstage{\textquotedblleft},
{\quotedblbase}Vorsorgeuntersuchung{\textquotedblleft}):


&
\subparagraph{\gCa{Symptome:}}

\begin{itemize}
    \item Kopfschmerz
    \item Sehstörungen
    \item Brustschmerz
    \item evt. Nasenbluten
    \item Übelkeit, Erbrechen, Schwindel
\end{itemize}

\gCave{Gefahr des Herzversagens -- Das Herz muss
plötzlich gegen großen Druck pumpen}

\\\addlinespace
Maßnahmen: Eine notfallmäßige Blutdrucksenkung ist (eher) nicht erforderlich.

&
Vitale Bedrohung

Eine frühzeitige Blutdrucksenkung (\caps{na}, \caps{nfs})
ist notwendig. 

Evt. Anwendung der Notkompetenz gem. SanG. 
\\\addlinespace\bottomrule
\end{tabularx}
\end{table}
\marginpar{\cite{WHOISH2003,JNC7Report,BLIM3,IntroClinEmergMed4,Vlcek2008}}

\gCo{\gZ{Die Unterscheidung zwischen "`Hypertensiver
Krise"' und "`Hypertensiven Notfall"' ist in der
Literatur nicht einheitlich. Oft wird auch nur
zwischen "`Hypertensiver Krise ohne Symptome"' bzw. "`mit
Symptomen"' unterschieden oder generell eine andere
Einteilung gewählt. Von diesen Spitzfindigkeiten abgesehen
gilt das Vorhandensein von Symptomen durchgehend als
entscheidendes Kriterium.
\cite{ClassenInnere5,KlinLeitAlle3,Fauci2008}}}
% Renz-Polster2006 p442
% IntroClinEmergMed4 Kap. 28, p 393



\gMassnahmenNG{Spezielle Maßnahmen}{Hypertensive Krise}{}
{
\begin{itemize}
    \item Lagerung: Oberkörper hoch
    \item evt. \ce{O2}-Gabe
    \item Beengende Kleidung öffnen
    \item Monitoring
    \item  Hospitalisierung oder Notarzt zur
    Behandlung/Belassung
\end{itemize}
Eine notfallmäßige Blutdrucksenkung ist nicht erforderlich.
}

\gMassnahmenNG{Spezielle Maßnahmen}{Hypertensiver Notfall}{}
{
\begin{itemize}
    \item \gTxMassVitBes
    \begin{itemize}
    \item Lagerung: Oberkörper hoch
    \end{itemize}
    \item Beengende Kleidung öffnen
\end{itemize}

Eine frühzeitige Blutdrucksenkung (\caps{na}, \caps{nfs})
ist notwendig. 

\ASBFLO{\gAbk{NFS} dürfen lt.
Algorithmus Nitrolingual-Spray verabreichen. Ist die Gabe erfolgreich kann u.U. von der
Anforderung eines Notarztes abgesehen werden.} 
}




\subsection{Gefäßaussackungen: Aneurysmen}

\subsubsection{Bauchaortenaneurysma}
\nocite{Meron2004}

%%%%%%%%%%%%%%%%%%%%%%%%%%%%%%%%%%%%%%%%%%%%%%%%%%%%%%%%%%%
% TODO : Überarbeitung
\gTODO{GABS}{????-??-??}{Ticket 34: Erklärung Aneurysma.}{Dieser Teil ist nicht gut formuliert und/oder unvollständig, 
   er benötigt eine Überarbeitung!}
%%%%%%%%%%%%%%%%%%%%%%%%%%%%%%%%%%%%%%%%%%%%%%%%%%%%%%%%%%%
\gLayout{0}{\gTxBeschreibung{} und Ursachen}
{%%% Gerhard Gabriel <gerhard.gabriel@grissu.org>
Bei der Entwicklung eines Bauchaortenaneurysmas ist zumeist
eine krankhafte Veränderung der Aortenwand
verantwortlich. Als Folge eines Unfallgeschehens,
oder auch nach Hebeversuchen schwerer Lasten, kann die
Gefäßwand im Bereich eines Aneurysmas einreißen. 
}
{\gSymbolText}{}{}{}

\gLayout{0}{\gTxSymptome}{
Je nach
Größe der Schädigung und entsprechendem Schweregrad der Blutung
im Läsionsbereich, kommt binnen kürzester Zeit unter
heftigem, akutem Bauchschmerz zum Tod. Kleinere Läsionen
sind ebenfalls von heftigem, akut einsetzendem Bauchschmerz
begleitet. Es bildet sich ein Blutungsschock (hypovolämischer
Schock) aus, der unbehandelt zum Tode führt.
}{
\begin{itemize}
    \item heftigste Schmerzen
    ("`Vernichtungsschmerz"')
    \item     hoher Blutverlust ${\rightarrow}$ schnell im
    Schock!
    \item evt. bereits bekanntes Aneurysma
    ${\rightarrow}$\,Anamnese
    \item evt. tastbare pulsierende Schwellung
\end{itemize}
}{}{}{}



Oft ist bei den Patienten das Aneurysma bereits bekannt und
war noch nicht groß genug um behandelt zu werden
(${\rightarrow}$\,Anamnese!). Ansonsten ist es schwer das
Aneurysma präklinisch zu diagnostizieren. Im Vordergrund
steht der Volumenverlust und der entstehende hypovolämische
Schock. 


\clearpage
\subsection{Gefäßverschlüsse in den Extremitäten}

\gLayout{0}{Unterscheidung}
{
Grundsätzlich unterscheidet man je nach betroffenem Gefäß
zwischen \gEs{arteriellen} und \gEs{venösen}
Gefäßverschlüssen. 

Beim \gE{arteriellen Verschluss} besteht
das Hauptproblem in der Unterversorgung der Extremität mit
Blut. Beim \gE{venösen Verschluss} kommt es zu einer Ab- und
Rückflussbehinderung. Zusätzlich besteht die Gefahr, dass sich
ein Gerinnsel losreißt und in einem anderen Teil des
Körpers Schaden anrichtet\footnote{Zum Beispiel in der
Lunge im Rahmen einer Lungenembolie}.
}{
\begin{itemize}
    \item \textbf{\textcolor[rgb]{0.7882353,0.0,0.08627451}{Arteriell}}
    \item \textbf{\textcolor{blue}{Venös}}:
    \index{tVT}\gT{tVT} (tiefe Venen Thrombose)
\end{itemize}
}{}{}{}



\subsubsection{Arterieller Gefäßverschluß}

\gLayout{0}{Beschreibung}
{
Beim \gE{arteriellen Gefäßverschluss} wird das
zuführende Gefäß verschlossen. Es kommt zu einer schmerzhaften 
Unterversorgung der Extremität mit Blut. Der komplette
Verschluss passiert plötzlich. 
}
{
\begin{itemize}
    \item Verschluß eines blut\gE{zuführenden} Gefäßes.
\end{itemize}
}{}{}{}

\gLayout{2}{Ursachen -- Warum kommt es zu einem Gefäßverschluss?}
{
Meistens liegt eine 
\gMargin{Chronische Gefäßschädigung}chronische Schädigung der Gefäße infolge anderer Erkankungen oder aufgrund eines ungesunden
Lebensstils vor. Ursachen für solch eine Gefäßschädigung
sind z.B. die Zuckerkrankheit (\gT{Diabetes mellitus},
s.\ref{topic:diabetes}), Hypertonie, erhöhte Blutfette oder
das Rauchen. Die Gefäße verkalken und verändern sich auch sonst
nicht zu ihrem Vorteil.
}
{
\begin{itemize}
  \item Chronische Gefäßschädigung, z.B. Bei Diabetes mellitus, Hypertonie \ldots
\end{itemize}
}{}{}{}



\gLayout{0}{pAVK}{
Wenn die Gefäße nur verengt sind spricht man von
der \gEs{peripheren
arteriellen Verschluss\-krankheit}, abgekürzt \gT{paVK}. Sie ist eine
häufige chronische Erkrankung, die man oft in Arztbriefen oder auf
Transportscheinen findet.
}{
\begin{itemize}
    \item Periphere arterielle
Ver\-schluss\-krank\-heit.
\end{itemize}

}{}{}{}

 (\gE{Achtung}: Obwohl es
"`\emph{Verschluss}krankheit"' heißt, sind die Gefäße nicht komplett
verschlossen, sondern nur verengt!)

\gLayout{}{}{}{}{}{}{}

\gLayout{0}{\gTxSymptome}
{
Die Symptome ergeben sich aus der Unterversorgung der
Extremität mit sauerstoffreichem Blut: sie ist blass,
kalt, tut höllisch weh und die Rekap-Zeit ist nicht
mehr sinnvoll messbar. Zusätzlich wird die Extremität
gefühllos und gelähmt. Der Patient kann ein Kribbeln
empfinden (\gE{"`Ameisenlaufen"'}. Tabelle \ref{tab:5p}
listet die klassischen Symptome auf \gZ{(die \emph{"`5
englischen P's"'})}.
}{
\gCaptionof{table}{Die 5 Zeichen des arteriellen Gefäßverschluss.}\label{tab:5p}

\begin{tabulary}{4cm}{CLL}
\midrule\addlinespace
\centering 1. & \noindent\gZ{Pain}        & \noindent\gEs{schmerzend}
\\
\centering 2. & \noindent\gZ{Pale}        & \noindent\gEs{blaß}
\\
\centering 3. & \noindent\gZ{Pulsless}    & \noindent\gEs{pulslos}
\\
\centering 4. & \noindent\gZ{Paresthesia} & \noindent\gEs{gefühllos}
\\
\centering 5. & \noindent\gZ{Paralysis}   & \noindent\gEs{gelähmt}
\\\addlinespace\bottomrule
\end{tabulary}
}{}{}{}


\gMassnahmenNG{Spezielle Maßnahmen}{Arterieller Gefäßverschluss}{}{
    \begin{itemize}
        \item Lagerung: Extremität hängen lassen, weich und warm lagern
        \item Schonend transportieren
        \item Nüchtern
        \item Hospitalisation: Abteilung für Chirurgie
        (Gefäßchirurgie)
    \end{itemize}
}

\subsubsection{Tiefe Venenthrombose}

\gDefinition{Blutgerinnsel (Thrombus\footnote{Thrombus: Blutgerinnsel. Im
Gegensatz dazu \gT{Embolus}: Losgelöster Thrombus, welcher woanders ein Gefäß
verstopft und eine \gT{Embolie} verursacht.}) verstopft Vene}


\gLayout{0}{Ursachen}
{
\begin{itemize}
    \item Venöse Insuffizienz (Krampfadern = Venenaussackungen aufgrund
    gestörten  Blutrückflusses, Venenklappeninsuffizienz), Rückstau
    \item Immobilisation
    \item Bettlägrigkeit
    \item Lange Reisen ohne Bewegung: \index{Economy class
    Syndrom}\gT{Economy class Syndrom}
\end{itemize}
}{
\begin{itemize}
    \item Venöse Insuffizienz (Krampfadern), Rückstau
    \item Immobilisation, Bettlägrigkeit
    \item Lange Reisen ohne Bewegung: (\gT{Economy class Syndrom})
\end{itemize}
}{}{}{}

\gLayout{0}{\gTxSymptome}
{
\begin{itemize}
    \item Schmerzen
    \item Schwellung (einseitig, anderes Bein vergleichen!)
    \item Dunkel-bläuliche (livide ) Verfärbung (einseitig, anderes Bein
    vergleichen!)
    \item Überwärmung (einseitig, anderes Bein vergleichen!)
\end{itemize}
}{
\begin{itemize}
    \item Schmerzen
    \item Schwellung (einseitig!)
    \item Dunkel-bläuliche (livide ) Verfärbung
    \item Überwärmung    
\end{itemize}
}{}{}{}

\gLayout{0}{\gTxGefahren}
{
Es kann zu einer \gE{Loslösung von Teilen des Thrombus} kommen, dieser
losgelöste Teil wird \gT{Embolus}\index{Embolus} genannt. Mit dem Blutstrohm gelangt der
Embolus über die Hohlvenen in das rechte Herz. Bis hierhin gibt es noch keine
Probleme, da die Blutgefäße an die Dicke zunehmen. Vom rechten Herzen gelangt
der Embolus jedoch in den Lungenkreislauf (Lungenarterien). Hier verzweigen
sich die Geäße jedoch wieder und der Durchmesser wird kleiner. Dies führt
dazu, dass der Embolus irgendwann zu groß ist und stecken bleibt. Dadurch wird
nun ein  Lungengefäß verstopft und ein Teil der Lunge nimmt nicht mehr am
Gasaustausch teil: Es entsteht eine \gEs{Lungenembolie}
(\prettyref{topic:lungenembolie}).
}{
\begin{itemize}
    \item Lösung des Thrombus →
    \gEs{Lungenembolie} (\gSee{topic:lungenembolie})
\end{itemize}
}{}{}{}




\begin{gWideParEnv}
\gDias{Bilderserie: Thrombosen und die "`echte Welt"'.}
%1
{\includegraphics[width=\linewidth]{./images/economyclass.jpg}
\gCaptionof{figure}{Die Economyclass. Sorgsam geschlichtet verbringen Menschen
hier Stunden damit Thrombosen zu basteln.
\gSourcePicture{Gabriel} \gAuthNotesPicLegalOk{}{}}}
%2
{\includegraphics[width=\linewidth]{./images/lovenox1.jpg}
\gCaptionof{figure}{Thromboseprophylaxe. \gZ{Niedermolekulares Heparin
(hier "`Lovenox\texttrademark") verhindert Thrombosen die z.B. durch lange
Immobilisation (Reisen, Bettlägrigkeit, Gips, \ldots)
entstehen können. Die Substanz wird unter die Haut ("`subkutan"') gespritzt.
\gSourcePicture{Gabriel} \gAuthNotesPicLegalOk{}{}}}}
%3
{~}
\end{gWideParEnv}






\gMassnahmenNG{Spezielle Maßnahmen}{Beinvenenthrombose}{}{
    \begin{itemize}
        \item Lagerung:  liegend, betroffenes Bein
        hochlagern.
        \item Bewegungsverbot (Emboliegefahr!).
    \end{itemize}
}

\section{Respiratorische Notfälle}


% \begin{center}
%  \marginpar{\includegraphics[width=\linewidth]{./images/noauth/AASdev20090228-img75.png}
%  \gAuthNotesPicLegalNotOk{img75}{Splash Pulmo Notfälle }}
% 
% \end{center}

\minisec{Querverweise}

\begin{itemize}
    \item Lungenentzündung (Pneumonie): \gSee{topic:pneumonie}
    \item Tuberkulose: \gSee{topic:tuberkulose}
\end{itemize}




\subsection{Mechanische Verlegung}

% TODO Fremdkörperaspiration: Erwachsen -- Kind überarbeiten und zusammenführen
% (# 39)

\gLayout{23}{}
{
Durch eine Verlegung der Atemwege ist die Ventilation behindert bzw. aufgehoben.

Spezialfall: \index{Aspiration}\gT{Aspiration} =
Fremdkörper/-säfte (Magensaft, Erbrochenes, Nahrung) gelangt in die
unteren Atemwege
}{}{./images/pharynx-schema.png}
{\gAuthNotesPicLegalNotOk{pharynx-schema}{}Atemwegsverlegung (Schema). }
{}



Besonders gefährdet hinsichtlich einer Atemwegsverlegung
sind:

\begin{itemize}
    \item Kinder
    \item bewusstseinseingeschränkte Personen
    \item Z.n. Bienenstich bei Allergiker (Kehldeckelödem)
\end{itemize}

\paragraph{\gTxSymptome}

\begin{itemize} 
    \item Dyspnoe, Husten
    \item in/exspiratorisches Pfeifen (Stridor)    
    \item Atemstillstand, in weiterer Folge Kreislaufstillstand
\end{itemize}

\gMassnahmenNG{Spezielle Maßnahmen}{Mechanische Atemwegsverlegung}{}{
    \begin{itemize}
        \item Atemwege freimachen
        \item Esmarch-Handgriff
        \item Schlag zwischen Schulterblätter
        \item Heimlich-Griff
        \begin{itemize}
            \item Durch Anwendung des Heimlich-Handgriffes kann es zu
            schweren inneren Verletzungen kommen. Ein Patient ist nach
            Anwendung unbedingt zu hospitalisieren!
        \end{itemize}
        \item \gTxBeiVit \gTxMassVitBes 
        \begin{itemize}
            \item Lagerung: Oberkörper hoch
            \item \ce{O2} 10\gLMin, bei Bedarf bis zu 15\gLMin
        \end{itemize}
    \end{itemize}
    
}



\subsection{Asthma bronchiale}
\label{topic:asthma-bronchiale}

\gDefinition{Chronische Erkrankung mit mehr oder weniger
häufigen Episoden von Anfällen mit massiver Atemnot oder
Atembehinderung durch Verengung der Bronchien.}

\gLayout{0}{\gTxBeschreibung{}}
{
Im akuten Asthmaanfall kommt es zu einer Verengung der
Bronchien und vermehrter Schleimbildung, die \gEs{Ausatmung}
wird massiv erschwert und der Patient hat das Gefühl nicht genug Luft zu bekommen.
}{ Wiederholte Anfälle von Atemnot durch

\begin{itemize}
    \item Verengte Bronchien.
    \item Vermehrte Schleimbildung.
    \item Behinderte Ausatmung.
\end{itemize}
}
{}
{}
{}





Asthma ist bei den meisten Patienten oft schon bekannt und
kann bzw. muss erfragt werden. Viele Patienten haben auch
einen oder mehrere Sprays oder Inhalationsgeräte zur Dauer-
oder Akutbehandlung. 

\gAuthNotes{GABS: Gabe von
Beta-Mimmetikum-Spray muss noch diskutiert werden.

KOCH: nein nur NKA

GABS: wenn, dann NFS (Medi-Liste 1).

Das Problem ist aber: Darf der Patient seinen (vom Arzt
vielleicht gerade für diese Situation
\underbar{verschriebenen}) Spray nehmen? wann darf er und
wann nicht? Macht ja keinen Sinn wenn man dem Pat seinen
Spray generell vorenthält.}

Besonders betroffen vom Asthma bronchiale sind Patienten
mit bekannten Allergien. Folglich kann alles, was bei dem
Patienten eine Allergie verursacht auch einen Asthmaanfall
auslösen (Pollen, Gräser, Katzen, Zigarettenrauch, \ldots). 

Ist bei Eintreffen des Rettungsdienstes noch keine
Besserung der Atemnot eingetreten ist diese Atemnot als
massiv und lebensbedrohend einzuschätzen. Es ist nicht
damit zu rechnen, dass der Zustand von alleine wieder
vergeht und bedarf einer medikamentösen Behandlung durch einen
Notarzt. 

\gCo{Unterschätze nie die Schwere eines Asthmaanfalles! Gefährlich sind
besonders die "`ruhigen"', wenig klagenden Patienten. Asthma-Patienten haben
oft ein anderes Luftnot-Empfinden. Diese Patienten sind
HochrisikoPatienten und machen auch bei der \ce{O2}-Gabe
Probleme! \ce{O2} nur bis max. 95\% Sauerstofsättigung
!\cite{Olschewski2009}\gA{\footnote{Vortrag Univ-Prof.Dr.
Horst Olschewski, 5. AGN-Kongress 2009; Notfall Rettungsmed (2009) 12:322}}}


\gA{Foto Inhaltation V}


\gLayout{0}{\gTxSymptome}
{
Das Leitsymptom ist die Atemnot. Oft ist ein pfeifendes,
brummendes oder giemendes Atemgeräusch besonders während der
Ausatmung (\gE{Exspiration}, \gE{exspiratorischer Stridor})
wahrnehmbar. Die Exspiration ist ausserdem verlängert. Trockerner Husten ist häufig

Der Patient nimmt meist automatisch und unbewusst eine
Position ein die seine Atemhilfsmuskulatur unterstützt: Er
sitzt aufrecht und stützt sich oft nach vorne oder hinten
ab. 

Die Atemfrequenz ist erhöht, oft auch die Herzfrequenz
(Tachypnoe und Tachykardie). Die Sauerstoffsättigung kann
erniedrigt sein. 

\subparagraph{\gCa{Schwerer Anfall}}

Im schweren Anfall ist oft \gEs{kein}, oder nur mehr ein
\gEs{leises Atemgeräusch} hörbar (da kaum mehr
Luft bewegt wird, \gZ{"`silent chest"'}). Es kommt zu einer
Kreislaufdepression, Herzfrequenz und Blutdruck stürzen ab.
Dies kann zu Bewusstseinsstörungen führen. Die
Sauerstoffsättigung ist stark vermindert und der Patient ist zyanotisch.


}{
\begin{itemize}
    \item Akute Atemnot mit trockenem Husten
    \item Exspiratorisches Atemgeräusch (Stridor)
    \item Einsatz der Atemhilfsmuskulatur
    \item Verlängerte Ausatmungsphase
    \item Schwerer Anfall:
    \begin{itemize}
        \item Kein oder leises Atemgeräusch.
        \item HF, RR $\downarrow$
        \item Sp\ce{O2} $\downarrow\downarrow$
        \item Bewusstseinsstörungen
    \end{itemize}
\end{itemize}
}
{}
{}
{}

\gLayout{2}{\gTxKomplikationen}
{
\begin{itemize}
    \item Zyanose
    \item Status asthmaticus
    \item Panik
\end{itemize}
}{

}
{}
{}
{}


\paragraph{Status Asthmaticus}
\index{Status!asthmaticus}\label{topic:status-asthmaticus}Der
"`Status"' ist eine gefürchtete Komplikation des
Asthmaanfalles bei der die Atemnot nicht wieder vergeht. Es
besteht die Gefahr, dass der Patient nach einiger Zeit keine
Kraft mehr zum Atmen hat (Erschöpfung). Der Status
asthmaticus kann tödlich enden.

% \begin{figure}[h]
%     \includegraphics{./images/copd-bronchus.png}
%     
%     \gCaptionof{figure}{Bronchus bei COPD.
%     \gAuthNotesPicLegalNotOk{copd-bronchus}{}}
% \end{figure}



\gMassnahmenNG{Spezielle Maßnahmen}{Akuter Asthmaanfall}{}{
     Einschätzen der vitalen Bedrohung.
     
     \gEs{\gTxBeiVit} \gTxMassVitBes
     
     \begin{itemize}
         \item Lagerung: s.u.
         \item \ce{O2}: s.u.
     \end{itemize}
     
     \gEs{weiters:}
     
     \begin{itemize}
         \item Lagerung: Oberkörper hoch, evt. abstützen
         lassen (Atemhilfsmuskulatur unterstützen)
         \item \ce{O2} bis eine Sauerstoffsättigung von 95\% erreicht ist,
         \gEs{nicht mehr!} Wenn keine Pulsoxymetrie vorhanden ist 8\gLMin.
         
         \gCave{Achtung: In seltenen Fällen
         kann es vorkommen, dass bei hochdosierter
         Sauerstoffgabe der Patient zu atmen aufhört! Siehe >\gAbk{COPD}<
         (\prettyref{topic:copd}). Bei sorgfältiger Überwachung
         des Patienten stellt die \ce{O2}-Gabe in der Praxis
         jedoch kein Problem
         dar\cite{ScholzNotfallmedizin2}.}
         
         \item Vitalwerte erheben \& Monitoring
         \item Beruhigen, voratmen, durch fast
         geschlossene Lippen ausatmen lassen (Lippenbremse)
         \item Sprays, die Patient evt. bei sich hat dürfen
         nicht mehr genommen werden.
     \end{itemize}
    \gAuthNotes{\#14: Einnahme von verordneten Sprays analog
    zu Nitro?\par $O2$-Gabe?}
    
 }

\gLayout{22}{}
{
Patientin mit einem akuten Asthmaanfall.

    Sie sitzt auf einer Treppe und stützt sich nach hinten
    mit den Armen ab. Die Erstmaßnahmen bei vital bedrohten Patienten wurden ergriffen: 
    \begin{compactitem}
            \item Situationsrechte Lagerung
            \item Sauerstoffgabe
            \item
            Reanimationsbereitschaft\footnote{Defibrillator eingeschaltet, sonstiges Zubehör nicht am Foto sichtbar}
            \item Notarztnachforderung\footnote{Paramedic
            vor Ort, am Foto nicht sichtbar.}
            \item Vitalwerte erheben und Monitoring:
            Blutdruck wird gerade gemessen, die Patientin
            ist am Monitor angeschlossen.
     \end{compactitem}
}{

}
{./images/asthma-paramedic.jpg}
{Patientin mit einem akuten Asthmaanfall.\label{fig:asthma-paramedic}}
{Gabriel.}





\subsection{Chronische Bronchitis und 
Chronisch-obstruktive Lungenerkrankung
(COPD)}
\label{topic:copd}
\nocite{Renz-Polster2006} 

% Renz-Polster2006 p442

\emph{Syn.: \gT{Chronic obstructive pulmonary disease}}

\gDefinition{\gT{Chronische Bronchitis}:
Chronisch entzündliche Schleimhautschädigung  der Atemwege.}
\gDefinition{\gT{COPD}: Chronisch-entzündliche
Schleimhautschädigung, welche durch inhalative Noxen
verursacht wird und im Verlauf eine zunehmende obstruktive
Atemwegseinschränkung aufweist.}

Am Anfang steht die chronische Bronchitis, welche durch
Husten mit schleimigem Auswurf gekennzeichnet ist
("`Raucherhusten"'). Es kommt dabei zu einer gesteigerten
Entzündungsantwort auf eingeatmete Stoffe (Zigarettenrauch,
Umweltschadstoffe, \ldots). Wenn man in der
Lungenfunktionsuntersuchung eine Atemwegseinschränkung
nachweisen kann, spricht man von der chronisch-obstruktiven
Lungenerkrankung (\gAbk{COPD}). 


\gLayout{0}{\gTxSymptome}{
Die \gAbk{COPD} ist durch eine voranschreitende
Verschlechterung der Atemleistung gekennzeichnet, erstes
Anzeichen ist ein Husten mit Auswurf, der klassische
"`Raucherhusten"'.

Je nach
Schweregrad kommt es zu Zeichen der Atemwegsverlegung
(Obstruktion) und Ateminsuffizienz:
Die \gE{Ausatmung} (Exspiration) ist erschwert, der Patient klagt
über \gE{Belastungsdyspnoe}, ein Engegefühl, \gE{Husten} und hat Tachypnoe und
evt. Zyanose. Man kann oft ein \gE{brummendes oder giemendes Atemgeräusch}
hören.

Durch die erschwerte Ausatmung kommt es zu einer chronischen Überblähung der
Lungenbläschen und zu einer "`Faß-förmigen"' Verformung des Brustkorbes. Im
Endstadium zeigen sich Zeichen einer Rechtsherzinsuffizienz (\gSee{topic:rechtsherzinsuffizienz}) aufgrund einer Störung im Lungenkreislauf.

Bei fortgeschrittener Ateminsuffizienz bekommt der Patient oft eine
\gE{Heimsauerstofftherapie} verordnet. 
}{
\begin{itemize}
    \item Husten mit Auswurf, anfänglich
    "`Raucherhusten"'.
    \item Voranschreitende, sich verschlechternde \gE{Ateminsuffizienz}:
    \begin{itemize}
        \item Belastungsdyspnoe (bei schwererer \gAbk{COPD}
        auch in Ruhe)
        \item Ausatmung (Exspiration) erschwert
        \item Evt. chronische Zyanose
        \item Brummendes, giemendes Ausatemgeräusch
        \item Benötigt evt. Heimsauerstoff
    \end{itemize}
\end{itemize}
}{}{}{}



\gLayout{0}{Exazerbationen}{
Der \gAbk{COPD}-Patient erleidet öfters plötzliche
Verschlechterungen, sog. \gE{Exazerbationen}. Diese sind
meist durch Infekte ausgelöst und durch vermehrte Atemnot,
Husten und Auswurf gekennzeichnet. 
}{
\begin{itemize}
    \item Plötzliche Verschlechterungen
    \item Meist Infekt-bedingt
\item Vermehrte Atemnot,
Husten und Auswurf   
\end{itemize}
}{}{}{}


% 
% \subparagraph*{Sammelbegriff für}

% \begin{itemize}
%     \item Lungenemphysem (Überblähung der Lungenbläschen)
%     \item chron. obstruktive Bronchitis
% \end{itemize}




\begin{figure}[H]
    
     \includegraphics[width=\linewidth]{./images/hirtler-aass/Alveolen-Bronchus-COPD-edited2.png} 
     \gCaptionof{figure}{\gAuthNotesPicLegalOk{}{}Veränderung der
     Atemwege bei der COPD. \gTabSub{Links:} Schema eines gesunden
     Bronchus und einer gesunden Alveole. \gTabSub{Oben rechts:} Bei der COPD sind
     die kleinen Luftwege verschleimt und verengt. \gTabSub{Unten rechts:} Die
     Lungenbläschen (Alveolen) sind überbläht, weil die Luft nur erschwert wieder entweichen kann.
 \gSourcePicture{Hirtler.}}
\end{figure}

\begin{figure}[h]
    \begin{center}
\includegraphics{./images/copd-stufen.png}
\gCaptionof{figure}{Eine COPD
entsteht nicht plötzlich: Ein COPD{\textquotesingle}ler hat eine
Karriere hinter sich. \gSourcePicture{Quelle unbekannt} \gAuthNotesPicLegalNotOk{copd-stufen.png}{}}
\end{center}
\end{figure}


\gLayout{0}{Probleme mit Sauerstoff bei
COPD-Patienten}
{
\index{Sauerstoff!COPD-Patient}
Normalerweise wird der Atemantrieb über die
\ce{CO2}-Konzentration im Blut gesteuert. Die
\ce{CO2}-Konzentration ist bei \gAbk{COPD}-Patienten jedoch chronisch
hoch, der Körper gewöhnt sich an diesen Zustand. Er
regelt dann den Atemantrieb direkt über den \ce{O2}-Gehalt im Blut.

Wird der \ce{O2}-Gehalt plötzlich künstlich erhöht, fehlt
der \index{COPD!Atemantrieb}Atemantrieb und der Patient kann
\gEs{Bewusstseinsstörungen} ("`\ce{CO2}-Narkose"') und einen
\gEs{Atemstillstand} erleiden.
}{
\gCa{Achtung!}

\gSymbolText
}{}{}{}




\gMassnahmenNG{Spezielle Maßnahmen}{COPD-Exazerbation}{}{
    \begin{itemize}
        \item
        Wenn bei Eintreffen noch keine Besserung
        eingetreten ist, ist von einer vitalen Bedrohung
        auszugehen, \gTxMassVit
        \item \ce{O2}-Gabe
\vspace{-1em}

\gCave{Sauerstoff vorsichtig dosieren,
anfänglich nur 2-3{\gLMin} und Atmung überwachen.
    
    Wenn der Patient bereits Heimsauerstoff benutzt, kann
    man bei Atemnot versuchen 1-2\gLMin{} höher zu dosieren. Die Atmung muss
    dabei überwacht werden! 
    
    Heimsauerstoff soll weiter gegeben werden!}
        
        \item Lagerung: Oberkörper hoch
        \item voratmen, Lippenbremse
        \item eigener Spray darf nicht mehr genommen werden
        
    \end{itemize}
    \gAuthNotes{Spray bei COPD -- Regelung wie bei Nitro??}
}

\gLayout{2}{\gZ{Literatur}}
{
\gZ{
\begin{inparadesc}
  \item[Kasuistiken:] \cite{Knacke200601}
\end{inparadesc}
}
}{}{}{}{}


\subsection{Lungenembolie}\label{topic:lungenembolie}

\gDefinition{Einschwemmen eines Gerinnsels in die  Strombahn des Lungenkreislaufs}


\gLayout{22}{Herkunft des Embolus}
{
Der Embolus, der die Lungenarterien verstopft, entsteht
normalerweise nicht in den Lungengefäßen selbst, sondern er
wird angeschwemmt. Sehr oft löst sich ein Thrombus bei
einer tiefen Beinvenenthrombose und schwimmt im Blut bis in
die Lunge. Daher gelten die
Risikofaktoren der tiefen Beinvenenthrombose
(Immobilisation, lange Reisen, \ldots) auch bei der
Lungenembolie. 
}{

}
{./images/noauth/AASdev20090228-img82.png}
{\gAuthNotesPicLegalNotOk{}{img82}Herkunft des Thrombus.}
{Quelle nicht eruierbar.}



\marginlabel{\begin{itemize}\item Tiefe
(Bein-)Venenthrombose\item Immobilisation (Gips, lange Reise)\item Herz: Bei Rhythmusstörungen\item Tumorpatienten
\end{itemize}}Auch bei Herzrhytmusstörungen wie einem
Vorhofflimmern kann ein Thrombus im Vorhof selber entstehen, 
sich losreißen und Lungengefäße verstopfen. Bei
Gerinnungstörungen oder auch Krebserkrankungen kann das
Gerinnungssystem verrückt spielen und Thromben bilden die
zu einem Embolus werden.

\gLayout{0}{\gTxSymptome}
{
\begin{itemize}
    \item Atemnot, Tachypnoe
    \item Akuter Thoraxschmerz, eher atemabhängig (Atemstopp)
    \item Angst
    \item Tachykardie
\end{itemize}

Eine Lungenembolie kann man leicht mit einem Herzinfarkt oder Angina
pectoris verwechseln. Dies hat aber bei der Erstversorgung kaum
Konsequenz, die Maßnahmen sind die gleichen.


}{
\begin{itemize}
    \item Atemnot, Tachypnoe
    \item Akuter Thoraxschmerz, eher atemabhängig
    \item Angst
    \item Tachkardie
\end{itemize}
}{}{}{}


\gLayout{0}{Typische Befunde}
{
\begin{itemize}
    \item Zeichen einer tiefen Beinvenenthombose (mögliche
    Quelle des Embolus): Bein: einseitig geschwollen,
    überwärmt, rot/rot-bläulich verfärbt, schmerzhaft
    \item evt. gestaute Halsvenen
\end{itemize}
}{

}{}{}{}


\gLayout{0}{Typische Anamnese}
{
\begin{itemize}
    \item Immobilisation oder Operation innerhalb
    der letzten 4 Wochen (\gEs{Lange Reise}, Flug, Auto,
    Bahn, Bettlägrigkeit, Gips, {\dots}
    \item \gZ{Schwangerschaft}
    \item \gZ{Raucher(in)}
    \item \gZ{Verhütungspille}
\end{itemize}
}{
\begin{itemize}
    \item Immobilisation
    \item \gZ{Schwangerschaft}
    \item \gZ{Raucher(in)}
    \item \gZ{Verhütungspille}
\end{itemize}
}{}{}{}


\gLayout{2}{Wann ist die Diagnose \emph{"`Lungenembolie"'}
 wahrscheinlich?}
{
nach \cite{Wells2000}

\begin{itemize}
    \item {\bfseries Klinischer Verdacht} basierend auf
    Erfahrungswerten
    \item {\bfseries Andere Diagnosen sind weniger
    wahrscheinlich}
    \item Herzfrequenz \textgreater 100\,\nicefrac{}{min}
    (Tachykardie)
    \item Immobilisation oder Operation \gZ{innerhalb
    der letzten 4 Wochen} (\gEs{Lange Reise}, Flug, Auto,
    Bahn, Bettlägrigkeit, Gips, {\dots}
    \item Bereits einmal durchgemachte tiefe
    Beinvenenthrombose oder Lungenembolie
    \item Bluthusten
    \item Krebserkrankung \gZ{innerhalb der letzten 6 Monate}
\end{itemize}

Die ersten beiden, fett gedruckten, Kriterien sind die
wichtigsten Kriterien. 
}{

}{}{}{}

\paragraph{} 

\gMassnahmenNG{Spezielle Maßnahmen}{Lungenembolie}{}{
\begin{itemize}
    \item \gTxMassVitBes
    \begin{itemize}
        \item Lagerung: Oberkörper hoch
        \item \ce{O2} 10\gLMin, bei Bedarf bis zu 15\gLMin
    \end{itemize}
    \begin{itemize}
        \item Striktes Bewegungsverbot
        \item Beengende Kleidung öffnen
    \end{itemize}
\end{itemize}
}


\subsection{Lungenödem}\label{topic:lungenoedem}

\gDefinition{Flüssigkeitsansammlung durch Stau in den
Alveolen und im Zwischengewebsraum der Lunge.}
\nopagebreak

\gLayout{2}{Ursachen}
{
\begin{itemize}
    \item Blutrrückstau: akute Linksherzinsuffizienz,
    Linksherzversagen
    \item z.B. bei MCI, zusätzliche Belastung bei vorgeschädigtem Herz
    (Infekt)
    \item Durch Entzündung: Inhalation toxischer Substanzen
\end{itemize}
\gCave{Das gefährliche am Lungenödem sind vor allem dessen Ursachen!}
}{

}{}{}{}



\gLayout{2}{\gTxSymptome}
{
\begin{itemize}
    \item Atemnot
    
    \item Zyanose
    
    \item \index{Brodelnde Atemgeräusche}Brodelnde Atemgeräusche, oft
    ohne Hilfsmittel hörbar, im Extremfall schon von der Wohnungstür
    aus.
    
    \begin{center}

\gT{Brodelnde Atemgeräusche}: Als ob jemand mit einem
Strohhalm in ein Glas Wasser bläst.
\end{center}

\end{itemize}
}{

}{}{}{}



\gMassnahmenNG{Spezielle Maßnahmen}{Lungenödem}{}{
    \begin{itemize}
      \item \gTxMassVitBes
      \begin{itemize}
        \item Lagerung: Oberkörper hoch, \gE{wenn möglich Beine
        herabhängen lassen}
        \item \ce{O2} 10\gLMin, bei Bedarf bis zu 15\gLMin
        \item  Notarzt
        \item Reanimationsbereitschaft
      \end{itemize}
      \item Striktes Bewegungsverbot
      \item Beengende Kleidung öffnen
      \item \gZ{evt. Gabe von Nitroglycerin-Spray durch einen
      Notfallsanitäter}
    \end{itemize} 
}

Patienten mit einem Lungenödem (v.a. bei kardialer Ursache) sind
\gEs{lebensbedrohlich krank}. Das Herz verbraucht oft
schon die allerletzten Reserven, schon die Aufregung durch den
Transport im Stiegenhaus zum Auto kann das Herz endgültig in den
Herz-Kreislaufstillstand treiben. 

Es kommt häufig vor, dass diese Patienten \gEs{im
Stiegenhaus} oder \gEs{im Auto}
\gEs{reanimationspflichtig} werden. 

\gCave{Kein Transport ohne Notarzt! 

Auch wenn das Spital
{\quotedblbase}um{\textquotesingle}s Eck{\textquotedblleft}
ist. Der Patient schafft es möglicherweise nicht 'mal mehr
ins Auto!}



\subsection{Hyperventilationstetanie}
\label{topic:hyperventilationstetanie}
\index{Hyperventilation!-stetanie}
\index{Hyperventilation!psychisch bedingte --}



\gLayout{0}{\gTxBeschreibung}{
Durch die Atmung wird der Säure-Basen-Haushalt
des Körpers wesentlich beeinflusst (\gSee{topic:saeure-basen-haushalt}). Atmet
ein sonst gesunder Mensch unnatürlich tief und schnell, so wird sein Blut-pH
sehr rasch basisch. Dadurch geraten die Elektrolyte im Blut in ein
Ungleichgewicht, es kommt zu Symptomen. 
}{
\begin{itemize}
    \item Tiefe und schnelle Atmung bei sonst gesundem Menschen
    \item Blut-pH wird basisch $\rightarrow$

\end{itemize}
}{}{}{}

\gLayout{0}{\gTxSymptome}
{
Die Patienten beklagen \gEs{Atemnot}, obwohl die \ce{O2}-Zufuhr nicht
beeinträchtigt und das Blut maximal gesättigt ist (Sp\ce{O2} von 100\,\%). Sie
haben eine \gEs{tiefe und schnelle Atmung} (\gE{Tachypnoe}).

Durch das Elektrolyt-Ungleichgewicht kommt es zuerst zu Schwindelgefühlen, etwas
später zu einem Kribbeln in den Fingern. In weiterer Folge zeigen sich tonische
Krämpfe, klassisch ist die sog. \gT{Pfötchenstellung}\index{Pfötchenstellung}:
Die Arme sind angezogen und Handhaltung erinnert an Pfoten eines Tieres. Wird
die Hyperventilation nicht rechtzeitig beendet, kann es zur
Eintrübung bis hin zur Bewusstlosigkeit kommen.

 }{
\begin{itemize}
    \item Subjektiv: \emph{"`Atemnot"'}, Sp\ce{O2} von 100\,\%
    \item Tiefe, schnelle Atmung
    \item Schwindel, Bewusstseinsstörungen
    \item Kribbeln in den Fingern
    \item Tonische Krämpfe ("`Pfötchenstellung"')
\end{itemize}
}{}{}{}

\gLayout{0}{Ursachen}
{
Die Hyperventilationstetanie ist meistens psychisch bedingt und meist Folge
einer \gEs{psychischen Belastungsreaktion} bzw. Folge einer Panikattacke. Auch
im Rahmen von (Pop-)Konzerten und ähnlichen Veranstaltungen sind oft Patienten mit
Hyperventilationstetanie anzutreffen.

}{
\begin{itemize}
  \item Meist psychisch bedingt 
  \begin{itemize}
      \item Psych. Belastungsreaktion, Panikattacke
      \item Veranstaltungen
%   \item Lungenembolie
%   \item diabetische Stoffwechselentgleisung (Ketoazidose, kompensatorisch)
\end{itemize}
\end{itemize}
}{}{}{}

\gLayout{0}{Differentialdiagnosen}
{
Die psychisch bedingte Hyperventilation muß unbedingt von der Hyperventilation
aufgrund anderer Ursachen abgegrenzt werden. Bei \gEs{Ateminsuffizienz} oder
\gEs{durch Erkrankungen verursachte Atemnot} kommt es natürlich auch oft zu
einer (kompensatorischen) Hyperventilation (Herzinfarkt, Lungenembolie, Pneumothorax,
Asthma-Anfall, \ldots). 

Bei einer stoffwechselbedingten \gEs{Übersäuerung} (Azidose) kommt es zu einer
Hyperventilation um \ce{CO2} abzuatmen (dadurch wird Säure im Körper abgebaut). Die \gE{Kussmaul'sche Atmung} beim \gEs{diabetischen Koma} wäre ein typisches
Beispiel dafür (\gSee{topic:koma-diabetisches}).

}{
\begin{itemize}
  \item Atemnot mit körperlicher Ursache 
  \item Ateminsuffizienz 
  \item Stoffwechselbedingte Azidose (z.B. Diabetisches Koma)
\end{itemize}
}{}{}{}


% \paragraph{\gTxSymptome}
% 
% \begin{itemize}
% \item Unruhe, Aufregung, evt. Angstgefühle, Panik
% \item Tachypnoe mit Alkalose $\rightarrow$ durch Änderung
% des pH-Werts Änderung des Kalziumspiegels im Blut, dadurch
% \begin{itemize}
%     \item Kribbeln, 
% \item Krämpfe (Gesichtsmuskulatur),  Pfötchenstellung
% \end{itemize}
% \end{itemize}

\gMassnahmenNG{Spezielle Maßnahmen}{Hyperventilationstetanie}{}{
\begin{itemize}
  \item Psychische Betreuung: Patient beruhigen, mit ruhiger, langsamer Stimme
  mit dem Patienten sprechen.
  \item Kommandoatmung.
  \item \ce{CO2}-Rückatmung: In ein Plastiksackerl über kurze Zeit Ein- und
  Ausatmen lassen. Ebenso eignet sich eine \gEs{\ce{O2}-Maske mit Reservoir}. 
  Handschuh)
  \item Grundsätzlich sollte kein \ce{O2}  verabreicht werden. Aus
  psychologischen Gründen kann es jedoch sinnvoll sein \gEs{\ce{O2} in geringer
  Dosierung (1\gLMin{})} über eine \ce{O2}-Maske zu
  verabreichen.
  \item Wenn nicht erfolgreich und Krämpfe bestehen bleiben bzw.
  Bewusstseinsstörungen auftreten: \gTxMassVit .
\end{itemize}
}





\clearpage
\section{Neurologische Erkrankungen und Notfälle}

\gDefinition{Erkrankungen, Notfälle und sonstige Störungen
die das Nervensystem betreffen. Neurologische Störungen
gehen oft mit einer Störung des Bewusstseins oder Schmerzen einher.}

Eine Übersicht und willkürliche Einteilung:

\begin{description}
    \item [Schlaganfall]
    \item [Vergiftungen] können fast alles machen, und neurologische Symptome erst recht
    \item [Schädel-Hirn-Trauma] Abkz. 'SHT'; Ein Schlag auf
    den Kopf erhöht \emph{nicht} das Denkvermögen \ldots
    \item [Stoffwechselstörungen] welche Einfluss auf das Nervensystem nehmen, z.B. Zuckerentgleisungen 
    \item [Infektionen] können das Nervensystem beeinflussen, z.B. Verwirrtheit bei Fieber
    \item [Alter] beeinflusst auch das Nervensystem, im Rahmen einer Demenz kommt es z.B. zum Abbau der Denkleistung
    \item [Krampfanfälle] können vom zentralen Nervensystem ausgehen
    \item [Erkrankungen des Stützapparates] Aufgrund der räumlichen Nähe zum Stütz- und Bewegungsapparates kann das Nervensystem in Mitleidenschaft gezogen werden, z.B. bei Bandscheibenvorfällen
\end{description}

\minisec{Querverweise}

\begin{itemize}
    \item Infektiöse (bakterielle) Meningitis:
    \gSee{topic:meningitis}
\end{itemize}


\subsection{Schlaganfall, Apoplektischer Insult, Hirninfarkt}
\index{Schlaganfall}
\index{Insult}
\label{topic:tia}

\emph{Syn.: \index{Insult}\index{Apoplektischer
Insult}Apoplektischer Insult, \index{Apoplexie}Apoplexie}

\gDefinition{siehe Arten.}

\gLayout{0}{Arten}
{
Grundsätzlich gibt es 2 Formen, die aber vor Ort nicht
unterschieden werden können:
\begin{labeling}[~--]{xxxxxxxxxxxxxxxxxx}
      \item[\gT{Trockener Insult}] Verschluss eines Hirngefäßes mit anschließendem
    Absterben der versorgten Hirnregion ({\quotedblbase}trockener
    Insult{\textquotedblleft}), oder
    \item[\gT{Blutiger Insult}] Hirnblutung, z.B. durch Platzen eines
    Hirngefäßes, kann spontan oder aufgrund eines Traumas auftreten.
\end{labeling}



}
{
\begin{itemize}
    \item "`Trockener"' Insult
    \item "`Blutiger"' Insult
\end{itemize}

}{}{}{}

\gLayout{0}{\gT{TIA}}
{\index{TIA}
\index{Transitorisch-Ischämische Attacke}
TIA ist die Abkürzung für Transitorisch-Ischämische Attacke. Sie ist eine Sonderform des
Schlaganfalls bei der die Symptome innerhalb von 24 Stunden vergehen und kein
Hirngewebe abstirbt.
    % TODO INHALT: TIA: Ischämie erklären (evt. in EIN)
}
{
% TODO MED: Slide fehlt!
\begin{itemize}
    \item Transitorisch-Ischämische Attacke.
    \item Symptome <\,24\,h.
\end{itemize}

}{}{}{}

% \gDias
% {
%     \includegraphics[width=\linewidth]{./images/noauth/AASdev20090228-img83.png}
%     \gCaptionof{figure}{\gAuthNotesPicLegalNotOk{img83}{}Insult: Eine spastische
%     Halbseitenlähmung kann sich als Spätfolge Wochen nach einem durchgemachten
%     Schlaganfall e ntwickeln.
%     \gSourcePicture{Quelle nicht eruierbar.}}
% }{
%     \includegraphics[width=\linewidth]{./images/noauth/AASdev20090228-img84.png}
%     \gCaptionof{figure}{\gAuthNotesPicLegalNotOk{img84}{}Insult:
%     Halbseitenlähmung im Gesicht.
%     \gSourcePicture{Quelle nicht eruierbar.}}
% }{}


\gLayout{0}{\gTxSymptome}
{
Die Symptome sind davon abhängig welche Hirnregion betroffen ist und 
können sehr unterschiedlich sein und umfassen:

%\begin{multicols}{2}
\begin{labeling}[~--]{xxxxxxxxxxxxxxxxxx}
    \item[Motorik] Störungen der Kraft und Beweglichkeit sind sehr häufig:

\gT{Halbseitenzeichen}: Aufgrund des Ausfalls eines Zentrums im
        Gehirn kann es zu einer \gE{Kraftminderung} einer Seite kommen, dies kann bis
        zu einer \gE{schlaffen, kompletten Lähmung} führen.
        
        Ein \gE{herabhängender Mundwinkel} ist ein Zeichen einer Lähmung
        der Gesichtsmuskulatur. 
        
        Durch die Lähmungen kann es zu \gEs{Schluckstörungen} kommen, es
        besteht dann eine erhöhte Aspirationsgefahr.
        
        Der \gT{Herdblick} ist eine "`Blickdiagnose"': Es kann zu Blicklähmungen
        kommen, der Patient blickt in Richtung des "`Herd des Geschehens"'.

    \item[Sprache] Störungen der Sprache sind ebenfalls sehr
    häufig

 Oft nimmt man eine \gEs{verwaschene Sprache} wahr, welche (fälschlich) an
 Trunkenheit erinnert.
 
Weiters kann es zu einer "`wirren Sprache"' kommen: Worte fehlen dabei oder
Silben werden vertauscht. Im Extremfall kann die Sprache gänzlich "`verloren"'
gehen. Auf der anderen Seite kann die Sprache zwar erhalten, aber das
Sprach\gE{verständnis} geschädigt sein

    \item[Wahrnehmung] 
    Es kann zu Wahrnehmungstörungen wie z.B. \gEs{Sehstörungen}, plötzliche
    Erblindung, Doppelbilder, Gesichtsfeldausfälle kommen.
    
    \item[Bewusstsein] Störungen des Bewusstseins sind gefährlich und müssen bei
    Insult-Patienten immer erwartet werden. Der Bewusstseinszustand kann sich
    auch plötzlich verschlechtern.
    
    \item[Allgemein] Weiters können unspezifische Symptome wie z.B.
    Schwindel und Kopfschmerzen auftreten.
    Ein plötzlicher, \gEs{donnerschlagartiger Kopfschmerz} ist ein klassisches
    Symptom für eine spontane Hirnblutung.
    
\end{labeling}
%\end{multicols}
}
{
Je nach betroffener Hirnregion sehr unterschiedlich.
\begin{itemize}
    \item Halbseitenzeichen (Schwäche, Lähmung)
    \item Herdblick
    \item Verwaschene Sprache
    \item Sehstörungen
    \item Bewusstseinsstörungen
    \item Kopfschmerzen
\end{itemize}
}{}{}{}


\gLayout{0}{Differentialdiagnose}
{
Die wichtigste Differntialdiagnose ist eine Störung des
Blutzucker-Stoffwechsels. Sowohl die Hyper(${\uparrow}$)-, als auch die
Hypoglykämie (${\downarrow}$) können einen Schlaganfall vortäuschen. Bei
neurologischen Symptomen muß der Blutzucker immer gemessen werden!
Ebenso ist immer an Vergiftungen zu denken.

Die Abgrenzung zu anderen neurologischen Erkrankungen ist oft schwierig bzw.
kaum möglich.
}
{
\begin{itemize}
    \item Blutzucker ${\downarrow}$, ${\uparrow}$.
    \item Sonstige neurologische Erkrankung.
    \item Vergiftungen
\end{itemize}

}{}{}{}


\gMassnahmenNG{Spezielle Maßnahmen}{Insult}{}{
    
    \begin{itemize}
        \item \gTxMassVitBes 
        \begin{itemize}
            \item Lagerung: \gE{leicht erhöhter Oberkörper}
            \item \ce{O2} 10\gLMin, bei Bedarf bis zu 15\gLMin
            % ERRATUM 0.8.0 Nachberunfung des Notarztes beim Schlaganfall
            
            \item Notarzt: Grundsätzlich ist ein Schlaganfall-Patient aufgrund seiner Diagnose als vital bedroht einzustufen. Es gibt jedoch eine
            Besonderheit: Eine definitive Therapie in einer Spezialabteilung
            (Stroke Unit) kann andere, noch nicht abgestorbene Teile des Hirns
            vor weiterem Schaden bewahren, d.h. es ist sehr wichtig, dass diese
            Therapie möglichst schnell erfolgt.    
            
            Daher soll von einer
            \gEs{Nachberufung eines Notarztes \underline{abgesehen} werden, 
            \gE{wenn folgende Umstände gegeben sind}}:
            \begin{itemize}
              \item Die Verdachtsdiagnose gilt als \gE{sehr sicher}.
              \item Der Patient ist sonst als \gE{stabil} eingestuft
              (\prettyref{m:ersteinschaetzung}) und
              bewusstseinsklar\footnote{Bei der Einschätzung des Bewusstseins
              darf man sich nicht durch die durch den Schlaganfall
              verursachten Symptome (Sprachstörungen, Blickstarre \ldots)
              beirren lassen.}.
              \item Es gibt sonst \gE{keine anderen Umstände} die die Berufung
              eines Notarztes erforderlich machen. 
              \item Es muss versucht werden, eine Aufnahme auf eine
              \gE{Spezialabteilung} für Schlaganfälle zu
              organisieren\footnote{Der Erfolg ist abhängig von der
              Verfügbarkeit eines entsprechenden freien Bettes.} (Stroke Unit,
              Kontaktaufnahme mit der Leitstelle).
            \end{itemize}
            \gA{HART: Muss es unbedingt eine Spezialabteilung sein?
            "`bewusstseinsklar"' könnte missverstanden werden.}
            \gA{\par Diskussion: }
            \gACh{KOCH}{2010-02-13}{ok.}
            \gACh{HART}{2010-02-13}{Muss es unbedingt eine Spezialabteilung sein?
            "`bewusstseinsklar"' könnte missverstanden werden.}
            \gACh{GABS}{2010-02-13}{Fußnote wegen Bewusstseinseinschätzung
            eingefügt. Spezialabteilung: muss nur \underline{versucht} werden
            eine Stroke Unit zu organisieren, je nach Verfügbarkeit.}
            \gACh{GABS/STEU}{2010-04-06}{Ticket 38:
            Besprechung gestern mit STEU: kann so übernommen werden, sieht
            allerdings Problem wegen Inkonsistenz und für den Nutzer
            möglicherweise verwirrender wenn-dann-Verzweigung. Bei Hirnblutung
            ist das Vorgehen ein Nachteil. (Einwand GABS: Ischämisch häufiger,
            Abwägung Benefit Narkose, etc. vs. schnelles CT)}
        \end{itemize}
        \item \gEs{Neurocheck} inkl. Messung des \gEs{Blutzuckers}!
        \item Hospitalisation: Wenn verfügbar: \gEs{Stroke Unit}, sonst
        Abklärung mit der Leitstelle wegen Alternativen (Abt. für
        Neurologie oder Innere Medizin, mit Möglichkeit zur
        Intensivüberwachung).
    \end{itemize}
}

\subsection{Vom Gehirn ausgehende Krämpfe --- Zerebrale Krampfanfälle}
\label{topic:zerebrale-krampfanfaelle}
\label{topic:krampfanfall}

\gDefinition{Vom Gehirn ausgehende motorische Krampfanfälle aufgrund von akuten
oder chronischen Erkrankungen des Gehirns. Dabei entladen sich unkontrolliert Nervenzell-Gruppen im Gehirn.}


\gLayout{0}{Einteilung}
{
Die zerebralen Krampfanfälle kann man auf verschiedene Arten einteilen

\begin{enumerate}
    \item Nach Art:
    \begin{labeling}[~--]{Generalisierter Anfall }
          \item[Tonisch] "`Verkrampfung"', erhöhter
        Muskeltonus
        \item[Klonisch] Zuckungen
        \item[Tonisch-klonisch] Kombination aus beidem, das
        klassische Erscheinungsbild eines Krampfanfalles
    \end{labeling}
    \item Nach Lokalisation:
    
    \begin{labeling}[~--]{Generalisierter Anfall }
        \item [Fokaler Anfall] Der Krampfanfall ist auf
        eine Körperregion beschränkt, siehe
        \ref{topic:krampfanfall-fokaler}
        
        \item [Generalisierter Anfall] Er betrifft den
        ganzen Körper, siehe
        \ref{topic:krampfanfall-generalisierter}
        
    \end{labeling}
    \end{enumerate}
}{
Die zerebralen Krampfanfälle kann man auf verschiedene Arten einteilen

\begin{itemize}
    \item Nach Art:
    \begin{description}
        \item Tonisch
        \item Klonisch
        \item Tonisch-klonisch
    \end{description}
    \item Nach Lokalisation:
    \begin{description}
        \item Fokaler Anfall
        \item Generalisierter Anfall
    \end{description}
    \end{itemize}
}{}{}{}

\gLayout{0}{Ursachen}
{
Für zerebrale Anfälle gibt es eine Reihe von unterschiedlichen Ursachen die
direkt oder indirekt das Gehirn betreffen. Einige betreffen die Hirn- und
Nervenstruktur, andere sorgen dafür, dass das Hirn zu wenig
Sauerstoff und andere Nährstoffe erhält. Tafel
\ref{tab:krampfanfaelle-ursachen} zeigt die häufigsten Ursachen und
Grunderkrankungen. 
}{
Tafel \ref{tab:krampfanfaelle-ursachen}.
}{}{}{}

\begin{labeling}{xxxxxxxxxxxxxxx}
  \item[Epilepsie] Die Epilepsie ist eine chronische Erkrankung welche durch wiederkehrende
  Episoden von zerebralen Krampfanfällen gekennzeichnet ist.
  \item[Vergiftungen] (Medikamente/Drogen/...) können wieder fast alles machen,
  also auch Krampfanfälle. \cite{Supady2009}
  \item[Hirntumor] Durch den Tumor wird das Hirn beschädigt, auch dadurch kann es
  zu Anfällen kommen.
  \item[Narbengewebe] nach Insult oder Hirnverletzung.
  \item[Hirnblutung]
  \item[Sauerstoffmangel] z.B. Kollaps, Ischämie aufgrund eines Gefäßverschlusses
  \item[BZ] Der Dauerbrenner. Eine Blutzuckerentgleisung kann fast alle
  neurologischen Symptome vortäuschen!
  \item[angeboren] ohne organisch fassbare  Ursache
  \item[Fieberkrampf] Fieber bei Kleinkindern
  \item[Infektionen]
  \item[Entzug] Entzug von Alkohol oder anderen Drogen
\end{labeling}




\subsubsection{Fokaler Anfall}
\index{Krampfanfall!fokaler}
\label{topic:krampfanfall-fokaler}
Der Krampfanfall ist auf eine Körperregion beschränkt, z.B.
auf eine Hand. Der Patient kann bei Bewusstsein sein. Ein
fokaler Anfall kann sich ausweiten und in einen
generalisierten Anfall übergehen. \gZ{Der fokale Anfall wird
auch \emph{"`petit mal"'} (\emph{franz.:} "`kleines
Unglück"') genannt.}
      


\subsubsection{Generalisierter Anfall}
\index{Krampfanfall!generalisierter}
\label{topic:krampfanfall-generalisierter}

\gLayout{0}{Beginn}
{
Er betrifft den ganzen Körper. Dem Anfall kann ein
\gE{Dämmerzustand} 
vorausgehen, \gZ{eine sog. "`Absence"'}. Der Patient wirkt
dabei plötzlich abwesend und \emph{"`irgendwie komisch"'}.
Dann verliert er das Bewusstsein und die Muskeln werden
anfangs schlaff. Wenn der Patient gestanden ist, fällt er
ungebremst auf den Boden (Cave: \gEs{Verletzungen!}).
Manchmal kommt es vor dass der Patient einen
Initialschrei\index{Initialschrei} von sich gibt, der
gruselig klingen kann. 
}{
\begin{itemize}
    \item Dämmerzustand -- Absence
    \item Initialschrei
    \item evt. Sturz
    \item evt. Folgeverletzungen
\end{itemize}

}{}{}{}

\gLayout{0}{Krampf}
{
Dann kommt es zu
\gEs{tonisch-klonischen Krämpfen},
die den ganzen Körper betreffen. Während des Anfalls ist der
Patient nicht bei Bewusstsein, es besteht außerdem i.d.R.
ein Atemstillstand. Er kann jedoch starke Kräfte entwickeln, es
besteht massive
Verletzungsgefahr.
}{
\begin{itemize}
    \item Tonisch-klonische Krämpfe
    \item Verletzungsgefahr
\end{itemize}
}{}{}{}


\gLayout{0}{Danach}
{
Der Krampf dauert meist wenige Minuten und hört von selbst
auf. Der Patient ist danach erstmal bewusstlos, im weiteren
Verlauf mit mehr oder weniger großem Aufwand erweckbar. Er
befindet sich dann in der sogenannten
\gT{Nachschlafphase}\index{Nachschlafphase}.

Der Patient \gE{klart im Verlauf mehr und mehr auf}. Es ist
wichtig zu versuchen mit dem Patienten in Kontakt zu treten
um seinen Wachheitsgrad und dessen
\gEs{Verlauf} beurteilen zu können. Für den Anfall selbst besteht
Amnesie! Daher muss man den Patienten aufklären was passiert
ist, wo er sich befindet, etc. Ein \gEs{Neuro-} und
\gEs{Traumacheck} ist nach \underline{\gE{jedem}} Anfall
unbedingt notwendig! 
}{
\begin{itemize}
    \item Nachschlafphase
    \item Aufklaren
    \item Amnesie
\item \gEs{Neurocheck}
\item \gEs{Traumacheck} nach \underline{\gE{jedem}} Krampfanfall    
\end{itemize}
}{}{}{}


\gLayout{0}{Sonstige Besonderheiten}
{
Häufig kommt es während des Anfalles zu einem Zungenbiss,
außerdem zum Einnässen oder Einkoten.
}{
\begin{itemize}
    \item Zungenbiss
    \item Einnässen
    \item Einkoten
\end{itemize}
}{}{}{}


\gZ{Der generalisierte Anfall wird auch \emph{"`grand mal"'} (\emph{franz.:} "`großes Unglück"') genannt.}

\subsubsection{Besondere Krampfanfälle}

\gLayout{0}{Fieberkrampf}
{
\label{topic:fieberkrampf}
\index{Fieberkrampf}
\index{Krampfanfall!Fieberkrampf}

Krampfanfall bei Kleinkindern bei schnellem, hohen Fieberanstieg
({\textgreater}\,39{\textcelsius}). Der Fieberkrampf sieht aus wie
generalisierter grand mal-Anfall, kommt allerdings nur bei Kindern {\textless} 4
Jahre vor. 
}{
\begin{itemize}
    \item Kinder {\textless} 4 Jahre.
    \item Hohes Fieber.
    \item Wie generalisierter Anfall.
\end{itemize}
}{}{}{}



\gLayout{0}{Eklamptischer Krampfanfall während der
Schwangerschaft}
{
In Folge einer Präeklampsie (EPH-Gestose)
kann es zu einem sog. \gE{eklamptischen Krampfanfall}
kommen. Dieser wird unter \gSee{topic:eklampsie}
besprochen.
}{
\begin{itemize}
    \item Folge einer Präemklampsie.
    \item \gSee{topic:eklampsie}.
\end{itemize}
}{}{}{}


\subsubsection{Status Epilepticus}
\index{Status!epilepticus}
\gLayout{0}{Lebensgefahr}
{
Der Status Epilepticus ist eine
lebensgefährliche Steigerung
des generalisierten Krampfanfalls. Er besteht dann, wenn der
Anfall nicht nach kurzer Zeit vergeht oder innerhalb kurzer
Abstände mehrere Anfälle auftreten (z. B. 2 Anfälle
innerhalb von 24 Stunden). Der Krampf muss dann mit Medikamenten
durchbrochen werden.
}
{
\begin{itemize}
    \item Lange Dauer.
    \item Wiederholtes Auftreten (>\,2\,/\,24\,h).
    \item Akute \gCa{Lebensgefahr!}
\end{itemize}
}{}{}{}


\gFazit{Im Rettungsdienst gilt praktisch jeder Krampfanfall,
den das Personal selber miterlebt, als Status-Verdächtig.}

\subsubsection{Maßnahmen}

\gMassnahmenNG{Spezielle Maßnahmen}{Krampfender Patient}{}{
\begin{itemize}
    \item \gTxMassVitBes
    \begin{itemize}
        \item Lagerung: sichern
        \begin{itemize}
            \item Patient vor Verletzungen schützen
            \item Möbel etc. entfernen oder polstern
            \item NICHT festhalten -- aber sichern, Eigenschutz
            beachten!\footnote{\textbf{Sichern}: Patient darf nicht von der
            Trage oder aus dem Bett fallen, der Kopf darf nicht am Boden oder Tischbein anschlagen, etc.}
        \end{itemize}
    \end{itemize}
    \item nach Anfall BAK-Schema, Atemwege freihalten
    \item stabile Seitenlage  (Erholungsschlaf wie  Bewusstlosigkeit)
    \item Besondere Diagnostik:
    \begin{itemize}
        \item Neurocheck inkl. Messung des \gEs{Blutzuckers}!
        \item Traumacheck
    \end{itemize}
    \item Reizabgeschirmter Transport
    \item Hospitalisierung: Abt.f. Innere Medizin. (evt. auch
    Intensivüberwachung)
\end{itemize}
}

\gMassnahmenNG{Spezielle Maßnahmen}{Nicht-(mehr) krampfender Patient}{}{
\begin{itemize}
    \item Umstände klären und entscheiden ob \gEs{Notarzt}
    notwendig ist:
    \begin{itemize}
        \item 2 oder mehr Krampfanfälle innerhalb von 24 Stunden
        \item Anfälle "`erheblich anders"' als sonst üblich
        \item Erster Krampfanfall überhaupt
    \end{itemize}
    \item Besondere Diagnostik:
    \begin{itemize}
        \item Neurocheck inkl. Messung des \gEs{Blutzuckers}!
        \item Traumacheck
    \end{itemize}
    \item Reizabgeschirmter Transport
    \item Hospitalisierung: Abt.f. Innere Medizin.
\end{itemize}
}



\subsection{Lumbago, Lumbo-ischialgie und Bandscheibenvorfall} 
\gZ{(Discusprolaps)} 

Unter Ischialgie (\gT{Lumbalgie}, Lumboischialgie,
\emph{"`\gT{Lumbago}"'}) versteht man Schmerzen im
Versorgungsbereich des Ischias-Nerv, welche u. a. durch eine Entzündung oder durch Reizung bzw. Kompression des Nervs oder seiner
Wurzeln verursacht werden können. 

\paragraph{Ursachen} können eine Reizung bzw. Druck
im Bereich der Lendenwirbel bzw. der Kreuzbeinwirbel,
ein Bandscheibenvorfall, sowie andere Erkrankungen der Wirbelsäule, aber auch ein Unfall sein. 

\paragraph{\gTxSymptome} sind Schmerzen in der
Lendengegend, die in das betroffene Bein bis zum Fußaußenrand ausstrahlen, 
Schonhaltung, lokale Druck- und Klopfempfindlichkeit der Wirbelsäule,
Sensibilitätsstörungen und eventuell Lähmung der Zehenmuskulatur. Je nach
Schwere und zeitlicher Entwicklung (langsames/plötzliches Auftreten) bezeichnet
man das Zustandsbild als Lumboischialgie oder Lumbago (Hexenschuss). 
Gegebenenfalls ist im Krankenhaus ein Bandscheibenvorfall abzuklären.




\clearpage
\section{Metabolische Störungen (Stoffwechselstörungen)}

Die relevanteste Stoffwechselstörung im Rettungsdienst ist der
\index{Diabetes mellitus}\gT{Diabetes mellitus}. Dieser
stellt eine Störung des Kohlehydrat-(zucker-)stoffwechsels dar.
\gEs{Zucker ist ein Kohlehydrat}. 

\subsection{Diabetes Mellitus}
\label{topic:diabetes}
\label{topic:diabetes-mellitus}

\gDefinition{Diabetes mellitus ist eine chronische
Stoffwechselerkrankungen bei der die Blutzuckerregulation gestört ist.} 



\gLayout{22}{Der Blutzuckerwert}
{~

\begin{center}
\begin{tabular}{cccl}
\toprule
\\\midrule
0 & -- & 20 & Lebensgefahr
\\\addlinespace
21 & -- & 60 & stark erniedrigt, Bewusstseinsstörungen
\\\addlinespace
61 & -- & 80 & erniedrigt, leichte Symptome
\\\addlinespace
81 & -- & 120 & Nüchtern-Blutzucker
\\\addlinespace
121 & -- & 200 & erhöht
\\\addlinespace
201 & -- & 399 & stark erhöht
\\\addlinespace
400 & -- & ?? & sehr stark erhöht
\\\addlinespace\bottomrule
\end{tabular}
\end{center}

Durch die Erkrankung oder die unsachgemäße Behandlung kann es zu 
folgenden Notfällen kommen:

\begin{description}
\item[\gT{Hyp\underline{\gE{o}}glykämie}] \ldots "`Unterzucker"'

\item[\gT{Hyp\underline{\gE{er}}glykämie}] \ldots "`Überzucker"'

\end{description}

}
{}
{./images/noauth/AASdev20090228-img88.png}
{\gAuthNotesPicLegalNotOk{img88}{}Interpretation des
BZ-Wertes: Richtwerte.}{Quelle nicht eruierbar.}

\gLayout{2}{Insulin}
{
\gT{Insulin} \index{Insulin} ist ein Hormon und wird in den \gT{Beta-Zellen}
der Bauchspeicheldrüse (Pankreas) produziert. Es ist für
den Zuckerstoffwechsel wichtig. Es sorgt dafür, dass Zucker
(\gT{Glukose}) aus dem Blut in die Zellen aufgenommen wird.

\gFazit{Insulin senkt den Blutzuckerspiegel.} 


Wenn man künstlich Insulin spritzt, besteht die Gefahr eines
Unterzuckers (zu viel Insulin gespritzt oder zu wenig
gegessen). ${\rightarrow}$  Spritz/Essabstand
er\-fra\-gen.
}{

}{}{}{}




Zwei Störungen sind möglich:

\begin{enumerate}
    \item Es wird zu wenig Insulin gebildet.
    \item Der Körper hat einen gesteigerten Bedarf an Insulin weil er schon 
    zu sehr daran gewöhnt ist. Der Körper reagiert auf die 'normale' Dosis
    nicht mehr so wie er sollte, er braucht größere Mengen Insulin um
    den gleichen Effekt wie ein Gesunder zu erreichen (Gewöhnungseffekt),
    irgendwann kann aber der Körper nicht mehr Insulin produzieren.
\end{enumerate}

Folglich unterscheidet man 2 Typen:



\gLayout{0}{Diabetes mellitus: Typen}
{
\begin{labeling}[~--]{xxxxxxx }
    \item[Typ I]  \gZ{(IDDM,
    insuline dependent)} Die Insulinproduktion ist gestört und 
   der Körper produziert zu wenig Insulin. Der Patient
   muss als Ersatz künstlich Insulin spritzen.  Die Patienten sind eher jung
   (Jugendliche, junge Erwachsene) und sportlich.
   
   \gZ{Ursache für die Störung ist eine Störung des
   Immunsystems, bei der es zu einer Zerstörung der Insulin-produzierenden
   ${\beta}$-Zellen des Pankreas kommt. Durch den Produktionsausfall kommt es
   zu einem absoluten Mangel.}
    
    \item[Typ II]
    \gZ{(NIDDM, non insuline dependent)} Durch eine dauerhaft zu hohe
    Zuckerzufuhr kommt es im Körper zu Gewöhnungseffekten und einem
    Wirkungsverlust des vom Körper produzierten Insulins
    \gZ{(Insulinresistenz)}. Als Folge muss der Körper immer mehr Insulin produzieren um den gleichen Effekt zu
    erreichen. Wenn das Produktionslimit erreicht ist kommt es zu
    einer Erhöhung des Blutzuckerspiegels und zum Typ-2-Diabetes.
        
    Den Typ-II-Diabetes kann man mittels
    Lebensstiländerungen und Medikamenten behandeln. Im
    Endstadium muss mitunter aber auch Insulin gespritzt
    werden.
    
    Der Typ-2-Diabetes betrifft eher Patienten im mittleren Alter und aufwärts,
    aufgrund der zunehmend fett- und kohlehydratreichen Ernährungsgewohnheiten
    werden die Patienten aber immer jünger. 
    \begin{itemize}
        \item späte Manifestation: {\textgreater}\,35 Jahre
        \item Übergewicht
        \item 95\% aller Diabetiker
    \end{itemize}
    
\end{labeling}
}
{
\begin{itemize}
    \item Typ I
    \begin{itemize}
        \item Insulinproduktion gestört.
        \item Insulin muss gespritzt werden.
        \item Eher Jugendliche und junge Erwachsene.
    \end{itemize}
    \item Typ II
    \begin{itemize}
        \item Lebensstil, Ernährung.
        \item Insulin-Gewöhnungseffekt.
        \item Eher ab mittleren Alter.
        \item Lebensstiländerung, Medikamente, Insulin.
    \end{itemize}
\end{itemize}
}{}{}{}

\gLayout{0}{Spätfolgen}
{
Ein hoher Blutzuckerspiegel über Jahre schädigt den gesamten
Körper. Es kommt zu:
\begin{labeling}[~--]{xxxxxxxxxxxxxxx }
    \item[Gefäßschäden] durch Verkalkung
    \begin{itemize}
        \item Gefäßverschlüsse (kleine und große Gefäße)
        \item \gEs{Beine}: {\quotedblbase}Diabetischer
        Fuß{\textquotedblleft}, in Folge der Mangeldurchblutung in der Peripherie heilen auch kleine
        Wunden nicht mehr, in weiterer Folge muss immer weiter amputiert
        werden.
        \item \gEs{Herz}: Die Herzkranzgefäße
        (Koronar\-arterien) sind ebenso betroffen. Diabetiker sind Risikopatienten für
        Herz-Kreislauf-Erkrankungen wie zB Herz\-infarkt.
        \item \gEs{Hirn}: Wie beim Herzen kommt es zu
        Gefäßschäden, daher ist das Risiko für
        Schlaganfälle erhöht.
        \item \gEs{Niere}: Die sehr kleinen Gefäße der
        Niere werden geschädigt. Im weiteren Verlauf stellt die Niere ihre
        Arbeit ein, es kommt zu einer {\quotedblbase}\gT{chronischen
        Niereninsuffizienz}{\textquotedblleft}, die im Endstadium mittels
        Blutwäsche (Dialyse) behandelt werden muss. Dialyse-Patienten machen
        einen großen Teil der Krankentransporte aus.
    \end{itemize}
    \item [Nervenschäden] \gZ{(periphere Polyneuropathie)}:
    Gefühlslosigkeit vor allem in den Extremitäten, aber auch am
    gesamten Körper: Ein Herzinfarkt kann zum Beispiel schmerzlos
    ablaufen!
    \item [Augenschäden] bis hin zur Erblindung
    
\end{labeling}

Krankheiten, die sonst schmerzhaft sind, können beim Diabetiker aufgrund
der Nervenschäden ohne Schmerzen ablaufen. 

\gCave{Achtung: "`Stummer"', schmerzarmer/-loser Herzinfarkt beim Diabetiker!}
}{
\begin{itemize}
    \item Gefäßschäden (Herz, Hirn, Nieren, Extremitäten, \ldots)
    \item Nervenschäden 
    \item \gCa{Cave "`Stummer Infarkt"'!}
    \item Augenschäden
\end{itemize}
}{}{}{}








\subsubsection{Hypoglykämie}
\label{topic:hypoglykaemie}
\index{Hypoglykämie}

\gDefinition{Maßgebliche Erniedrigung des Blutzuckerspiegels.}

\gLayout{0}{Ursachen}
{
Die häufigsten Ursachen für eine plötzliche Hypoglykämie sind Diätfehler
und/oder eine falsche Anwendung von Insulin oder bestimmten anderen
Diabetes-Medikamenten. Aufgrund besonderer Anstrengungen kann es zu einem
gesteigerten Zuckerbedarf des Körper kommen, wodurch der Blutzuckerspiegel
sinkt. Dies kommt häufig im Rahmen von Begleiterkrankungen, wie z.B. grippalen
Infekten, vor. Auch die Alkohol-Intoxikation kann zu einer
Hypoglykämie führen\opt{NIVDOC}{ (Das beim Ethanolabbau
entstehende \gAbk{NADH} hemmt die
Glukoneogenese.)}\cite{LoefflerBiochemie7}. }{
\begin{itemize}
    \item zu viel Insulin (spritzen ohne zu essen).
    \item Fasten (spritzen ohne zu Essen).
    \item Besondere Anstrengungen, Begleiterkrankungen,
    insb. Infekte. 
    \item Alkohol-Intoxikation. 
\end{itemize}
}{}{}{}






\gLayout{0}{\gTxSymptome}
{
Der Patient mit Hypoglykämie ist meistens unruhig, desorientiert und agitiert,
oft auch unkooperativ und aggressiv. Oft fällt ein Zittern und eine schweißige
Haut auf. Auch Kopfschmerzen sind häufig. Eine Tachykardie und Tachypnoe ist
häufig feststellbar. Manchmal klagt der Patient über Übelkeit und Heißhunger.

Je tiefer der Blutzuckerspiegel sinkt, desto deutlicher werden die
\gEs{Bewusstseinsstörungen}. Von der anfänglichen
Agitiertheit/Desorientiertheit kann es rasch zu einer Somnolenz, zum Sopor und
zum Koma kommen!

Eine Hypoglykämie kann fast alle neurologischen  Störungen
aus\-lösen/imitieren! Häufig sind Kopfschmerzen und Bewusstseinsstörungen, evt.
auch Krampfanfälle. Bei jeder neurologischen Symptomatik muss man daher an eine Hypoglykämie denken!

\gCave{Bei der Hypoglykämie können Symptome auftreten, die
den Patient wie alkoholisiert wirken lassen! \gE{Hüte dich vor vorschnellen
Diagnosen!}}

\gFazit{${\rightarrow}$ BZ-Messung bei jeder neurologischen Störung. }

}{
\begin{itemize}
    \item Unruhe, Zittern, Schweißausbruch
    \item Aggressives Verhalten, unkooperativ
    \item Neurologische Symptomatik.
    \begin{itemize}
        \item Bewusstseinsstörungen, Verwirrtheit, Kopfschmerzen, Krampfanfälle,
        Bewusstlosigkeit.
    \end{itemize}
    \item Heißhunger, Übelkeit.
    \item Tachykardie, Tachypnoe.
    \item Haut feucht/blass.
\end{itemize}
}{}{}{}












\gMassnahmenNG{Spezielle Maßnahmen}{Hypoglykämie}{}{
\begin{itemize}
    \item Vitale Bedrohung einschätzen. \gTxBeiVit \gTxMassVit
    
    \item BZ messen! Zur Diagnosestellung bzw. \gE{vor} und \gE{nach}
    jeder Zuckergabe
    \item Zucker zuführen: Wenn der Patient
    bewusstseinsklar ist Traubenzucker oder Zuckerwasser geben. Danach "`feste
    Speisen"' nachessen lassen, da der Traubenzucker nur eine kurze Wirkdauer hat.
    
    Wenn der Patient nicht mehr bewusstseinsklar ist kann die Zuckergabe
    nur über eine venöse Gabe erfolgen (getrübte Patienten dürfen aufgrund
    der Aspirationsgefahr nichts mehr essen).
    
\end{itemize}
}

\subsubsection{Hyperglykämie}
\label{topic:hyperglykaemie}
\index{Hyperglykämie}

% Done: XXX INHALT-NEU: Hyperglykämie, Text und Überarbeitung,
% POI 11
\gA{\#11: Muss noch überarbeitet werden}

\gDefinition{Maßgebliche Erhöhung des Blutzuckerspiegels.}

\gLayout{0}{Hyperglykämie: Zustand oder Notfall?}
{
Eine Erhöhung des Blutzuckerspiegels ist nicht automatisch ein Notfall. Die
Bedrohlichkeit ist z. B. stark davon abhängig wie
rasch sich diese Erhöhung entwickelt hat. So hat ein
Patient, bei dem seit Jahren ein unerkannter Diabetes mellitus besteht auch mit sehr hohen Blutzuckerwerten (akut) keine
Probleme. Der Körper hat sich über die Zeit daran \gE{gewöhnt}. 

Ein anderer
Patient, bei dem der Zustand plötzlich eingetreten ist (z.B.
Erstmanifestation eines Diabetes mellitus Typ I bei einem 19-jährigen, oder wenn
das Insulin nicht gespritzt wurde, \ldots), kann mit dem gleichen
Blutzuckerwert bereits im Koma liegen. 

\gFazit{Beurteile nicht nur den Messwert, sondern immer auch den Zustand des
Patienten.} }
{←}
{}{}{}



\gLayout{0}{Auswirkungen}
{
Im Grunde bereitet der Zucker im Akutfall dem Körper zwei
Probleme. Einerseits wirkt der Zucker wie ein Elektrolyt auf den \gEs{Flüssigkeitshaushalt} ein und
stört diesen. Es kommt zu einer ungünstigen Wasserumverteilung. Andererseits
mangelt es den Körperzellen an Zucker wenn das Insulin fehlt um es in die
Zellen hineinzubekommen, es entsteht ein \gEs{Nährstoffmangel}.}{
\begin{itemize}
    \item Störung des Wasser- und Elektrolythaushaltes.
    \begin{itemize}
        \item Zucker wirkt wie ein Elektrolyt.
    \end{itemize}
    \item Störung der Nährstoffversorgung.
\end{itemize}
}{}{}{}

\gLayout{0}{Ursachen}
{
Die Hyperglykämie betrifft praktisch nur Diabetes mellitus-Patienten. Bei
Patienten, die einen noch unerkannten DM~Typ~I haben, ist ein hyperglykämischer
Notfall oft die Erstmanifestation (die Diagnose Diabetes ist noch nicht
bekannt!). Sonst sind Diätfehler oder eine mangelnde Insulinzufuhr häufige
Ursachen. Auch Stress und Begleiterkrankungen wie z.B. akute Infektionen
kommen in Betracht, da hier der Insulinbedarf steigt.
}{
\begin{itemize}
    \item Erstmanifestation eines DM~Typ~I
    \item Diätfehler
    \item Insulinzufuhr unzureichend/vergessen
    \item Stress
    \item Infektionen (Insulinbedarf ${\uparrow}$)
\end{itemize}
}{}{}{}









\gLayout{0}{\gTxSymptome}
{
Der Patient hat i.d.R. ein erhöhtes Durstgefühl und gleichzeitig eine
\gEs{vermehrte Harnproduktion} (Zucker wird ausgeschieden). Er klagt über
Abgeschlagenheit, in weiterer Folge kann es zu \gEs{Bewusstseinsstörungen} bis
hin zum \gEs{Koma} (s.u.) kommen. Die Geschwindigkeit der Veschlechterung ist je
nach Patient sehr unterschiedlich.

\subparagraph{Koma} \label{topic:koma-diabetisches} \gZ{Wie oben bereits
beschrieben wurde, kann es sowohl zu einem Problem des
Flüssigkeitshaushaltes und/oder der Nährstoffversorung der
Zellen kommen. Je nachdem welches Problem im Vordergrund steht, kommt es zu} \gE{unterschiedlichen Symptomen}. 

Ist vorwiegend der Flüssigkeitshaushalt gestört kommt es zum sog.
\gZ{Hyperosmolaren Koma}, dessen Leitsymptom erwartungsgemäß die
Bewusstseinsstörung ist.

Steht der Nährstoffmangel im Vordergrund passiert etwas Interessantes: Der
Körper produziert sog. \gZ{Ketonkörer}. Da diese \gE{sauer} sind kommt es zu
einer stoffwechselbedingten Übersäuerung (metabolische \gEs{Azidose}
\gSee{topic:azidose-metabolische}).

Unter \gSee{topic:alkalose-respiratorische} wurde besprochen dass der pH-Wert
bei einer Hyperventilation in Richtung basisch wandert. Bei der diabetischen
Ketoazidose (sauer) versucht der Körper diesen Mechanismus auszunutzen: Er \gEs{hyperventiliert} um sich "`vom sauren pH
    zum normalen pH zu atmen"'. Es ist daher bei diesen Patienten eine
    \gEs{schnelle, tiefe Atmung} feststellbar\gZ{\footnote{\gZ{Diese
    charakteristische Form der Atmung wird auch "`Kußmaul-Atmung"' genannt}}}.

Mitunter kann auch ein blumig/fruchtiger oder Azeton-artiger Atemgeruch
wahrgenommen werden. Dieser darf nicht mit einer "`Alkoholfahne"' verwechselt
werden! 
}{
\begin{itemize}
    \item Durstgefühl
    \item Vermehrte Harnproduktion
    \item Abgeschlagenheit, Bewusstseinstrübung bis
    \item Bewusstlosigkeit (Koma)
    \begin{itemize}
        \item Problem Flüssigkeitshaushalt
        \begin{itemize}
            \item \gZ{Hyperosmolares Koma}
        \end{itemize}
        \item Problem Nährstoffmangel
        \begin{itemize}
            \item \gZ{Ketoazidotisches Koma}
            \item Azidose → Hyperventilation zur Kompensation
        \end{itemize}
    \end{itemize}   
    \item Azeton-Geruch (Nicht mit Alkoholfahne verwechseln!)
\end{itemize}
}{}{}{}



\gMassnahmenNG{Spezielle Maßnahmen}{Hyperglykämisches Koma}{}{
\begin{itemize}
    \item \gTxMassVitBes
    \begin{itemize}
        \item Lagerung: Stabile Seitenlage
        \item \ce{O2}: ab 10\gLMin
    \end{itemize}
\end{itemize}
}





\clearpage
%%%%%%%%%%%%%%%%%%%%%%%%%%%%%%%%%%%%%%%%%%%%%%%%%%%%%%%%%%%
%%% Infektionskrankheiten                               %%%
%%%%%%%%%%%%%%%%%%%%%%%%%%%%%%%%%%%%%%%%%%%%%%%%%%%%%%%%%%%

\section{Infektions- und Entzündungskrankheiten}


%%%%%%%%%%%%%%%%%%%%%%%%%%%%%%%%%%%%%%%%%%%%%%%%%%%%%%%%%%%
%%%%%%%%%%%%%%%%%%%%%%%%%%%%%%%%%%%%%%%%%%%%%%%%%%%%%%%%%%%
%%%%%%%%%%%%%%%%%%%%%%%%%%%%%%%%%%%%%%%%%%%%%%%%%%%%%%%%%%%
\subsection{Grippe -- Influenza}
\index{Grippe}\index{Influenza}\index{Influenza}
\label{topic:grippe}

\gDefinition{Eine schwere, durch Viren verursachte
Infektionskrankheit.}

\gLayout{2}{\gTxBeschreibung}
{
Die Grippe ist eine \gE{virale} Erkrankung, welche je nach
Patient und Virenstamm (s.u.) einen sehr schweren
Verlauf nehmen und auch tödlich enden kann. Gefürchtet sind
\gEs{Komplikationen} wie Lungen- oder
Herzmuskelentzündungen oder weitere (Folge-)Infektionen.
Besonders gefährdet sind ältere und immunsupprimierte Personen.
}{}{}{}{}


\gFazit{Da es sich bei der Grippe um eine durch Viren
verursachte Krankheit handelt sind \gEs{Antibiotika
wirkungslos}.}

Die Grippe tritt besonders in der kalten Jahreszeit auf und tritt
regelmäßig als \index{Epidemie}\gEs{Epidemie} auf. Es
kann auch zum Ausbruch einer weltweiten \index{Pandemie}Pandemie kommen, \gZ{z.B.
Spanische Grippe (1918, 25~Mio~Tote), Asiatische Grippe
(1957, 1~Mio~Tote) und  Hongkong-Grippe (1968,
800.000~Tote)\footnote{http://de.wikipedia.org/wiki/Influenza}}.
\marginlabel{\cite{wikiSpanischeGrippe}}

In Österreich erkranken pro Jahr ca. 380.000 Personen an der
Influenza, ca. 4.500 Personen müssen stationär behandelt werden und
ca. 120 Personen sterben an den
Folgen\cite{Strauss0000}.

%\footnote{http://www.bmg.gv.at/cms/site/standard.html?channel=CH0742\&doc=CMS1075899626901}.

\gFazit{Die Grippe tritt in Wellen auf und kann tödlich
enden.}

\gLayout{2}{\gTxSymptome}
{
Typisch für die Grippe ist das \gEs{rasche, hohe
Anfiebern (38{\textdegree}-40{\textdegree}C) mit
Schüttelfrost,  Kopf-, Hals-, Muskel-, Gelenksschmerzen, Husten und stark reduziertem
Allgemeinzustand} sowie lang dauernde Abgeschlagenheit nach
Ende der Erkrankung. Bei schon vorher geschwächten
Patienten kann der Verlauf weniger spektakulär aussehen (z.B. kaum Fieber), das Risiko einer Komplikationen ist allerdings natürlich höher. 
}
{
% TODO VIT: Slide fehlt!
}
{}
{}
{}


\gLayout{2}{Stämme}
{
Grippeviren kommen in verschiedenen
Stämmen (\gCode{A}, \gCode{B} und \gCode{C}) vor, welche
wiederum durch unterschiedliche Arten von
\gZ{Hämagglutinin} (\gCode{H}) und \gZ{Neuramidase}
(\gCode{N}) unterschieden werden. Dadurch kommen so
klingende Namen wie "`\gCode{A(H5N1)}"' oder
"`\gCode{A(H1N1)}"' zu Stande. Je nach Unterart verhalten sich die Viren anders und man beobachtet Unterschiede hinsichtlich der
Ansteckungsgefahr, Schwere der Erkrankung usw. 
\gZ{}
}
{

}
{}
{}
{}

\gLayout{2}{Schutzimpfung}
{
Gegen die Grippe wird eine
\gEs{Schutzimpfung} angeboten. Sie basiert auf der Voraussage
welcher Stamm in der jeweils bevorstehenden Grippewelle
vorherrschen wird und wirkt daher nur für ein Jahr und muss
jede Saison wiederholt werden. Die Impfung bietet keinen
hundertprozentigen Schutz, wird aber trotzdem für
Risikopatienten und Personen die einem erhöhten
Infektionsrisiko ausgesetzt sind (darunter fällt auch
Personal im Gesundheitsbereich) empfohlen.
}
{
% TODO VIT: Slide fehlt!
}
{}
{}
{}

\cite{Gallaher2009,H1N1InvTeam2009,Peiris2009,Wolf2009,Cutler2009}



\gLayout{2}{Schweinegrippe, Vogelgrippe -- Warum die große
Aufregung?}
{
Grippeviren befallen nicht nur den Menschen, auch
(Wasser-)Vögel und andere Tiere erkranken an der Grippe.
Normalerweise erfolgt keine Übertragung vom Tier zum
Menschen. Überschreitet ein Erreger diese Grenze wird es
gefährlich: Die Tierwelt beinhaltet dann einen nahezu
unbegrenzten Vorrat an Keimen (Keimreservoir). Dieser kann 
praktisch nicht kontrolliert werden, schließlich lassen
sich Zugvögel nicht so leicht unter Quarantäne stellen
\ldots 

Glücklicherweise entwickelten sich in letzter Zeit nur
Viren, welche zwar vom Tier auf den Menschen und dann
weiter auf den Menschen übertragbar waren, jedoch waren sie
(bis jetzt) nicht übermäßig infektiös oder zeigten keinen
besonders schweren Verlauf. 

\subparagraph{Die Gefahr} Die
Tier\,→\,Mensch\,→\,Mensch-Übertragung stellt eine ständige
Gefahr dar. 

Ein weiteres Problem ist jedoch die mögliche Kreuzung
mit Stämmen die sehr infektiös und schwer verlaufend sind. Dies kann passieren wenn ein
Patient mit mehreren Stämmen gleichzeitig infiziert ist.
Daher ist es wichtig die Infektionen so rasch als möglich
einzudämmen. 
}
{
% TODO VIT: Slide fehlt!
}
{}
{}
{}



\gFazit{Die große Gefahr ist die Entwicklung eines hoch
infektiösen und schwer verlaufenden Stammes, welcher eine
Tier-Mensch-Mensch-Übertragung aufweist.}
\gFazit{Eine rasche Isolierung der Infzierten soll die
Wahrscheinlichkeit einer bösartigen Kreuzung minimieren!}

\subparagraph{Neue Grippe ("`Schweinegrippe"')} Die jüngste
Entwicklung war die Übertragung des Stammes \gCode{A(H1N1)}
vom Schwein auf den Menschen in Mexiko im Frühjahr 2009. Der
neue Stamm breitete sich sehr rasch um die ganze Welt aus,
\gE{die Erkrankung verlief jedoch recht mild}. Bemerkenswert
war jedoch, dass -- anders als bei einer "`normalen
Grippe"' üblich -- besonders junge Menschen infiziert wurden. 

Man darf gespannt sein was die Zukunft bringt \ldots

% \begin{tabulary}{\linewidth}{lLL}\toprule
% 
% \gTabHeading{Übertragung}
% &
% \gTabHeading{~}
% &
% \gTabHeading{Einschätzung}
% \\\midrule
% Mensch → Mensch
% &
% &
% \\
% Tier → Mensch
% &
% &
% \\
% Tier → Mensch → Mensch
% &
% &
% \\
% Tier → Mensch → Mensch
% &
% hohe Ansteckungsrate
% 
% schwerer Verlauf
% &
% sehr schlimm
% \\\bottomrule
% \end{tabulary}

%%%%%%%%%%%%%%%%%%%%%%%%%%%%%%%%%%%%%%%%%%%%%%%%%%%%%%%%%%%
%%%%%%%%%%%%%%%%%%%%%%%%%%%%%%%%%%%%%%%%%%%%%%%%%%%%%%%%%%%
%%%%%%%%%%%%%%%%%%%%%%%%%%%%%%%%%%%%%%%%%%%%%%%%%%%%%%%%%%%
\subsection{Grippaler Infekt}\index{grippaler
Infekt}

\emph{Syn.: \index{Verkühlung}Verkühlung,
 \index{Erkältung}Erkältung} 

\gLayout{2}{\gTxBeschreibung}
{
Die Verkühlung ist eine Infektion der oberen Luftwege mit Viren,
allerdings mit anderen als bei der {\quotedblbase}echten
Grippe{\textquotedblleft} (Influenza). Die Symptome ähneln denen der
Grippe, allerdings sind sie bei weitem nicht so stark ausgeprägt. Im
Vordergrund stehen Halsschmerzen, Husten, rinnende oder verstopfte Nase
und eine eher leicht erhöhte Temperatur. 
}{}{}{}{}



%%%%%%%%%%%%%%%%%%%%%%%%%%%%%%%%%%%%%%%%%%%%%%%%%%%%%%%%%%%
%%%%%%%%%%%%%%%%%%%%%%%%%%%%%%%%%%%%%%%%%%%%%%%%%%%%%%%%%%%
%%%%%%%%%%%%%%%%%%%%%%%%%%%%%%%%%%%%%%%%%%%%%%%%%%%%%%%%%%%
\subsection{Lungenentzündung -- \gT{Pneumonie}}
\label{topic:pneumonie}

Streng genommen ist die Lungenentzündung eine Erkrankung der
Atemwege. Meistens steht bei den Symptomen eher die
entzündliche Komponente im Vordergrund (Fieber, allgemeine
Schwäche, \ldots).

\gLayout{2}{Ursachen}
{
\begin{itemize}
    \item Infektiös: Bakterien, Pilze
    \gZ{(Pneumokokken,Chlamyden, Candida bei
    HIV/Immunsupression)}
    \item Chemische Schädigung: Spätfolge nach
    \index{Pneumonie!Aspiration}Aspiration
\end{itemize}

Höheres Lebensalter begünstigt Infektionen wie z.B.
Pneumonien, da ältere Patient ein zunehmend schwächeres
Abwehrsystem haben.
}{}{}{}{}


\gLayout{2}{\gTxSymptome}
{
Die Symptome sind oft nicht charakteristisch, eine Diagnose kann meist
erst im Krankenhaus gestellt werden. 

\begin{itemize}
    \item Dyspnoe, Husten
    \item Fieber (muss aber nicht sein)
    \item Reduzierter Allgemeinzustand (AZ)
\end{itemize}
}{

}{}{}{}





\gMassnahmenNG{Spezielle Maßnahmen}{Pneumonie}{}{
    
    \begin{itemize}
        \item Symptomatisch
        \item Wärmeerhalt
        \item Hospitalisierung, Abteilung: Innere Medizin
    \end{itemize}    
}

\subsection{Bakterielle Meningitis}
\index{Meningitis}
\label{topic:meningitis}

\gDefinition{Eitrige Entzündung der Hirnhäute aufgrund
einer Infektion mit Bakterien.}

Die bakterielle Meningitis ist eine gefürchtete Erkrankung. Einerseits ist sie
hoch ansteckend, andererseits kann sie rasch tödlich enden. Gefährdet sind
grundsätzlich alle Altersgruppen, besonders aber Kinder, Jugendliche und alte
Menschen.
 
 
\gLayout{0}{\gTxSymptome}
{
Leitsymptome sind \gEs{Fieber}, \gEs{Nackensteifigkeit} (Meningismus) und
\gEs{Lichtscheu}. Darüber hinaus kann es zu anderen unspezifischen
Symptomen wie Kopfschmerzen, Übelkeit, Erbrechen, etc. kommen. 

\gE{Die Meningitis wird oft mit einem einfachen grippalen Infekt oder einer
Grippe verwechselt!}

\gZ{Im fortgeschrittenen Stadium kann es zu massiven
Gerinnungsstörungen kommen, am Körper bilden sich dann dunkle
Hämatome (s. Abb. \ref{fig:meningitis-baby}). Einer der
häufigsten Erreger sind die Meningokokken, man spricht dann
von einer Meningokokken-Meningitis.}

}
{
\begin{itemize}
    \item Fieber, Kopfschmerzen.
    \item Nackensteifigkeit.
    \item Lichtscheu.
    \item Erbrechen, Übelkeit.
    \item Hämatome am ganzen Körper (im Endstadium).
    \item Verwechslungsgefahr mit "`normalem Infekt"'!
\end{itemize}

}{}{}{}

\gLayout{22}{}
{
Endstadium einer bakteriellen Meningitis.
Die Flecken auf der Haut sind Blutungen in Folge einer nicht mehr
funktionierenden Blutgerinnung. Der Tod ist bei dieser
Patientin innerhalb weniger Stunden zu erwarten.
}{}
{./images/meningitis-kind-waterhouse.png}
{\label{fig:meningitis-baby}}
{unbekannt}




\gAuthNotes{FEHLT: Querverweis Meningimus\\FEHLT: Maßnahmen}


\gMassnahmenNG{Spezielle Maßnahmen}{Meningitis}{}{
\begin{itemize}
    \item Schutzmaske für Personal und Patient (wenn zumutbar).
    \item Rückmeldung an Leitstelle vor Transport, Voranmeldung im Krankenhaus.
    \item Patient erst in das Spitalsgebäude bringen, wenn
    Ziel und Weg dorthin klar ist.
    \item Desinfektionsmaßnhamen gem. Hygieneplan. 
    \item Personal bekommt im Spital evt. eine Antibiotika-Prophylaxe. 
\end{itemize}

}

\subsection{Wundstarrkrampf -- Tetanus}
\index{Tetanus}\index{Wundstarrkrampf}


%\includegraphics[width=5.984cm,height=3.539cm]{images/AASdev20090228-img97.gif}
%\gCaptionof{figure}{ \gAuthNotes{Tetanus Bild in gif:img97} }



Der Erreger des Wundstarrkrampfes kommt \gE{überall in der
Umgebung} vor. Er produziert einen Giftstoff (Toxin),
welches einen tonischen Krampf der Skelettmuskulatur
auslöst. \gZ{Die Inkubationszeit beträgt ca. 4-21 Tage.
Erste Vorzeichen sind Müdigkeit und Inappetenz.}

\gMarriedLiii{\gTxSymptome}
{
Es kommt zu tonischen \gEs{Krämpfen der
quergestreiften Muskulatur}, oft beginnend bei der
Kaumuskulatur. Der Patient hat starke Schmerzen und
entwickelt in Folge der Krämpfe einen \gEs{Atemlähmung}. Die
\gE{Sterblichkeit ist hoch}, die Letalität beträgt ca.
50\,\%.
}
{
% TODO VIT: Slide fehlt!
}
{./images/hygiene/6373_lores_tetanus.jpg}
{\gCaptionof{figure}{\gAuthNotesPicLegalOk{PD}{}Patient
mit Tetanus. \gSourcePicture{Public Health Image Library,
ID\#: 6373, Public Domain.}}}\emph{\gZ{Clostridium tetani}}

\gLayout{0}{Schutzimpfung}
{
\index{Impfschutz} 
Es existiert
eine Schutzimpfung, diese muss dem Patienten \gE{rechtzeitig} nach der
Verletzung gegeben werden. Der Impfschutz beträgt
normalerweise ca \gE{10 Jahre}, bei einer
Verletzung wird sicherheitshalber allerdings schon bei einer
\gEs{länger als 5 Jahre} zurückliegenden Impfung eine neuerliche Gabe verabreicht.

Auch bei kleinen Verletzungen muss nach dem Impfschutz
gefragt werden, dies muss auch unbedingt \gEs{dokumentiert
werden!} Ist das Datum der letzten Impfung nicht
sicher erinnerlich soll der Patient unverzüglich zu Hause im
Impfpass nachschauen. Erforderlichenfalls soll der Patient
unverzüglich eine Unfallabteilung aufsuchen. 

Die entsprechende Aufklärung muss am Revers bestätigt werden!
}
{
% TODO VIT: Slide fehlt!
}
{}
{}
{}

\gCave{Eine mangelhafte Aufklärung oder Dokumentation
bezüglich Tetanus und dessen Impfschutz ist grob fahrlässig und ein unverzeihbarer
Kunstfehler.}


%%%%%%%%%%%%%%%%%%%%%%%%%%%%%%%%%%%%%%%%%%%%%%%%%%%%%%%%%%%
%%%%%%%%%%%%%%%%%%%%%%%%%%%%%%%%%%%%%%%%%%%%%%%%%%%%%%%%%%%
%%%%%%%%%%%%%%%%%%%%%%%%%%%%%%%%%%%%%%%%%%%%%%%%%%%%%%%%%%%


\subsection{Tuberkulose}
\label{topic:tuberkulose}

\emph{\gZ{Mykobacterium Tuberculosis, "`Schwindsucht"'.}}
Abk.: "`\gT{Tbc}"'

\gLayout{28}{Beschreibung}
{
Die Tuberkulose kann praktisch alle Organe befallen, der
Verlauf wird vor allem davon bestimmt, welches
Organ befallen ist \gZ{(Lunge, Skelett, Darm, Haut,
Genitalien, ...)}. Ausgangspunkt ist fast immer die \gEs{Lunge}.

Die Übertragung erfolgt vor allem \gEs{aerogen}, aber unter
Umständen auch oral, durch Hautkontakt oder plazentär. 

Zunehmend treten auf Antibiotika resistente Erreger auf \gZ{(MDR-TB,
XDR-TB)\footnote{Quelle: AGES}}
% TODO LIT, HYG: MDR/XDR-TB Quelle bei AGES
.
    
\gZ{Die Inkubationszeit beträgt 4-12 Wochen.}
}{}{./images/AASdev20090228-img98.jpg}{\gAuthNotesPicLegalNotOk{img98}{Durch Foto von Hauer
ersetzen}Röntgenbild eines TB-Pat\-ient\-en.}{}


\gLayout{0}{\gTxSymptome}
{
Es dominieren vor allem Allgemeinsymtptome wie erhöhte
Temperatur, Nachtschweiß, Lymphknotenschwellung und Husten.
Auffallend ist manchmal die Schwere des Hustens, Blutbeimengungen gelten
als Alarmsignal. 

In der Anamnese ist mitunter eine bereits druchgemachte
Tuberkulose erkennbar.
}
{
\begin{itemize}
    \item Husten, Bluthusten.
    \item erhöhte Temperatur.
    \item \gEs{Nachtschweiß, Lymphknotenschwellung}
    \item Anamnese!
\end{itemize}
}
{}
{}
{}


\gLayout{2}{Offene Tuberkulose}
{
\index{Offene Tb} 
Wenn der Infektionsherd in der Lunge \gE{Luftkontakt} hat, spricht man von der \gT{Offenen
Tuberkulose}. Der Patient ist dann infektiös,
eine Tröpfcheninfektion ist möglich. Die Verwendung von
Schutzmasken ist anzuraten. 

Es ist äußerlich \gE{nicht sicher erkennbar} ob der Patient
an einer offenen oder geschlossenen Tuberkulose leidet.
Bluthusten wäre allerdings ein deutlich Hinweiszeichen.
}
{
% TODO VIT: Slide fehlt!
}
{}
{}
{}


\gZ{
\paragraph{Ansteckungsgefahr}
Bei Einhaltung der Schutzmaßnahmen (s. Spezielle Maßnahmen) ist eine Ansteckung
von sonst gesunden und immunkompetenten Personal unwahrscheinlich. Für eine
Ansteckung  wäre ein ca. 8-stündiger Kontakt erforderlich
\footnote{National Institute for Clinical Excellence (NICE) Tuberculosis.
National clinical Guideline for Diagnosis, management, prevention and control
(2006)}
% TODO HYG, LIT: National Institute for Clinical Excellence (NICE)
% Tuberculosis. National clinical Guideline for Diagnosis, management,
% prevention and control (2006)
.
}

\gMassnahmen{
\paragraph{Spezielle Massnahmen: Tuberkulose}~

\begin{itemize}
    \item Bei Verdacht auf eine offene Tuberkulose
    Schutzmasken verwenden: FFP~3 für das Personal, "`OP-Maske"' für den
    Patienten, wenn zumutbar.
    \item \gEs{Meldung an Leitstelle \& Betriebsleitung!}
    \item Desinfektion laut Hygieneplan.
\end{itemize}
}


%%%%%%%%%%%%%%%%%%%%%%%%%%%%%%%%%%%%%%%%%%%%%%%%%%%%%%%%%%%
%%%%%%%%%%%%%%%%%%%%%%%%%%%%%%%%%%%%%%%%%%%%%%%%%%%%%%%%%%%
%%%%%%%%%%%%%%%%%%%%%%%%%%%%%%%%%%%%%%%%%%%%%%%%%%%%%%%%%%%
\subsection{HIV und AIDS}
\index{HIV}\index{Humanes Immundefizienz
Virus}\label{topic:hiv}\label{topic:aids}

\gT{HIV}: \gE{Humanes Immundefizienz Virus}. 
\gT{AIDS}: \gE{Aquired Immune Deficiency Syndrom},
\gE{erworbenes Immunschwäche-Syndrom}

\gLayout{22}{HIV und AIDS}
{
\subparagraph{HI-Virus} Das HIV befällt vor allem die Zellen des
Abwehrsystems. Es vermehrt sich in ihnen, setzt sie außer Funktion und zerstört sie schließlich. Das körpereigene Abwehrsystem kann HIV nicht aus dem Körper
    entfernen.
Es kommt zu einer Immunschwäche.

\subparagraph{AIDS} HIV verursacht das \gE{AIDS} (Aquired Immune Deficiency
Syndrom, erworbenes Immunschwäche-Syndrom), welches die
eigentliche Erkrankung darstellt. Dabei kommt es durch das
\gE{geschwächte Immunsystem} zu Folgeerkrankungen, an denen der Patient versterben kann:

\begin{itemize}
    \item \gE{Infektionen} mit normalerweise harmlosen
    Erregern
\begin{itemize}
    \item Pilze, {\textquotedblleft}fakultative{\textquotedblright} Erreger,
    {\textquotedblleft}Opportunisten{\textquotedblright}, seltene Arten der
    Lungenentzündung
\end{itemize}
    \item \gE{Tumore}
\end{itemize}

Abhängig von der durchgeführten Therapie bricht das AIDS
einige Jahre nach der eigentlichen HIV-Infektion aus.
}{
% TODO VIT: Slide fehlt!
% \begin{itemize}
%     \item \gZ{Point}
% \end{itemize}
}
{./images/hygiene/Niaid-hiv-virion-mod.jpg}
{\gAuthNotesPicLegalOk{PD}{}Schema
eines HI-Virus.}
{US National Institute of Health, PD-USGov-HHS-NIH.}









\gLayout{2}{Übertragung}
{
Übertragen wird das HI-Virus v.a. durch Körperflüssigkeiten
wie Blut, Blutprodukte, Sperma und Scheidensekret.
\gE{"`Ungeschützter Sex"'}, Transfusionen und
\gE{Nadelstichverletzungen} gelten als klassische Risiken.
Es erfolgt \gEs{keine Übertragung durch normale
Sozialkontakte} wie Händeschütteln, wohnen im
gemeinsamen Haushalt, gemeinsame Benutzung von Geschirr und
Toilettenanlagen, etc.
}
{
% TODO VIT: Slide fehlt!
}
{}
{}
{}

\gLayout{2}{Test}
{
Erst nach 6-12 Wochen sind Antikörper nachweisbar. Das
bedeutet dass der \gEs{HIV-Test nicht
    {\textquotedblleft}aktuell{\textquotedblright}, sondern
    {\textquotedblleft}12 Wochen blind{\textquotedblright}}
    ist!
}
{
% TODO VIT: Slide fehlt!
}
{}
{}
{}


\gLayout{2}{Therapie}
{
\begin{itemize}
    \item Derzeit ist \gEs{keine Heilung} möglich, jedoch
    kann eine \gE{deutliche Lebensverlängerung} durch
    Medikamente erreicht werden.
\end{itemize}
}
{
% TODO VIT: Slide fehlt!
}
{}
{}
{}

\gLayout{2}{Prophylaxe}
{
% TODO VIT: Fließtext fehlt!
\begin{itemize}
    \item Allgemeine Selbstschutzmaßnahmen bei der
    Patientenversorgung: Schutzhandschuhe, Kontakt mit
    Körperflüssigkeiten vermeiden.
    \item Postexpositionsprohylaxe bei Nadelstichverletzungen
    \item Kondome
\end{itemize}
}
{
% TODO VIT: Slide fehlt!
}
{}
{}
{}




%%%%%%%%%%%%%%%%%%%%%%%%%%%%%%%%%%%%%%%%%%%%%%%%%%%%%%%%%%%
%%%%%%%%%%%%%%%%%%%%%%%%%%%%%%%%%%%%%%%%%%%%%%%%%%%%%%%%%%%
%%%%%%%%%%%%%%%%%%%%%%%%%%%%%%%%%%%%%%%%%%%%%%%%%%%%%%%%%%%
\subsection{Hepatitis}
\label{topic:hepatitis-virale}\label{topic:hepatitis-toxisch}\label{topic:hepatitis}
\index{Hepatitis!infektiöse}\index{Hepatitis!toxische}

\gDefinition{Leberentzündung durch verschiedene (virale)
Erreger (A -- E) oder durch sonstige
(chronische) Schädigung.}



\gLayout{0}{\gTxBeschreibung{} und Ursachen}
{%%% Gerhard Gabriel <gerhard.gabriel@grissu.org> 
Hepatitis ist eine
entzündliche Erkrankung der Leber. Diese Entzündung kann durch Erreger wie
Viren, oder Mikroorganismen hervorgerufen werden
(infektiöse, virale Hepatitis). Aber auch
regelmäßige Schädigungen der Leber durch Giftstoffe wie
Alkohol, Medikamente \gZ{(z.B. Paracetamol)}, etc. können
zu einer (chronischen) Entzündung führen. Auch ein länger
bestehender Rückstau der Gallenflüssigkeit, z.B. bei einer
Verstopfung des Gallenganges, kann zu einer Entzündung
führen.

Je nach Ursache der Hepatitis kann es vorkommen, dass es zu
keiner Abheilung kommt und die Hepatitis dauerhaft bestehen
bleibt. Man spricht dann von der \gT{chronischen
Hepatitis}.} 
{
\begin{itemize}
    \item Virusinfektionen (Hep A, B, C, \ldots ).
    \item Gallen-Abflussstau.
    \item Vergiftungen Lebensmittel(Pilze)~/
    Medikamente (Paracetamol)~/ chron. Alkoholismus.
    \item Chronische Verlaufsform möglich.
\end{itemize}
}{}{}{}

\gLayout{0}{\gTxSymptome}
{
Die Symptomatik einer Hepatitis ist
uncharakteristisch und eher allgemein: Unspezifisches mehr oder minder
ausgeprägtes Krankheitsgefühl, ähnlich einem grippalem Infekt, Appetitlosigkeit.
Auch erhöhte Körpertemperatur und Schmerzen im
rechten Oberbauch kommen bisweilen vor. Die Leber ist oft
vergrößert.

Im
fortgeschrittenen Stadium kann es zu einer Gelbfärbung
des Augenweiß und eine gelbliche Verfärbung der Haut
(Ikterus) kommen. 

Bei einem chronischen
Verlauf kommt es zu einem Umbau von
funktionierendem Lebergewebe zu
funktionslosem Bindegewebe. Man spricht
dann von der
\gT{Leberzirrhose}\label{topic:leberzirrhose}\index{Leberzirrhose}.
Schließlich kann es zu einem Leberversagen und in kurzer
Folge zum Tod kommen, da die Leber ein überlebenswichtiges
Organ ist. 
}{
\begin{itemize}
    \item uncharakteristische Allgemeinsymptome (grippale Symptome,
    Appetitlosigkeit).
    \item Schmerzen im re. Oberbauch.
    \item Gelbfärbung Augen/Haut (Ikterus ).
\item Vergrößerte, schmerzhafte Leber.
    \item Ikterus (Gelbsucht).
    \item Unterscheide {\textquotedblleft}akut{\textquotedblright} und
    {\textquotedblleft}chronisch{\textquotedblright}.
    \begin{itemize}
        \item Chronisch: wenn Hepatitis nicht abheilt.
    \end{itemize}
    \item Ende: Leberversagen, Tod.
\end{itemize}
}{}{}{}



\paragraph{Verschiedene Arten der viralen Hepatitis}~

Abhängig vom jeweiligen Virus unterscheidet man
verschiedene Arten der viralen Hepatitis:

\begin{table}[H]
\begin{gWideParEnv}
\gCaptionof{table}{Übersicht der wichtigsten viralen Hepatitis-Arten.}

\begin{tabular}{p{3cm}p{\linewidth-4cm}}
\gTabHeading{Typ}
&
\gTabHeading{Bemerkungen}
\\\midrule\addlinespace
\gTabSub{Hepatitis-A} 

Hepatitis-A-Virus (HAV)
&
\begin{itemize}
    \item \gZ{Inkubationszeit 10 - 40 Tage}
    \item typische Reisekrankheit (Mittelmeerraum, Südamerika, Orient),
    meist epidemisch,
    \item Dauer einige Wochen,
    \item kein Übergang in chronischen Verlauf
\end{itemize}

\\\addlinespace
\gTabSub{Hepatitis-B} 

Hepatitis-B-Virus (HBV)
&
\begin{itemize}
    \item \gZ{Inkubationszeit 4 - 12 Wochen,}
    \item \gZ{meist sporadisches Auftreten}
    \item Übertragung durch Speichel, Blut(produkte),
    Sperma und andere Körperflüssigkeiten,
    \item chronischer Verlauf in ca. 5 \%
\end{itemize}
\\\addlinespace
\gTabSub{Hepatitis C}

Hepatitis-C-Virus (HCV)
&
\begin{itemize}
    \item \gZ{Inkubationszeit 6 - 12 Wochen,}
    \item häufig nach Bluttransfusion oder i.v. Drogenabusus,
    \item \gEs{oft chronischer Verlauf} (in 70 -- 80~\% der
    Fälle)
\end{itemize}
\\\bottomrule
\end{tabular}
\end{gWideParEnv}
\end{table}



\gLayout{0}{Impfung}
{
Für die Hepatitis A und B ist eine Impfung
verfügbar (z.B. Twinrix{\texttrademark}). Für die Hepatitis
C ist derzeit noch keine Impfung regulär am Markt. Gegen
die toxischen Formen der Hepatitis gibt es natürlich auch keine Impfung.

Achtung: Non-Responder möglich!
% FEATURE Hepatitis Laborwerte Interpretation einfügen.
}
{
\begin{itemize}
    \item Für Hep. A \& B.
    \item Nicht für Hep. C.
\end{itemize}
}
{}
{}
{}


\gFazit{Für den Rettungsdienst sind besonders die Hepatitis
B und C wichtig: Hohe Infektiosität, chronischer Verlauf!

Gegen Hepatitis C gibt es derzeit keine Impfung!

Für Risikopersonal (dazu zählt auch Sanitätspersonal)
übernimmt i.d.R. die Unfallversicherungsanstalt die
Kosten für die Schutzimpfung.}



%%%%%%%%%%%%%%%%%%%%%%%%%%%%%%%%%%%%%%%%%%%%%%%%%%%%%%%%%%%
%%%%%%%%%%%%%%%%%%%%%%%%%%%%%%%%%%%%%%%%%%%%%%%%%%%%%%%%%%%
%%%%%%%%%%%%%%%%%%%%%%%%%%%%%%%%%%%%%%%%%%%%%%%%%%%%%%%%%%%
\subsection{Sepsis}
\label{topic:sepsis}

Siehe auch: Septischer Schock und Toxischer Schock,
Kap.\,\ref{topic:schock-septischer}.

\begin{itemize}
\item Wenn Keime und deren Produkt außer Kontrolle geraten,
erfolgt Amoklauf des Körpers:

\item Ungehemmte Freisetzung von Stoffen des Entzündungs-, Gerinnungs-
und Immunsystems

\item Auch ohne Infektion möglich:


\begin{itemize}
\item \gZ{Systemic Inflammatory Response Syndrome (\gAbk{SIRS})} 
\item nach schwerem Trauma, Verbrennungen, Gewebshypoxie
\end{itemize}\end{itemize}


\gLayout{2}{\gTxSymptome}
{
Die Symptome sind sehr vom Fortschritt und
körperlichen Abbau abhängig. Grundsätzlich:


\begin{itemize}
    \item Sehr stark reduzierter Allgemeinzustand
    \item Fieber ({\textgreater}\,38{\textdegree} C) oder
    Hypothermie ({\textless}\,36{\textdegree} C)
    \item Herzfrequenz {\textgreater}\,90/min
    \item Atemfrequenz {\textgreater}\,20/min oder
    Hyperventilation
    \item unkontrollierbare Gerinnung/Blutung
    \item Hypotonie ${\rightarrow}$ Septischer Schock
    \item Organversagen
\end{itemize}

}{}{}{}{}
\paragraph{} 


%%%%%%%%%%%%%%%%%%%%%%%%%%%%%%%%%%%%%%%%%%%%%%%%%%%%%%%%%%%
%%%%%%%%%%%%%%%%%%%%%%%%%%%%%%%%%%%%%%%%%%%%%%%%%%%%%%%%%%%
%%%%%%%%%%%%%%%%%%%%%%%%%%%%%%%%%%%%%%%%%%%%%%%%%%%%%%%%%%%





\clearpage
%%%%%%%%%%%%%%%%%%%%%%%%%%%%%%%%%%%%%%%%%%%%%%%%%%%%%%%%%%%
%%% Abdomen                                             %%%
%%%%%%%%%%%%%%%%%%%%%%%%%%%%%%%%%%%%%%%%%%%%%%%%%%%%%%%%%%%

\section{Abdominelle Erkrankungen}



\gLayout{22}{Der Bauch als Herausforderung}
{
Der Bauch ist eine echte Herausforderung. Er
beinhaltet viele unterschiedliche Organe und Strukturen
welche weh tun und
erkranken können. Viele dieser, oft auch mit starken Schmerzen verbundenen, Störungen
sind nicht übermäßig schlimm, einige sind jedoch sehr gefährlich oder können rasch
gefährlich werden. 
}
{}
{./images/operation.jpg}
{\gAuthNotesPicLegalOk{}{}Oft die
    Abteilung der Wahl für einen abdominellen Notfall: Die
    Chirurgie. Aber auch die Abteilungen für
    Kinderheilkunde, Kinderchirurgie, Gynäkologie (!) oder
    die Interne können manchmal passende Ziele sein.
    }{Gabriel}
% \begin{wrapfigure}{r}{5cm}
%     \includegraphics[width=\linewidth]{}
%     \gCaptionof{figure}{
%     \gSourcePicture{}}
% \end{wrapfigure}


\paragraph{Anamnese und Untersuchung} Es ist daher wichtig
typische Symptome zu erheben und Alarmzeichen sofort zu erkennen.

Wesentlich ist bei allen unklaren Bauchbeschwerden die Erhebung einer exakten Stuhlanamnese:

\begin{itemize}
    \item Auffälligkeiten? (Geruch, Farbe, Konsistenz, Frequenz)
    \item Wann erfolgte der letzte Stuhlgang?
    \item Winde?
    \item Genaue Beschreibung von Harn und eventuell Erbrochenem.
    \item Bei Frauen im gebärfähigem Alter Regelanamnese.
\end{itemize}

Der Patient muss bis zur endgültigen Klärung der Situation
\gEs{nüchtern} belassen werden!

Wenn es der Patient toleriert, muss eine Betrachtung und
Abtastung des unbekleideten Abdomens beim liegenden
Patienten erfolgen wie unter
\ref{topic:untersuchung-abdomen-abtasten} beschrieben.

\subsection{Der schmerzende Bauch}

\gLayout{0}{Der Bauch -- unendliche Weiten}
{
Gleich vorweg: In vielen
Fällen wird der Grund des Bauchschmerzes nicht mit
hinreichender Sicherheit feststellbar sein und es bietet sich
keine sichere Verdachtsdiagnose an. Statt einer Diagnose wird
dann das \gEs{Beschwerdebild beschrieben} (Ort des Schmerzes, Art des
Schmerzes, Dauer, Verlauf, Austrahlung, Ergebnis der
abdominellen Tastuntersuchung: Abdomen weich oder hart,
Druckschmerz, Loslassschmerz).
}{
\begin{itemize}
    \item In vielen Fällen Krankheitsbild unklar.
    \item Im Zweifel: Beschreibung statt Diagnose.
\end{itemize}
}{}{}{}

\gLayout{2}{Beispiele}
{
\begin{quote}
{\sffamily\bfseries Diagnose:} {\ttfamily Unklarer
Bauchschmerz: Kolikartige Schmerzen im linken Unterbauch seit 1/2h }

{\sffamily\bfseries Anmerkungen:} {\ttfamily keine Ausstrahlung, Druckschmerz im li. Unterbauch, Abdomen sonst  weich.}
\vspace{0.5em}

{\sffamily\bfseries Diagnose:} {\ttfamily Unklarer
Bauchschmerz: Brennend-stechender Schmerz im mittleren
Oberbauch seit 2h mit Ausstrahlung in den Rücken}

{\sffamily\bfseries Anmerkungen:} {\ttfamily Abdomen weich,
keine Druck- oder Loslassschmerzen}
\vspace{0.5em}

{\sffamily\bfseries Diagnose:} {\ttfamily Unklarer
Bauchschmerz: „Dumpfer“ Schmerz im linken und rechten
Oberbauch seit gestern. }

{\sffamily\bfseries Anmerkungen:} {\ttfamily Abdomen
weich, keine Druck- oder Loslassschmerzen}
\end{quote}
}{

}{}{}{}



Dennoch muss man bei jeden Patienten versuchen, so gut als
möglich die in Frage kommenden Krankheitsbilder einzugrenzen
und gefährliche Erkrankungen nach Möglichkeit auszuschließen. Im
folgenden sind einige häufige Krankheitsbilder beschrieben,
die \gE{erkennbar}  und relativ häufig sind.

\gMassnahmenNG{Allgemeine Maßnahmen}{Abdominalerkrankung}
{\label{m:am-abdominalerkrankungen}}{
\begin{itemize}
    \item \gEs{Vitale Bedrohung einschätzen.} \gTxBeiVit
    \gTxMassVit
    \item Lagerung: grundsätzlich bauchdeckenentspannend
    (Knierolle, Beine angezogen), oder nach Wunsch des Patienten
    \item nüchtern lassen (in Hinblick auf möglicherweise bevorstehende
    Operation, nichts essen oder trinken lassen)
%     \item Monitoring. Der Zustand kann sich abrupt
%     verschlechtern.
    \item Hospitalisierung: Situationsgerecht. Bei
    weiblichen Patienten Abt. für Gynäkologie (evt. im >Hintergrund<) erwägen.
\end{itemize}
}


\subsection{Häufige und gut erkennbare Erkrankungen}



\subsubsection{Das Akute Abdomen --- Alarmzeichen mit
bretthartem Bauch}\label{topic:akutes-abdomen}



\gDefinition{\gT{Akutes Abdomen}: Akute lebensbedrohliche
Erkrankung des Bauchraumes mit plötzlichen, starken
\gEs{Bauchschmerzen} und
\gEs{bretthartem Bauch}. Es sind viele Ursachen  möglich!}

\gZ{Die Definition des Akuten Abdomens ist oft nicht ganz
eindeutig: Einerseits kann man da\-runter jede plötzlich auftretende
Baucherkrankung verstehen\footnote{Longmore et.al.: Oxford Handbook of
Clinical Medicine, 7\textsuperscript{th} edt., Oxford University Press,
2007\cite{Longmore2007}}\footnote{\cite{Renz-Polster2006}
S.611}.} 

Allerdings ver\-stehen viele Leute darunter eher
eine bedrohliche Bauch\-erkrankung, bei der man akut (schnell) handeln muss (Ab\-wehrspannung, harter Bauch). In der
Notfallmedizin halten wir uns an letztere Definition. 


\gLayout{2}{\gTxBeschreibung}
{Das „\gT{Akute Abdomen}“ ist ein interventionspflichtiger
Symptomenkomplex, der auf einem sich akut oder chronisch
entwickelnden (Stunden, Tage oder auch Wochen) pathologischen
Prozess im Bauchraum beruht.}
{

}{}{}{}

Es hat grundverschiedene Ursachen, die eine gemeinsame
Symptomatik haben.

\gLayout{0}{\gTxSymptome}
{
Der Patient liegt meistens im Bett. Oft ist
eine Schonhaltung (angezogene Beine, gekrümmte
Körperhaltung) zu beobachten. Er klagt über Schmerzen im
Bauch die können als diffus (im gesamten Bauch) oder auf ein Bauchsegment
beschränkt, angegeben werden. 

Die harte Bauchdecke („\gE{bretthart}“) ist neben den
Schmerzen das Leitsymptom des Akuten Abdomens. 

Es besteht oft Übelkeit oder auch Erbrechen. Der Kreislauf
kann kompensiert sein, oder Zeichen von Schwäche zeigen,
bis hin zu einer Schocksymptomatik (Blutverlust, Sepsis).
} 
{
\begin{itemize}
    \item \gEs{Schmerz}: diffus (ganzer Bauch) oder
    lokalisierbar
    \item \gEs{Harte  Bauchdecke} (bretthart)
    \item Schonhaltung  (angezogene Beine, Patient
    krümmt sich)
    \item Übelkeit, Erbrechen
    \item evt. Stuhl-, Windverhalten
    \item evt. Fieber
    \item evt. Schock
\end{itemize}
}{}{}{}


Diese Auflistung der Symptome ist in der Praxis nicht
unbedingt vollständig bei einem Patienten nachzuweisen. Die
Intensität und die Anzahl der angeführten Symptome ist
einerseits vom Allgemeinzustand, vom Lebensalter, von
anderen vorbestehenden Erkrankungen des Patienten und der
Dauer des aktuellen Zustandes abhängig.

Wesentlich nach Eintreffen beim Patienten, ist die
Beantwortung der Fragen \emph{"`Wie schlecht, bzw. wie gut
geht es dem Patienten? Ist der Patient akut vital
bedroht?"'}.


\gLayout{0}{Ursachen}
{
Grundsätzlich kommt fast jede Erkrankung oder Verletzung in
Betracht die Strukturen und Organe im Bauchraum betrifft,
entsprechend vielfältig sind auch die Ursachen: 

\begin{multicols}{2}
\begin{itemize}
    \item Bauchfellentzündung (Peritonitis)
    \item Darmverschluss (Ileus)
    \item Blinddarmentzündung (Appendizitis)
    \item Entzündung der Bauchspeicheldrüse  (Pankreatitis)
    \item Entzündung des Magens / Magengeschwür
    \item Durchbruch / Zerreißung eines Hohlorganes
    \item Blutungen in den Bauchraum
    \item Bauchaortenaneurysma
    \item Mesenterialinfarkt
    \item andere abdominelle Erkrankungen
    \item Entzündung der Leber (Hepatitis )
    \item Gynäkologische Erkrankungen
    \item \ldots
\end{itemize}
\end{multicols}
}{
\begin{itemize}
    \item Bauchfellentzündung (Peritonitis)
    \item Blinddarmentzündung (Appendizitis)
    \item Entzündung der Bauchspeicheldrüse  (Pankreatitis)
    \item Durchbruch / Zerreißung eines Hohlorganes
    \item Blutungen in den Bauchraum
    \item Bauchaortenaneurysma
    \item andere abdominelle Erkrankungen
    \item Gynäkologische Erkrankungen
    \item u.v.a.m. \ldots
\end{itemize}
}{}{}{}





\gMassnahmenNG{Spezielle Maßnahmen}{Akutes Abdomen}
{\label{m:am-abdominalerkrankungen}}
{
\begin{itemize}
  \item Allgemeine Maßnahmen bei Abdominalerkrankungen
  (\prettyref{m:am-abdominalerkrankungen})
    \item \gEs{Vitale Bedrohung einschätzen.} \gTxBeiVit
    \gTxMassVitBes
    \begin{itemize}
        \item \ce{O2} 10\gLMin
        \item Lagerung: je nach Situation: bauchdeckenentspannend,
        Beine hoch, \ldots
    \end{itemize}
    \item Lagerung: grundsätzlich bauchdeckenentspannend
    (Knierolle, Beine angezogen), oder nach Wunsch des Patienten
    \item nüchtern lassen (in Hinblick auf möglicherweise bevorstehende
    Operation, nichts essen oder trinken lassen)
    \item Monitoring. Der Zustand kann sich abrupt
    verschlechtern.
    \item Hospitalisierung: Abt. f. Chirurgie, bei Kindern ${\rightarrow}$
    Abt. f. Kinderheilkunde, bei Frauen evt. Gynäkologie oder Chirurgie mit
    Abt. f. Gynäkologie im Haus \opt{REGVIE}{(z.B. Donauspital, Rudolfstifung,
    Wilhelminenspital, etc.)}
\end{itemize}
}


\subsubsection{Darmverschluss -- Ileus}\index{Ileus}

\gDefinition{Unterbrechung der normalen  Darmpassage}

\gZ{\index{Ileus}\gT{Ileus}  (vollständig)

\index{Subileus}\gT{Subileus}  (unvollständig)}

\gLayout{2}{\gTxBeschreibung}
{

Der Verschluss kann mechanisch bedingt sein (durch
Hindernisse wie Verwachsungen oder Tumore), oder durch eine
Lähmung der Darmmuskulatur (Medikamente, Gifte, Entzündung
\ldots) verursacht sein. Bestimmte Schmerzmittel
\gZ{(Opiate)}, welche oft bei Krebspatienten eingesetzt
werden, sind bekannt dafür Darmlähmungen zu erzeugen. 
}{
\begin{itemize}
  \item Mechanische Blockade oder Lähmung.
\end{itemize}
}{}{}{}


\gLayout{0}{\gTxSymptome}
{
Typisch für den Darmverschluss ist die mangelhafte
Stuhlproduktion. Es kommt zu diffusen Bauchschmerzen und einem
Blähbauch, sowie zu Übelkeit und Erbrechen bis hin zum
Erbrechen von zurück gestautem Kot (selten).
}{
\begin{itemize}
    \item Diffuse Bauchschmerzen
    \item Stuhlverhalten
    \item Übelkeit, Erbrechen (im schlimmsten Fall Koterbrechen)
\end{itemize}
}{}{}{}



\gMassnahmenNG{Spezielle Maßnahmen}{Darmverschluß}{}{
\begin{itemize}
  \item Allgemeine Maßnahmen bei Abdominalerkrankungen
  (\prettyref{m:am-abdominalerkrankungen})
    \item \gEs{Vitale Bedrohung einschätzen.} \gTxBeiVit
    \gTxMassVitBes
    \begin{itemize}
        \item \ce{O2} 10\gLMin
        \item Lagerung: je nach Situation: bauchdeckenentspannend,
        Beine hoch, \ldots
    \end{itemize}
    \item Patient nüchtern lassen.
    \item Hospitalisierung: Abt. f. Chirurgie.
\end{itemize}
}


\subsubsection{Blinddarmentzündung -- Appendizitis}



\gDefinition{Entzündung des \gEs{Wurmfortsatzes}
(\gEs{Appendix} \gZ{vermiformis}) des Blinddarms}

\gLayout{0}{\gTxSymptome{} und Verlauf}
{
\gACh{GABG}{2009-05-07}{NEU}Im Vordergrund der akuten Appendizitis
steht anfänglich meist ein als eher diffus empfundener, heftiger Schmerzzustand im Bereich der
Nabelgegend\gACh{GABS}{2009-05-07}{Magengegend --> Nabelgegend}, der
innerhalb weniger Stunden in den rechten Unterbauch
wandert. Übelkeit, Brechreiz und Verstopfung
\gACh{GABS}{2009-05-07}{Obstipation --> Verstopfung} können das Krankheitsbild begleiten. Meist ist die Körpertemperatur bis 38\textcelsius{} leicht erhöht
\gACh{GABS}{2009-05-07}{gelöscht da unerheblich für RS: , wobei eine rektale und
axilläre Messung eine Differenz um ca. 1\textcelsius aufweist}. 
\gACh{GABS}{2009-05-07}{alt: Unbehandelt kann der weitere Verlauf durch eine
Abszeßbildung und Perforation in
die Bauchhöhle gekennzeichnet sein. neu: Unbehandelt kann es zu einer Eiterung
und einem Durchbruch in die Bauchhöhle kommen.} Unbehandelt kann es zu einer
Eiterung und einem Durchbruch in die Bauchhöhle kommen. In diesem Stadium entwickelt sich ein
\gACh{GABS}{2009-05-07}{alt: Schockzustand, der eine absolute Lebensbedrohung
darstellt; neu: akutes Abdomen, siehe \ref{topic:akutes-abdomen}} akutes
Abdomen (\prettyref{topic:akutes-abdomen}). 

Der Verlauf einer akuten Appendizitis ist
jeweils durch individuelle Faktoren wie Lebensalter und Allgemeinzustand des
Patienten geprägt. 
}{
\begin{itemize}
    \item Schmerzen im re. Unterbauch
    \item \gE{Druck}schmerzen im re. Unterbauch
    \item \gE{Loslass}schmerz li. Bauch
    \item Abwehrspannung
    \item Übelkeit, Erbrechen
    \item evt. Fieber
\end{itemize}
}{}{}{}

\gMassnahmenNG{Spezielle Maßnahmen}{Appendizitis}{}{
\begin{itemize}
  \item Allgemeine Maßnahmen bei Abdominalerkrankungen
  (\prettyref{m:am-abdominalerkrankungen})
    \item Patient nüchtern lassen.
    \item Hospitalisierung: Abt. f. Chirurgie; bei Kindern
    Kinderchirurgie oder Abt. f. Kinderheilkunde (je nach
    regionaler Regelung)
\end{itemize}
}




\subsubsection{Gallenkolik}
\label{topic:gallenkolik}
\index{Kolik!Gallen-}
\index{Gallenkolik}

\gDefinition{Kolik-artige Schmerzen infolge einer Blockade
der Gallenwege durch Gallensteine}

\gZ{\emph{Vgl. Kolik: \gSee{topic:kolik}}}

\gLayout{0}{\gTxBeschreibung} 
{Gallensteine verursachen eine Verstopfung des
Gallengangs und folglich einen Rückstau von
Gallenflüssigkeit in die Leber.
} 
{
\begin{itemize}
    \item Verstopfung der Gallenwege (Stein)
    \item → Rückstau und Dehnung
\end{itemize}
}{}{}{}

\gLayout{0}{\gTxSymptome}
{
Durch den Rückstau und der Dehnung des Gallengangs kommt es zu
starken, auf- und abschwellenden, krampfartigen
(\gEs{kolikartigen}) \gEs{Schmerzen} im \gEs{rechten Oberbauch},
klassischerweise mit Ausstrahlung in die rechte Schulterregion. Begleitendes Fieber ist möglich.

Der rechte Oberbauch ist äußerst druckschmerzhaft, der
Patient ist sehr \gEs{unruhig} und agitiert.
}{
\begin{itemize}
    \item Auf- und abschwellende, krampfartige
    (kolikartige) Schmerzen im re. Oberbauch
    \item Besonders nach fettreichen Mahlzeiten
    \item evt. Gelbfärbung des Augenweiß (Gelbsucht (Ikterus))
    \item evt. Übelkeit, Erbrechen
\end{itemize}
}{}{}{}

\gMassnahmenNG{Spezielle Maßnahmen}{Gallenkolik}{}{
\begin{itemize}
  \item Allgemeine Maßnahmen bei Abdominalerkrankungen
  (\prettyref{m:am-abdominalerkrankungen}).
  \item Notarzt zur Schmerztherapie erwägen.
    \item Patient nüchtern lassen.
    \item Hospitalisierung: Abt. f. Chirurgie; bei Kindern
    Kinderchirurgie oder Abt. f. Kinderheilkunde (je nach
    regionaler Regelung).
\end{itemize}
}


\subsubsection{Ösophagusvarizenblutung}

\gDefinition{Blutende "`Krampfadern"' in der
Speiseröhre.}\nopagebreak[4]

\gLayout{0}{\gTxBeschreibung}
{Bei Patienten mit einer fortgeschrittenen Lebererkrankung
(Leberzirrhose, \ref{topic:leberzirrhose}) kann es zur
Bildung von Krampfader-artigen Gefäßen in der Speiseröhre
kommen. Diese Gefäße können massive Blutungen verursachen
die tödlich enden können. 

Je nach Stärke  äußern sich die Blutungen durch
schwallartiges Erbrechen von Kaffeesatz-artigem
Mageninhalt\footnote{\gT{Kaffeesatz-artiges Erbrechen}:
Durch die Magensäure kommt es zur Veränderung des Blutes,
es wird bräunlich, ähnlich dem Kaffeesatz im Kaffeefilter.}
(wenn das Blut von der Magensäure bereits verändert wurde) oder von rotem, frischen Blut.
}
{
\begin{itemize}
    \item Oft bei Lebererkrankungen
    \item Hoher Blutverlust möglich -- Lebensgefahr!
\end{itemize}
}{}{}{}

\gLayout{0}{\gTxSymptome}
{
\begin{itemize}
    \item Bluterbrechen im Schwall (Hämatemesis)
\end{itemize}
}{
\begin{itemize}
    \item Bluterbrechen im Schwall (Hämatemesis)
\end{itemize}
}{}{}{}



\gFazit{Der Patient ist aufgrund des Blutverlustes akut
vital bedroht, es besteht Schockgefahr.}

\gMassnahmenNG{Spezielle Maßnahmen}{Oesophagusvarizenblutung}{}{
    
    \begin{itemize}
    \item \gTxMassVitBes
    \begin{itemize}
        \item \ce{O2} 10-15\gLMin
        \item Lagerung:
        Situationsgerecht \gAuthNotes{Lagerung? Beine hoch?
        OK hoch?}
    \end{itemize}
    \item Allgemeine Maßnahmen der Schockbekämpfung
    (\prettyref{m:am-schock})
    \item Spezielle Maßnahmen bei Hypovolämischem Schock
    (\prettyref{topic:schock-hypovolaemischer})
    \item \gZ{Notarzt: evt. Sengstaken-Blakemore-Sonde}
    \item Hospitalisation: Abteilung für Chirurgie, Voranmeldung
\end{itemize}
}






\subsubsection{Magen-Darm-Grippe -- Gastroenteritis}

\emph{Syn.: Brechdurchfall, \gZ{Gastroduodenitis}}

Die Magen-Darm-Grippe ist eine chronisch oder akut
verlaufende Entzündung der Magen- und der Dünndarmschleimhaut. 

\gLayout{0}{\gTxBeschreibung{} und Ursachen}
{
Die Gastroenteritis wird durch Erreger
wie Bakterien oder Viren hervorgerufen. Es gibt völlig
harmlose aber auch lebensbedrohliche Verlaufsformen, dies ist abhängig vom Erreger und
letztendlich auch vom Alter des Patienten. Durch den
Durchfall und das Erbrechen kann es zu einem maßgeblichen
Verlust von Wasser und Elektrolyten kommen. Wenn der Patient
nicht ausreichend trinken kann besteht die Gefahr der Austrocknung (Exsikkose).
}{
\begin{itemize}
    \item Vielfältig.
    \item Häufig Infektionen.
\end{itemize}
}{}{}{}


\gLayout{2}{\gTxSymptome}
{
\begin{itemize}
    \item Erbrechen (Emesis)
    \item Durchfall (Diarrhoe)
    \item Wasser-/Elektrolytverlust
\end{itemize}
}{
\begin{itemize}
    \item Erbrechen (Emesis)
    \item Durchfall (Diarrhoe)
    \item Wasser-/Elektrolytverlust
\end{itemize}
}{}{}{}

\gLayout{0}{Besonders gefährdet}
{
Besonders betroffen davon sind Kleinkinder, da sie
normalerweise einen höheren Wasseranteil am Körpergewicht
haben.  Die Symptome sind klassisch: Oberbauchkrämpfe,
die ziemlich heftig sein können, Übelkeit, Erbrechen, Durchfälle.
Exsikkose-Gefahr! 
}{
\begin{itemize}
    \item (Klein-)Kinder
    \item alte Menschen
\end{itemize}
}{}{}{}



\gMassnahmenNG{Spezielle Maßnahmen}{Gastroenteritis}{}{
\begin{itemize}
  \item Allgemeine Maßnahmen bei Abdominalerkrankungen
  (\prettyref{m:am-abdominalerkrankungen})
    \item Patient nüchtern lassen.
    \item Hospitalisierung: Abt. f. Innere Medizin; bei
    Kindern Abt. f. Kinderheilkunde.
    \item Strikte Desinfektionsmaßnahmen! 
\end{itemize}
}







\subsubsection{Flüssigkeitsmangel -- Exsikkose}
\gACh{WITM}{????-??-??}{NEU: Schlagworte}
\gACh{GABG}{2009-05-11}{NEU: Fließtexte}
\gACh{GABS}{2009-05-14}{Fließtexte stark modifiziert und angepasst. }

\emph{Syn.: Dehydratation}

\gDefinition{„Austrocknung“: Es besteht ein hochgradiges Defizit von Flüssigkeit im Körper.}

Grundsätzlich wird (wurde) entweder \gE{zu wenig Flüssigkeit zugeführt}
oder \gE{zu viel Flüssigkeit an die Außenwelt abgegeben} (eine Kombination aus
beidem ist natürlich auch möglich). Dafür gibt es jeweils sehr verschiedene Ursachen: 

\gLayout{0}{Ursachen}
{
Ältere Menschen haben oft ein mangelndes Durstgefühl und daher
einen verminderten Trinkantrieb. Auch durch die Einnahme von
Drogen (Ecstacy) kann auf den Durst "`vergessen"' werden.  

Andererseits kann es durch \gE{Erbrechen}, \gE{Durchfall}
sowie \gE{starkes Schwitzen} zu einem stark erhöhten
Flüssigkeitsverlust kommen. Bei \gE{Fieber} wird auch
vermehrt Feuchtigkeit abgeatmet, und natürlichen können
auch \gE{Medikamente} "`entwässern"'.
}{
\begin{itemize} %GABS: umgebaut
    \item Flüssigkeitszufuhr ${\downarrow}$
    \begin{itemize}
        \item Verminderter Trinkantrieb
    \end{itemize}
    \item Flüssigkeitsverlust ${\uparrow}$
    \begin{itemize}
        \item Erbrechen, Durchfall
        \item Starkes Schwitzen
        \item Fieber (Feuchtigkeit wird vermehrt abgeatmet)
        \item (Entwässerungs)medikamente
    \end{itemize}
\end{itemize}
}{}{}{}

Die genannten Ursachen sind natürlich nur ein kleiner Teil von sehr vielen
Möglichkeiten.

\gLayout{0}{\gTxSymptome}
{
Der Patient wirkt allgemein \gE{geschwächt bis apathisch},
oft \gEs{verwirrt}, klagt über mehr oder minder heftiges
Durstgefühl, ev. über Wadenkrämpfe und Erbrechen. Die Haut
ist auffallend trocken, \gEs{Falten sind abhebbar und
bleiben bestehen}. \gEs{Mundhöhle und Zunge sind trocken},
evt. borkig oder mit eingetrocknetem Schleim belegt. Die Augen sind charakteristischerweise tiefliegend, die Nase wirkt auffallend spitz. Der RR ist mitunter sehr nieder,
die HF eher hoch. \gE{Der Patient ist bis zur Klärung der
Situation nüchtern zu belassen!}
}{
\begin{itemize}
    \item RR ${\downarrow}$
    \item stehende Hautfalten
    \item trockene, belegte Zunge
    \item Elektrolytentgleisung
    \item \gEs{Verwirrtheit, Schwindel, Apathie}
\end{itemize}
}{}{}{}


% \par
% \hrule
% \par
% 
%
% \subsubsection{Flüssigkeitsmangel -- Exsikkose}
% \gACh{WITM}{????-??-??}{NEU: Schlagworte}
% \gACh{GABG}{2009-05-11}{NEU: Fließtexte}
% \gACh{GABS}{2009-05-14}{Fließtexte stark modifiziert und angepasst. }
% 
% \emph{Syn.: Dehydratation}
% 
% \gDefinition{„Austrocknung“: Es besteht ein hochgradiges Defizit von Flüssigkeit im Körper.}
% 
% 
% \gTagPar{Ursachen}Grundsätzlich wird (wurde) entweder
% \gEs{zu wenig Flüssigkeit zugeführt} oder \gEs{zu viel Flüssigkeit an die
% Außenwelt abgegeben} (eine Kombination aus beidem ist natürlich auch möglich). Dafür gibt es jeweils sehr verschiedene Ursachen: 
% 
% Ältere Menschen haben oft ein mangelndes Durstgefühl und daher
% einen verminderten Trinkantrieb. Auch durch die Einnahme von
% Drogen (Ecstacy) kann auf den Durst "`vergessen"' werden.  
% 
% Andererseits kann es durch \gE{Erbrechen}, \gE{Durchfall}
% sowie \gE{starkes Schwitzen} zu einem stark erhöhtem
% Flüssigkeitsverlust kommen. Bei \gE{Fieber} wird auch
% vermehrt Feuchtigkeit abgeatmet, und natürlichen können
% auch \gE{Medikamente} "`entwässern"'.
% 
% 
% 
% Die genannten Ursachen sind natürlich nur ein kleiner Teil von sehr vielen
% Möglichkeiten.
% 
% \gTagPar{\gTxSymptome}Der Patient wirkt allgemein \gE{geschwächt
% bis apathisch}, oft \gEs{verwirrt}, klagt über mehr oder minder heftiges
% Durstgefühl, ev. über Wadenkrämpfe und Erbrechen. Die Haut
% ist auffallend trocken, \gEs{Falten sind abhebbar und
% bleiben bestehen}. \gEs{Mundhöhle und Zunge sind trocken},
% evt. borkig oder mit eingetrocknetem Schleim belegt. Die Augen sind charakteristischerweise tiefliegend, die Nase wirkt auffallend spitz. Der RR ist mitunter sehr nieder,
% die HF eher hoch. \gE{Der Patient ist bis zur Klärung der
% Situation nüchtern zu belassen!} 
% 
% 
% \subsubsection{Flüssigkeitsmangel -- Exsikkose}
% \gACh{WITM}{????-??-??}{NEU: Schlagworte}
% \gACh{GABG}{2009-05-11}{NEU: Fließtexte}
% \gACh{GABS}{2009-05-14}{Fließtexte stark modifiziert und angepasst. }
% 
% \emph{Syn.: Dehydratation}
% \gA{Version mit multicols}
% 
% \gDefinition{„Austrocknung“: Es besteht ein hochgradiges Defizit von Flüssigkeit im Körper.}
% 
% \begin{multicols}{2}
% \gTagPar{Ursachen}Grundsätzlich wird (wurde) entweder
% \gEs{zu wenig Flüssigkeit zugeführt} oder \gEs{zu viel Flüssigkeit an die
% Außenwelt abgegeben} (eine Kombination aus beidem ist natürlich auch möglich). Dafür gibt es jeweils sehr verschiedene Ursachen: 
% 
% Ältere Menschen haben oft ein mangelndes Durstgefühl und daher
% einen verminderten Trinkantrieb. Auch durch die Einnahme von
% Drogen (Ecstacy) kann auf den Durst "`vergessen"' werden.  
% 
% Andererseits kann es durch \gE{Erbrechen}, \gE{Durchfall}
% sowie \gE{starkes Schwitzen} zu einem stark erhöhtem
% Flüssigkeitsverlust kommen. Bei \gE{Fieber} wird auch
% vermehrt Feuchtigkeit abgeatmet, und natürlichen können
% auch \gE{Medikamente} "`entwässern"'.
% 
% 
% 
% Die genannten Ursachen sind natürlich nur ein kleiner Teil von sehr vielen
% Möglichkeiten.
% 
% \gTagPar{\gTxSymptome}Der Patient wirkt allgemein \gE{geschwächt
% bis apathisch}, oft \gEs{verwirrt}, klagt über mehr oder minder heftiges
% Durstgefühl, ev. über Wadenkrämpfe und Erbrechen. Die Haut
% ist auffallend trocken, \gEs{Falten sind abhebbar und
% bleiben bestehen}. \gEs{Mundhöhle und Zunge sind trocken},
% evt. borkig oder mit eingetrocknetem Schleim belegt. Die Augen sind charakteristischerweise tiefliegend, die Nase wirkt auffallend spitz. Der RR ist mitunter sehr nieder,
% die HF eher hoch. \gE{Der Patient ist bis zur Klärung der
% Situation nüchtern zu belassen!} 
% \end{multicols}
% 
% 
% \subsubsection{Flüssigkeitsmangel -- Exsikkose}
% \gACh{WITM}{????-??-??}{NEU: Schlagworte}
% \gACh{GABG}{2009-05-11}{NEU: Fließtexte}
% \gACh{GABS}{2009-05-14}{Fließtexte stark modifiziert und angepasst. }
% 
% \emph{Syn.: Dehydratation}
% 
% \gDefinition{„Austrocknung“: Es besteht ein hochgradiges Defizit von Flüssigkeit im Körper.}
% 
% 
% \gMarried{
% \gTagPar{Ursachen}Grundsätzlich wird (wurde) entweder
% \gEs{zu wenig Flüssigkeit zugeführt} oder \gEs{zu viel Flüssigkeit an die
% Außenwelt abgegeben} (eine Kombination aus beidem ist natürlich auch möglich). Dafür gibt es jeweils sehr verschiedene Ursachen: 
% 
% Ältere Menschen haben oft ein mangelndes Durstgefühl und daher
% einen verminderten Trinkantrieb. Auch durch die Einnahme von
% Drogen (Ecstacy) kann auf den Durst "`vergessen"' werden.  
% 
% Andererseits kann es durch \gE{Erbrechen}, \gE{Durchfall}
% sowie \gE{starkes Schwitzen} zu einem stark erhöhtem
% Flüssigkeitsverlust kommen. Bei \gE{Fieber} wird auch
% vermehrt Feuchtigkeit abgeatmet, und natürlichen können
% auch \gE{Medikamente} "`entwässern"'.
% 
% 
% 
% Die genannten Ursachen sind natürlich nur ein kleiner Teil von sehr vielen
% Möglichkeiten.
% }{
% \gTagPar{\gTxSymptome}Der Patient wirkt allgemein \gE{geschwächt
% bis apathisch}, oft \gEs{verwirrt}, klagt über mehr oder minder heftiges
% Durstgefühl, ev. über Wadenkrämpfe und Erbrechen. Die Haut
% ist auffallend trocken, \gEs{Falten sind abhebbar und
% bleiben bestehen}. \gEs{Mundhöhle und Zunge sind trocken},
% evt. borkig oder mit eingetrocknetem Schleim belegt. Die Augen sind charakteristischerweise tiefliegend, die Nase wirkt auffallend spitz. Der RR ist mitunter sehr nieder,
% die HF eher hoch. \gE{Der Patient ist bis zur Klärung der
% Situation nüchtern zu belassen!} 
% }
% 




\subsection{Weniger häufige oder nicht-so-gut-erkennbare Erkrankungen}

\minisec{Querverweise}

\begin{itemize}
    \item Hepatitis: \gSee{topic:hepatitis}
\end{itemize}




\subsubsection{Magen- und Zwölffingerdarmgeschwür}

\gZ{\emph{Syn. Magengeschwür: Gastritis}}

\gDefinition{Magengeschwür: Entzündung der
Magenschleimhaut mit Substanzdefekt.}

\gLayout{2}{Ursachen}
{
\begin{itemize}
    \item Bakteriell \gZ{(Helicobacter pylori)}
    \item Medikamente\footnote{Nicht-steroidale Antirheumatika (NSAR), wie z.B.
    Voltaren{\texttrademark} (Diclofenac), Aspirin{\texttrademark}
    (Acetylsalicylsäure, ASS),etc., schädigen die Magenschleimhaut.}
    \item Chemisch-toxisch.
    \item \gZ{Gallensäurenrückfluß (Reflux)}
    \item \gZ{Autoimmunerkrankungen}
\end{itemize}
}{
\begin{itemize}
    \item Bakteriell \gZ{(Helicobacter pylori)}
    \item Medikamente
    \item Chemisch-toxisch.
    \item \gZ{Gallensäurenrückfluß (Reflux)}
    \item \gZ{Autoimmunerkrankungen}
\end{itemize}
}{}{}{}

\gLayout{0}{\gTxSymptome}
{\gAuthNotes{Fließtext fehlt: insbes. \underline{Wie} ist
der Schmerz, wann}
Es kommt zu Appetitlosigkeit, Völlegefühl, Übelkeit. Wenn
gleichzeitig eine Blutung besteht, dann sind Teerstühle
oder kaffeesatz-artiges Erbrechen zu beobachten.
}
{
\begin{itemize}
    \item Schmerzen (Oberbauch), Übelkeit, Erbrechen
    \item bei Blutung: Kaffesatz-Erbrechen,  Teerstuhl
\end{itemize}
}{}{}{}









\subsubsection{Bauchfellentzündung}

\gZ{\emph{Syn.: Peritonitis}}

\gDefinition{Entzündung und Keimbesiedelung des 
Bauchfells}

\gMarried
{Durch die Infektion und Entzündung des Bauchfelles kommt es
zu Symptomen eines Akuten Abdomens. Oft ist die Perforation
eines Hohlorganes  (Magendurchbruch, Darmperforation,
Blinddarmdurchbruch, \ldots) Ursache der Entzündung. }
{\gSlideHeading{Ursache}

\begin{itemize}
    \item Perforation eines Hohlorganes
    \item Infektion über Blut- / Lymphwege oder eingewandert
    über Fremdkörper
\end{itemize}

\gSlideHeading{\gTxSymptome}

\begin{itemize}
    \item Symptome eines Akuten Abdomens
    \item Schmerzen (diffus)
    \item evt. reflektorischer Darmverschluss
    \item evt. Sepsis- und/oder Schockzeichen
\end{itemize}}


% \gZ{
% \subsection{Reserviert \gA{Sonstige -- nicht akute --
% Erkrankungen}}
% \subsubsection{Reserviert \gA{Chronisch entzündliche
% Darmerkrankungen}} }
% 



\section{Urologische Erkrankungen}

\gDefinition{Erkrankungen des harnableitenden Systems.}


\subsection{Nierenkolik}
\label{topic:nierenkolik}
\index{Kolik!Nieren-}
\index{Nierenkolik}

\gDefinition{Blockade des Harnleiters durch Nierensteine
mit massiven kolikartigen Schmerzen aufgrund des Harnstaus.}

\gZ{\emph{Vgl. Kolik: \ref{topic:kolik}}}

\gLayout{0}{Ursachen}{
Die Nierenkolik ist eine äußerst schmerzhafte
akute Erkrankung. Sie entsteht dadurch, dass Nierensteine
den Harnleiter blockieren und verstopfen. Dadurch wird der Harn
gestaut und der Harnleiter gedehnt. 
}{
\begin{itemize}
    \item Nierenstein(e)
\end{itemize}
}{}{}{}

\gLayout{0}{\gTxSymptome}{
Die Folge sind massive,
kolikartige Schmerzen, in der Flanke bzw. im Unterbauch der
jeweiligen Seite. Im Gegensatz zu anderen abdominellen
Erkrankungen ist der Patient extrem unruhig und agitiert.
Oft läuft er gestikulierend und wehklagend durch die
Gegend. Mitunter ist eine entsprechende Vorerkrankung
\gE{(Nierensteine)} bereits bekannt und kann erfragt werden. 
}{
\begin{itemize}
    \item Rücken-/Flankenschmerzen
    \item evt. Unterbauchschmerzen bei  Stein-/Sandabgang
    \item {\quotedblbase}der Stein wandert, der Schmerz wandert, der Patient
    wandert{\textquotedblleft}
    \item Blut im Harn (Hämaturie)
\end{itemize}
}{}{}{}

\gLayout{0}{Versorgung}{
Im Vordergrund bei der Versorgung
steht der rasche zügige Transport an eine Abteilung für Urologie. Wenn die
Schmerzen sehr massiv sind, muss eventuell eine
Schmerztherapie durch den Notarzt erfolgen.
}
{\gSymbolText}{}{}{}


\gMassnahmenNG{Spezielle Maßnahmen}{Nierenkolik}{}{
    \begin{itemize}
        \item Lagerung nach Wunsch des Patienten
        \item schonender Transport
        \item Bei unerträglichen Schmerzen Notarzt zur
        Schmerztherapie beiziehen
        \item Hospitalisation: Abt. für Urologie
    \end{itemize}
}


\subsection{Nierenbeckenentzündung}
\label{topic:nierenbeckenentzuendung}
% TODO INHALT: fehlt komplett: Nierenbeckenentzündung 

\subsection{Harnwegsinfekt}
\label{topic:harnwegsinfekt}
% TODO INHALT: fehlt komplett: Harnwegsinfekt 

\subsection{Akutes Harnverhalten}

\gDefinition{Abflussbehinderung des Harns aus der
Harnblase, verschiedene Ursachen sind möglich.}



Das akute Harnverhalten ist besonders bei Männern im "`besten"'
oder fortgeschrittenem Alter häufig anzutreffen. 

\gLayout{0}{Ursachen}
{
Hauptursache  dafür ist eine
altersbedingte Vergrößerung der Prostata, die die Harnröhre abschnürt. Davon abgesehen
können auch alte Narben, Medikamente oder eine Störung der
Schließmuskel der Harnblase\footnote{zum Beispiel im Rahmen eines
Querschnitt-Syndroms nach Wirbelsäulenverletzungen} zu
Harnverhalten führen. 
}
{
\begin{itemize}
     \item Vergrößerung der \gEs{Prostata}
    \item Lähmung (Querschnittlähmung)
    \item Schließmuskel-Krampf
    \item Verlegung der Harnröhre
    \item Medikamente
\end{itemize}
}{}{}{}

Je nach Füllungszustand der Harnblase kann das
Harnverhalten sehr schmerzhaft sein.

Die Maßnahmen beschränken sich auf einen zügigen Transport
an eine Abteilung für Urologie. Eine Notarzt-Nachforderung
ist kaum sinnvoll, da ein \gAbk{NEF} oder \gAbk{NAW} selten geeignetes
Material zur Harnableitung mitführt. 

\gMassnahmenNG{Spezielle Maßnahmen}{Akutes Harnverhalten}{}
{
\begin{itemize}
    \item Zügiger Transport in eine Urologische Ambulanz
    \item \gZ{vorsichtig ablassen über Harnkatheter oder
    Blasenpunktion (Spital)}
\end{itemize}
}




\subsection{Niereninsuffizienz und Nierenversagen}
\label{topic:cni-therapie-dialyse}
\label{topic:dialysepatienten-gefahren}

\gAuthNotes{Abkz. CNI}

Meistens kommt es im Rahmen anderer Krankheiten
(Diabetes mellitus, Hypertonie, erhöhte Blutfette \ldots) zu einer massiven
Schädigung der kleinen Nierengefäße und in weiterer Folge zum Nierenversagen. Da
die zuvor erwähnten Krankheiten sehr häufig sind, ist in weiterer Folge die
Niereninsuffizienz ebenfalls häufig anzutreffen. 

Als Therapie steht der
Medizin anfänglich eine medikamentöse Behandlung, später jedoch nur die
Nierentransplantation oder die \gE{Blutwäsche}
(\gT{Dialyse}\index{Dialyse}) zur Verfügung. Bei der
klassischen Dialyse wird der Patient jeden zweiten oder
dritten Tag für einige Stunden an eine Maschine angeschlossen, welche sein
Blut reinigen kann. Die "`Dialysefahrten"' sind eine der
Hauptaufgaben des Krankentransportdienstes.

\index{Dialyse-Patient!Gefahren}
\gEs{Der Ausfall der normalen Nierenfunktion hat
Auswirkungen} auf den Wasserhaushalt, den
Elektrolythaushalt und den Säurebasenhaushalt. Diese Funktionen der Niere sind zum Überleben elementar (Vitalfunktionen zweiter
Ordnung). Dialysepatienten müssen daher sehr darauf achten
nicht zu viel zu trinken. Weiters können \gE{Elektrolytstörungen} sogar
\gE{lebensbedrohliche Notfälle} nach sich ziehen (bis hin
zum Herzstillstand).

Es ist daher wichtig auf den
Dialysepatienten -- auch im "`normalem"' Krankentransport --
immer ein wachsames Auge zu haben.







% \subsection{Hodentorsion}
% 
% Verdrehung eine Hodens im Hodensack  ={\textgreater}
% 
% Abschnürung von Samenstrang und  Blutgefäßen
% Oft nach Eintritt in Pubertät, ev. {\textgreater} 20 J.
% 
% \paragraph{Symptome:}
% 
% {\mdseries
% akuter heftiger Schmerz}
% 
% {\mdseries
% Schwellung}
% 
% \paragraph{Therapie:}
% 
% hochlagern
% 
% umgehende urologisch-chirurgische  Intervention (ad URO!!!)




\clearpage
\chapter{[MED-GYN] \emph{Spezielle Patientengruppe:} Frauen und Schwangere}
\minitoc


\section{Schwangerschaft und Geburtshilfe}

\subsection{Weiblicher Zyklus}

Der Begriff Zyklus oder Periode beschreibt in diesem Zusammenhang den Vorgang
des rund alle \gEs{28 Tage} stattfindenden Reifens
mindestens einer Eizelle, dem Eisprung und den damit im Zusammenhang stehenden Veränderungen der
Gebärmutterschleimhaut. Der Zyklus ist im Kapitel Anatomie und Physiologie
unter \ref{topic:weiblicher-zyklus} beschrieben.








\subsection{Schwangerschaft -- Gestation}

\paragraph{Schwangerschaftsdauer} Die normale
Schwangerschaft dauert \gEs{40 Wochen $\pm$  2 Wochen}. Das sind rund
\gEs{neun Kalendermonate} oder 10 Mondmonate
(Lunarmonate\footnote{\gZ{lat. luna: Mond.}}). Der
Fortschritt der Schwangerschaft wird in
\gT{Schwangerschaftswochen}, Abk. \gT{SSW}, angegeben. Von
einer \gT{Frühgeburt} spricht man wenn die Geburt
\gE{zwischen der 26. und bis zur vollendeten 37. SSW}
erfolgt. Von einer \gE{Übertragung} spricht man ab der
\gE{42.~SSW}\gZ{, die Geburt wird dann künstlich
eingeleitet}.

\gLayout{22}{Schwangerschaftsverlauf}
{
Nach dem Eisprung ist die Eizelle 8-10\,Std.
befruchtungsfähig. Sie benötigt für die Wanderung Eierstock -- Eileiter --
Uterus (Gebärmutter) 4-6 Tage. Wurde die Eizelle befruchtet, kommt es zur
\gEs{Einnistung in der Gebärmutterschleimhaut}. Dort bildet
sich der \gT{Mutterkuchen} (\gT{Plazenta}), welcher die
Frucht\footnote{Mit dem Begriff \emph{"`Frucht"'} ist das
Kind gemeint.} mit Nährstoffen und \ce{O2} versorgt. Weiters
bildet sich um die Frucht eine \gT{Fruchtblase}, welche mit \gT{Fruchtwasser} gefüllt ist und sie so ideal vor äußeren
Erschütterungen schützt.

Die 'Frucht' hat - je nach Schwangerschaftsdauer - verschiedene Namen:

% \gEs{Ebryo} Frucht vom Tage der Befruchtung bis zur 12 SSW \gEs{Fetus} Frucht ab
% der 12 SSW bis zur Geburt

\begin{description}
    \item[\gT{Embryo}] Frucht bis zur 12. SSW
    \item[\gT{Fetus}] Frucht ab der 12. SSW
\end{description}

Zwischen der 18. und der 20. SSW kann die Mutter die erste
Kindesbewegungen spüren. \gEs{Ab der 24.
SSW gilt das Kind als überlebensfähig} und ggfs. bei einer
Frühgeburt als reanimationspflichtig.
}{}
{./images/wikipedia-pd/Nabelschnur.png}
{Die Lage des Fetus in der Gebärmutter.}
{Gray's Anatomy, 20th ed., PD, Copyright expired.}




%%%%%%%%%%%%%%% GAB
% 
% Nach einer Befruchtung
% wandert innerhalb von 4-6 Tagen die Eizelle normalerweise weiter in den Uterus
% und nistet sich dort ein. 
% 
% Es bildet sich dann der Mutterkuchen, die \gT{Placenta} bildet sich. Sie dient
% der Nährstoff- und Sauerstoffversorgung des Fetus. 
% 
% Hormon: ${\beta}$-HCG (in Blut und Urin der Mutter nachweisbar)
% 
% \paragraph{Schwangerschaft -- Verlauf}
% 
% Die normale Schwangerschaft
% dauert im Mittel \gEs{40 Wochen} $\pm$ 2 Wochen, das sind
% rund Kalendermonate oder 10 Lunarmonate\footnote{\gZ{lat. luna: Mond.} Eine
% Mondphase hat 28 Tage.}. Der Fortschritt der Schwangerschaft wird in 
% \gEs{Schwangerschaftswochen}, Abk. \gEs{SSW}, angegeben. 
% 
% Die \gE{'Frucht'} hat -- je nach Schwangerschaftsdauer -- verschiedene Namen:
% 
% \begin{description}
%     \item[\gT{Embryo}] Frucht bis zur 12. SSW
%     \item[\gT{Fetus}] Frucht bis zur 12. SSW
% \end{description}
% 
% Erste Kindesbewegungen können von der Mutter in der 18.-20. SSW gespürt werden.
% Die Lungenreife wird vom Kind in der 26. SSW erreicht. 
% 
% Ab der \gEs{SSW ~24} gilt das Kind als
% überlebensfähig (und ggfs. reanimationsplflichtig).
% 
% \subparagraph*{Normale Schwangerschaftsdauer} Die normale Schwangerschaft
% dauert \gEs{38-42 Wochen}, das sind rund Kalendermonate oder 10
% Lunarmonate\footnote{\gZ{lat. luna: Mond. Eine Mondphase hat 28 Tage.}} Der
% Fortschritt der Schwangerschaft wird in \gE{Schwangerschaftswochen}, Abk. SSW,
% angegeben.
% 
% Die Schwangerschaft wird in 3 Drittel unterteilt, welche
% \gT{Trimenon}\index{Trimenon} heißen.
% 
% Von einer \gEs{Frühgeburt} spricht man wenn die Geburt bis zur vollendeten 37.
% SSW erfolgt.



\subsection{Notfälle -- Frühschwangerschaft}

\subsubsection{Fehlgeburt (Abortus)}

\gDefinition{Beendigung einer Schwangerschaft bei einem
Kindesgewicht von unter 500\,g und fehlenden Lebenszeichen.
\cite{HaagGyn2007}}


Aufgrund verschiedener Ursachen kann es zur Fehlgeburt (Ausstoßung der Frucht)
kommen. In den meisten Fällen kommt es zum Absterben der Frucht, welche nach
einigen Tagen ausgestoßen wird. Ein sehr früher
Fruchtabgang wird oft nicht einmal bemerkt bzw. mit
Regelunregelmäßigkeiten verwechselt. 

\gEs{Symptome:}
\begin{itemize}
	\item Vaginale Blutung
	\item Abgang von Gewebsteilen
	\item Unterbauchschmerzen
\end{itemize}

\gMassnahmenNG{Spezielle Maßnahmen}{Abort}{}{
\begin{itemize}
    \item Lagerung: sterile Vorlage und Lagerung nach
    Fritsch 
    % TODO Querverweis zu Lagerungsarten einfügen
    \item \gTxBeiVit: \gTxMassVitBes
    \begin{itemize}
        \item Lagerung: Situationsgerecht
    \end{itemize}
    \item Hospitalisation: Abt. f. Gynäkologie und
    Geburtshilfe.
\end{itemize}
} 

\gCave{Achtung! Ab der 24. SSW ist die theoretische
Lebensfähigkeit erreicht, das Kind muss reanimiert werden!}



\subsubsection{Eileiterschwangerschaft -- Tubaria}

\emph{Syn.: extrauterine Schwangerschaft}

Wenn sich die befruchtete Eizelle nicht in der Gebärmutter, sondern schon im
Ei\-leiter einnistet, kommt es zu einer
\index{Eileiter\-schwanger\-schaft}\gEs{Eileiter\-schwanger\-schaft}
die lebens\-bedrohlich ist. 

\gMarried{
Mit zunehmender Größe der Frucht wird der Eileiter
gedehnt, dabei kommt es zu Schmerzen. Schließlich reißt der Eileiter 
und die Frau erleidet massive innere Blutungen. Es zeigen sich dann Anzeichen
eines akuten Abdomens und eines hypovolämischen Schocks. 
}{
\begin{itemize}
    \item Einnistung der Frucht in Eileiter (statt in Gebärmutter)
    \item Schmerz durch Dehnung des Eileiters
    \item Schließlich Zerreißung des Eileiters
    \begin{itemize}
        \item Blutung
        \item Akutes Abdomen
        \item Schock
    \end{itemize}
\end{itemize}
}




${\rightarrow}$ Bei Frauen mit akutem Abdomen immer nach möglicher
Schwangerschaft fragen!

\gMassnahmenNG{Spezielle Maßnahmen}{Verdacht auf Eileiterschwangerschaft}{}{
\begin{itemize}
    \item Hospitalisation: Abt. f. Gynäkologie
\end{itemize}
}

\gMassnahmenNG{Spezielle Maßnahmen}{Rupturierte Eileiterschwangerschaft}{}{

\begin{itemize}
    \item \gTxMassVitBes
    \begin{itemize}
        \item Lagerung: Bauchdeckenentspannend.
        \item Schockbekämpfung
        \item \ce{O2}: 10-15\gLMin
    \end{itemize}
\end{itemize}
}

\subsection{Notfälle -- Spätschwangerschaft}

\subsubsection{Vorzeitige Plazentalösung}

\gDefinition{Ablösung des Mutterkuchens (Plazenta) vor Beendigung der Geburt}

Aufgrund einer Einblutung zwischen Plazenta und Gebärmutter kann es zu
einer vorzeitigen Ablösung der Plazenta mit dann ungehinderter
Blutung kommen. Es besteht \gEs{Lebensgefahr} für Mutter und Kind durch
Blutung der Mutter und Mangelversorgung des Kindes. 

Die Diagnose ist durch einen Sanitäter nicht sicher zu stellen. Im
Vordergrund steht daher die Behandlung der Symptome (Schock) und ein
zügiger Transport. 

\paragraph{\gTxSymptome}

\begin{itemize}
    \item unspezifisch:
    
    \begin{itemize}
        \item Schock
        \item akutes Abdomen
    \end{itemize}
\end{itemize}

\gMassnahmenNG{Spezielle Maßnahmen}{Vorzeitige Plazentalösung}{}{
    \begin{itemize}
        \item \gTxMassVitBes
        \begin{itemize}
            \item Lagerung: sterile Vorlage und Lagerung nach Fritsch.
            \item Schockbekämpfung
        \end{itemize}
        \item Keine vaginale Untersuchung oder  Scheidentamponade
        \item Zügiger Transport mit Voranmeldung
    \end{itemize}
}


\subsubsection{Vena-cava-Syndrom}
\label{topic:vena-cava-syndrom}

Es kann vorkommen, dass das \gEs{Kind die untere
Hohlvene der Mutter abdrückt}. Dadurch wird der Blutstrom zum Herzen
vermindert und es kann zu einem \gEs{plötzlichen
Kollaps der Mutter} und einer Unterversorgung des Kindes kommen. Die
Therapie besteht in einer Linksseitenlagerung, da die Hohlvene eher
rechts der Körpermitte verläuft. 

\gFazit{\gEs{Jede hochschwangere Patientin }sollte daher in
Linksseitenlage und nicht in Rückenlage transportiert werden. }

\gMassnahmenNG{Spezielle Maßnahmen}{Vena-cava-Syndrom}{}{

\begin{itemize}
    \item Linksseitenlage, grundsätzlich immer bei
    hochschwangeren Patientinnen.
\end{itemize}
}

\subsubsection{Präeklampsie, EPH-Gestose, HELLP-Syndrom, Eklampsie}
\label{topic:eph-gestose}\label{topic:praeeklampsie}\label{topic:eklampsie}\index{EPH-Gestose}\index{Präeklampsie}\index{Eklampsie}\index{Krampfanfall!eklamptischer}

Synonym: \gT{Präeklampsie}
% FEATURE INHALT: HELLP,Syndrom, Leberspannungsschmerz

Die \gT{EPH-Gestose}, auch \gT{Präeklampsie} genannt,
ist eine der gefürchtetsten Komplikationen der
Spätschwangerschaft. Unbehandelt kann sie zu einem \gE{eklamptischen
Krampfanfall} (\gT{Eklampsie}, ähnlich einem "`normalen"'
zerebralen Krampfanfall) und in weiterer Folge zum Tod von Mutter und Kind
führen.

\gFazit{Die \emph{Prä}eklampsie (EPH-Gestose) kann zu einem
lebensgefährlichen Krampfanfall (Eklampsie) führen.}

\paragraph{\gTxSymptome}

\begin{enumerate}
    \item \gEs{E}dema (Ödeme)
    \item \gEs{P}roteinuria (Proteinurie = Eiweiß im Harn)
    \item \gEs{H}ypertonie (RR  ${\uparrow}$)
\end{enumerate}


RR {\textgreater}140/90 bei einer Schwangeren ist zu  hoch!!!

Mutter-Kind-Pass Untersuchungen zur  Früherkennung



\gMassnahmenNG{Spezielle Maßnahmen}{Verdacht auf
Präeklampsie}{}{
\begin{itemize}
    \item schonender Transport
    \item Überwachung, evt. \ce{O2}
    \item Reizabschirmung
    %\item Notarzt: Krampfdurchbrechung, RR-Senkung
    %\item (im Vollbild: vorzeitige Beendigung der  Schwangerschaft)
\end{itemize}
}

\gMassnahmenNG{Spezielle Maßnahmen}{Eklampsie -- Eklamptischer Krampfanfall}{}{
Grundsätzlich ist so wie bei jedem anderen
Krampfanfall zu verfahren , vgl. Kap. \ref{topic:krampfanfall}. Bedenke:
Der eklamptische Krampfanfall ist für Mutter und Kind lebensgefährlich!
}


\subsubsection{Vorzeitiger Fruchtwasserabgang}

Der Fetus schwimmt während  Schwangerschaft in steriler
Flüssigkeit (Fruchtwasser). Bei vorzeitiger Eröffnung der
Fruchtblase ergeben sich für ihn die Gefahr
eines \gEs{Nabelschnurvorfalls} oder einer \gEs{Infektion}
(Keimbesiedlung der Fruchthöhle). Die Patientin ist an
einer geburtshilflichen Abteilung vorzustellen.

\gMassnahmenNG{Spezielle Maßnahmen}{Vorzeitiger Fruchtwasserabgang}{}{
\begin{itemize}
    \item Lagerung: Linksseitenlagerung, Becken erhöht
    (Schwerkraft verzögert Geburt), Lagerung nach Fritsch
    \item Nicht mehr aufstehen lassen! Liegender Transport!
    \item Hospitalisation: Abteilung für Gynäkologie und
    Geburtshilfe, Kreißsaal
\end{itemize}
}



\subsection{Die Geburt}
\label{topic:geburt}

Als \gEs{Beginn der Geburt} zählt:
\begin{itemize}
\item wenn regelmäßige Wehen in 10-15 Minuten Abstand auftreten und dieser
Zustand länger als eine Stunde dauert \item wenn die Fruchtblase springt und
Fruchtwasser abgeht
\end{itemize}

Die Geburt verläuft in 3 Phasen. Der \gEs{Eröffnungsphase}
mit den Eröffnungswehen und dem Blasensprung. Der
\gEs{Austreibungsphase} mit den Presswehen und der Geburt
des Kindes. Und der \gEs{Nachgeburtsphase} mit der Placenta als
Nachgeburt.


\gLayout{0}{Eröffnungsphase}
{
Durch rhythmisches Zusammenziehen der Gebärmuttermuskulatur (Wehen) wird das Kind
langsam in den Geburtskanal nach unten gepresst. Dadurch verkürzt sich der
Gebärmutterhals (Cervix), und wird durch den Druck des Kindskopfes zu einer
flachen, runden Öffnung (Muttermund). 

Spätestens am Ende der Eröffnungsphase
kommt es zum Platzen der Fruchtblase (\gT{Blasensprung}),
welches den Abgang von Fruchtwasser und eventuell die Absonderung von blutigen Sekret mit sich führt. 

\gCa{Ab
dem Blasensprung darf die Gebärende nicht mehr aufstehen,
sondern muss liegend transportiert werden. Es besteht die Gefahr eines Nabelschnurvorfalls in den
Geburtskanal!}
}{
\begin{itemize}
     \item Wehen im Abstand von 10-15 Minuten
     \item Gebärmutterhals verkürzt sich, Muttermund öffnet sich
     \item Blasensprung mit Fruchtwasserabgang
     \item Ab Blasensprung: \gEs{Gefahr eines
     Nabelschnurvorfalles}
     \begin{itemize}
         \item \gCa{Liegender Transport!}
     \end{itemize}
 \end{itemize}
}{}{}{}

\gLayout{0}{Austreibungsphase}
{
Nachdem der Muttermund vollständig eröffnet ist (10\,cm),
beginnt die Austreibungsphase. Die Wehen kommen nun alle 2-3 Minuten und sind nun keine
Senkwehen mehr, sondern richtige \gEs{Presswehen}. Das
Hinterhaupt des Kindes wird nun zwischen den Schamlippen  sichtbar. 
}{
\begin{itemize}
    \item Wehen im Abstand von 2-3 Minuten
    \item Presswehen
    \item Geburt des Kindes
\end{itemize}
}{}{}{}


Mittels \gT{Dammschutz} kann die
Wahrscheinlichkeit eines Dammriss minimiert werden. Der \gT{Damm} ist das Stück
Haut zwischen Scheide und Anus. Gleichzeitig wird das Kind vor dem Stuhl der
Mutter geschützt\footnote{Durch das Pressen wird nicht nur
das Kind geboren, sondern in der Regel auch Stuhl abgesetzt.}. 

Wenn der Kopf geboren ist, folgt der restliche Körper
in der Regel schnell nach. (Vorsicht: Das Kind ist durch die Fruchtschmiere
rutschig!). 

\gMassnahmenNG{Spezielle Maßnahmen}{Versorgung des Neugeborenen}{}{
\begin{enumerate}
    \item Nun wird das Neugeborenen mittels Schleimabsauger
    abgesaugt und kann für das Abnabeln vorbereitet werden.
    \item Dabei wird eine Handbreite vom Kind entfernt die
    erste Nabelklemme gesetzt.
    \item Dann wird die Nabelschnur Richtung Mutter ausgestreift und
    wieder eine Handbreit entfernt die zweite Klemme gesetzt.
    \item Zwischen den beiden
    Klemmen wird die Nabelschnur durchtrennt.
    \item Zu guter letzt wird das Neugeborene
    warm eingepackt (nur das Gesicht darf frei bleiben) und der Mutter überreicht.
\end{enumerate}
}


\gLayout{0}{Nachgeburtsphase}
{
Nach 10-20 Minuten kommt es neuerlich zu Wehen, welche den Presswehen ähnlich
sind. Es kommt zur Geburt der \gE{Placenta}
(\gE{Mutterkuchen}). Die Gefahr hierbei ist, dass es zu
starkem Blutverlust bei der Mutter kommen kann. \gEs{Die
Placenta ist zur Kontrolle der Vollständigkeit ins
Krankenhaus mitzunehmen!} 
}{
\begin{itemize}  
    \item nach 10-20 min. neuerlich Wehen 
    \item Geburt der Placenta
\end{itemize}
}{}{}{}





% % %%%%%%%%%%% GAB
% 
% \paragraph{Verlauf}~
% 
% Die Geburt verläuft in \gE{3} Phasen:
% 
% %\begin{gWideParEnv}
% \begin{center}
% \begin{tabulary}{\linewidth}{LLL}
% \toprule
% Phase
%  &
% Beschreibung
%  &
% Wehen
% \\\midrule
% \gT{Eröffnungs\-phase}
%  &
% Eröffnungs\-wehen
% 
% Gebärmutterhals (Cervix) verkürzt sich,  Muttermund
% öffnet sich
% 
% Blasensprung
% 
% \gE{Liegender Transport}, sonst Beschleunigung der  Geburt
% und Gefahr eines  \gEs{Nabelschnurvorfalls}!! 
% 
% Kindskopf tritt in Geburtskanal ein
%  &
% 10-15 min.
% \\\hline
% \gT{Aus\-treibungs\-phase}
%  &
% Preßwehen
% 
% Pausen wichtig (Erholungsphase für Fetus)
% 
% Muttermund vollständig eröffnet
% 
% \gT{Preßwehen}: Kontraktion Bauch- plus 
% Uterusmuskulatur
% 
% Kopf wird sichtbar -- normalerweise  Hinterhauptslage (alle
% anderen Lagen stellen  Komplikationen dar)
% 
% \gT{Dammschutz} ${\rightarrow}$ sonst
% Dammriß, Schutz des Kindes vor Stuhl der Mutter
% 
% Kopf geboren ${\rightarrow}$ Körper folgt schnell nach
% ${\rightarrow}$  auffangen!
% 
% Absaugen ${\rightarrow}$ abnabeln ${\rightarrow}$ warm
% einpacken ${\rightarrow}$ der Mutter auf die Brust legen
% &
% 2-3 min.
% \\\hline
% \gT{Nach\-geburts\-phase}
%  &
% Nach 10-20 min. neuerliche Wehen
% 
% Geburt der Placenta  (Nachgeburt)
%  &
% \\\bottomrule
% \end{tabulary}
% \end{center}
% %\end{gWideParEnv}
% 
% \subparagraph{Neugeborenes - Vitalitätsbeurteilung}
% 
% \subparagraph{\gZ{APGAR-Score}}
% 
% \gZ{
% \begin{enumerate}
%     \item A ussehen
%     \item P uls
%     \item G esichtsbewegung
%     \item A ktivität
%     \item R espiration
% \end{enumerate}
% 
% je 0-2 Pkt., jeweils nach 1\,min\,/\,5\,min\,/\,10\,min}









\subsubsection{Die Geburt im Rettungsdienst}

\gEs{Fragen des Sanitäters an die Gebärende:}

\begin{itemize}
    \item \emph{"`In der wievielten SSW befinden Sie
    sich? Wann ist der geplante Geburtstermin?"'} 
	\item \emph{"`Die wievielte Geburt ist das?"'} Bei
	Mehrgebärenden ist der Geburtsverlauf oft rascher.
	\item \emph{"`In welchem Abstand sind die Wehen?"'} Je
	nachdem ist erkennbar in welcher Geburtsphase sich die Gebärende befindet
	\item \emph{"`Gab es schon einen Blasensprung?"'} Wenn ja
	darf die Gebärende nur mehr liegend transportiert werden! Gefahr: Nabelschnurvorfall!
	\item \emph{"`Gab es während der Schwangerschaft
	Komplikationen?"'}
	\item Mutterkindpass verlangen. Darin stehen mit
	unter wichtige Informationen über die Schwangerschaft bzw. die  Geburt.
	\item \emph{"`In welchem Krankenhaus ist die Geburt
	angemeldet?"'} Sofern nicht im Mutterkindpass vermerkt.
\end{itemize}

\paragraph{Der Transport}~

\begin{wrapfigure}{r}{5cm}
\includegraphics[width=\linewidth]{./images/med/birth-teaser-welcome-0500.jpg}
\end{wrapfigure}
Befindet sich die Mutter noch in der Eröffnungsphase wird ein schnellstmöglicher,
schonender Transport ins Krankenhaus angestrebt. Grundsätzlich wird die Gebärende
so transportiert wie es für Sie am angenehmsten ist. Ab dem
Blasensprung darf sie jedoch nur mehr liegend transportiert werden. Sollten Sie die Mutter stehend
angetroffen haben, wird am Transportschein dokumentiert: "`stehend angetroffen;
liegend transportiert". 

Um die Gefahr eines \gE{Vena-cava-Kompressionssyndroms}
(\gSee{topic:vena-cava-syndrom}) zu vermeiden, sollte die
Patientin liegend in leichter \gEs{Linkseitenlage}
transportiert werden. Befindet sich die Mutter schon in der
Austreibungsphase und der Kindskopf ist zwischen den
Schamlippen sichtbar, ist der weitere Transport zu
unterlassen und die Sanitäter müssen bei der Geburt
Hilfestellung leisten und einen Notarzt nachfordern.

Am
Ende der Geburt, müssen die \gEs{Geburtszeit} und der
\gEs{Geburtsort} (Straße mit Hausnummer; Kreuzung; Autobahn mit km-Angabe, etc.) \gEs{dokumentiert} werden.
Der Geburtsort ist insofern wichtig, da für die Ausstellung der Geburtsurkunde  das für den
Geburtsort zuständige Bezirksamt verantwortlich ist. Wenn Mutter und Kind wohlauf
sind, kann der Transport in Krankenhaus fortgesetzt werden. 

Es soll nicht auf
die Nachgeburt gewartet werden. Sollte die Plazenta vor Eintreffen im
Krankenhaus geboren werden, ist diese unbedingt sicherzustellen und mitzunehmen.
Sie wird im Krankenhaus auf Vollständigkeit und Reife untersucht.







\subsection{Notfälle während der Geburt}

\subsubsection{Nabelschnurvorfall}

Bei einem Nabelschnurvorfall in den Geburtskanal bestehen zwei Gefahren:

Einerseits wird während des Geburtsvorganges die Nabelschnur durch den
Kindskopf abgeklemmt, und somit ist die Sauerstoffversorgung des Kindes nicht mehr
gewährleistet. Andererseits kann sich die Nabelschnur  um den Hals
des Kindes wickeln und es kommt während der Geburt zur Strangulation des
Kindes.

\gMassnahmenNG{Spezielle Maßnahmen}{Nabelschnurvorfall}{}{
\begin{itemize}
	\item Beckenhochlagerung
	\item Mutter soll \underline{nicht} pressen (wenn möglich Geburt erst im
	Krankenhaus)
	\item Wehenhemmung (hecheln; \gZ{medikamentös durch Notarzt})
	\item Nabelschnurumschlingung des kindlichen Halses während der Geburt sofort lösen
\end{itemize}
}

\subsubsection{Placenta praevia}

Normalerweise liegt die Placenta nicht vor dem Geburtskanal. Bei einer sehr
tiefen Einnistung des befruchteten Eis in der Gebärmutter, kann es jedoch
dazukommen, dass sich die Placenta \emph{vor} dem Geburtskanal bildet. Dies ist
natürlich ein sehr großes Geburtshindernis, außerdem kann
es durch vorzeitiges Lösen des Mutterkuchens zu massiven Blutungen kommen.

Eine Plazenta praevia ist meistens schon durch die Kontrolluntersuchungen beim
Frauenarzt bekannt und kann erfragt bzw. im Mutter-Kind-Pass nachgelesen werden.

\gEs{Symptome:}
unter Wehentätigkeit starke Blutungen aus der Scheide

\gMassnahmenNG{Spezielle Maßnahmen}{Placenta praevia und bevorstehende
Geburt}{}{
\begin{itemize}
    \item \gTxMassVitBes
    \begin{itemize}
        \item Lagerung: sterile Vorlage und Lagerung nach Fritsch
        \item ggfs. Schockbekämpfung
    \end{itemize}
\end{itemize}
}

% \paragraph{Gefahren:}
% 
% Abklemmung der Nabelschnur  (Sauerstoffversorgung  [F0EA?] )
% 
% Strangulation (Gehirnperfusion [F0EA?] )
% 
% \paragraph{Therapie:}
% 
% Linksseitenlagerung, Beckenhochlagerung
% 
% Mutter soll NICHT pressen (Geburt w.m. erst  im KH!!!)
% 
% Wehenhemmung (hecheln, medikamentös) 
% 
% Nabelschnurumschlingung des  kindlichen Halses während der Geburt 
% sofort lösen!
% 
% Placenta praevia
% 
% Plazenta liegt unterhalb des Fetus in  Geburtskanal
% 
% Geburtshindernis
% 
% massive Blutungen
% 
% \paragraph{Maßnahmen}
% 
% Lagerung nach Fritsch
% 
% Schockbekämpfung
% 
% NAW
% 
% im KH: Notsectio

\subsubsection{Pathologische Geburtslagen}


Alle Lagen des Kindes, bei denen nicht zuerst der Kopf im Geburtskanal erscheint.
\begin{itemize}
	\item Steißlage bzw. Beckenendlage
	\item Armvorfall
\end{itemize}

Die größte Gefahr bei allen Lageanomalien ist, dass die Nabelschnur abgedrückt
wird. Daher sollte unter allen Umständen die Geburt im KH, mittels Kaiserschnitt,
durchgeführt werden. Die Geburt in jedem Fall verhindern bzw. verzögern.


% 
% alle Lagen, bei denen nicht zuerst der  Kopf im Geburtskanal erscheint 
% 
% Steißlage
% 
% Beckenendlage
% 
% Armvorfall
% 
% \paragraph{Gefahr}
% 
% Abdrücken der Nabelschnur
% 
% Geburtsstopp
% 
% Geburt nur im KH durchführbar (Sectio- Bereitschaft!)
% 
% In jedem Fall Geburt  verhindern/hinauszögern!

\subsubsection{Uterusatonie}


Normalerweise zieht sich in der Nachgeburtsphase die Uterusmuskulatur zusammen. In diesem Fall kommt es aber zur Erschlaffung (Atonie) der Gebärmutter. Die Folge sind lebensbedrohliche Blutungen durch das fehlende Kontrahieren des Uterus nach der Plazentalösung.

\gMassnahmenNG{Spezielle Maßnahmen}{Uterusatonie}{}{
    \begin{itemize}
        \item \gTxMassVitBes
        \begin{itemize}
            \item Schockbekämpfung
            \item Lagerung: Fritsch, Beine hoch
        \end{itemize}
    \end{itemize}
}

% normalerweise zieht sich in der  Nachgeburtsphase die Uterusmuskulatur 
% zusammen
% 
% hier: Erschlaffung der Uterusmuskulatur  nach Plazentalösung
% ={\textgreater} großflächige  Blutung!
% 
% \paragraph{Maßnahmen}
% 
% \begin{itemize}
%     \item Allgemeine Maßnahmen bei lebensbedrohlichen Zuständen
%     \item Schockbekämpfung
%     \item Lagerung: Fritsch, Beine hoch
%     \item Evt. Druck mit beiden Fäusten oberhalb  Schambein ${\rightarrow}$
%     Blutungsverminderung
%     \item Zügiger Transport mit Voranmeldung
% \end{itemize}


\subsubsection{Uterusruptur}

Die Uterusruptur kommt meistens nur nach vorangegangenem
Kaiserschnitt vor. Es treten plötzlich Schmerzen während
der Wehen auf. Die Wehen nehmen dann ab oder hören
plötzlich auf. Damit wird der Geburtsvorgang gestoppt und es herrscht Lebensgefahr für Mutter, wegen lebensbedrohlichen Blutungen, und Kind, wegen \ce{O2}-Mangel.

\gMassnahmenNG{Spezielle Maßnahmen}{Uterusruptur}{}{
\begin{itemize}
	\item \gTxMassVit
	\item Schockbekämpfung
	\item Lagerung: Fritsch, Beine hoch
\end{itemize}
}

% 
% {\mdseries
% plötzlicher Schmerz während Wehen,  dann keine Wehen mehr
% ${\rightarrow}$ Geburtsstop}
% 
% \gFazit{Lebensgefahr für Mutter (Blutung) und  Kind (O2
% -Mangel!)}
% 
% \paragraph{Therapie}
% 
% \begin{itemize}
%     \item Allgemeine Maßnahmen bei lebensbedrohlichen Zuständen
%     \item Schockbekämpfung
%     \item Zügiger Transport mit Voranmeldung
% \end{itemize}

\subsubsection{Asphyxie}

Neugeborenes atmet nicht oder nicht  ausreichend, siehe APGAR - Score

\gMassnahmenNG{Spezielle Maßnahmen}{Asphyxie des Neugebornen}{}{
    \begin{itemize}
        \item absaugen
        \item Fußsohlen und Rücken klopfen, reiben
        \item \ce{O2} - Gabe
        \item Reanimation, Notarzt
    \end{itemize}
}

\section{Sonstige gynäkologische Erkrankungen und
Notfälle}

\subsection{Vaginale Blutungen}

\gMarried
{Die vaginale Blutung ist präklinisch schwer abzuklären und
kann viele Ursachen haben. Wichtig ist festzuhalten, dass
jede vaginale Blutung abgeklärt gehört, da sich dahinter
auch ein Tumorgeschehen (besonders bei Patientinnen in der
Menopause) verstecken kann. Darüber hinaus kommt eine
verstärkte Monatsblutung, Verletzung (nach Sturz, 
Geschlechtsverkehr oder nach Vergewaltigungen),
Defloration\footnote{Entjungferung} oder Blutungen im Rahmen
einer Schwangerschaft in Betracht. Bei vaginalen Blutungen in Kindern darf man \gE{nicht}
automatisch von einem erfolgten Missbrauch ausgehen. }
{
\begin{itemize}
    \item Verstärkte Monatsblutung
    \item Verletzungen (Sturz, Geschlechtsverkehr, Vergewaltigung,
    \ldots)
    \item Defloration
    \item Tumor
    \item Im Rahmen einer Schwangerschaft
\end{itemize}
}

Eine akute vitale Bedrohung besteht seltenst, dennoch muss
sie routinemäßig vor allem anhand des Blutverlustes und der
Vitalwerte eingeschätzt werden. 


\gMassnahmenNG{Spezielle Maßnahmen}{Vaginaler Blutung}{}{
\begin{itemize}
    \item Lagerung nach Fritsch
    \item Hospitalisation: Abteilung für
    Gynäkologie, bei Wehen Kreißsaal
\end{itemize}
}





\chapter{[MED-KIN] \emph{Spezielle Patientengruppe:} Kinder}
\minitoc
%%%%%%%%%%%%%%%%%%%%%%%%%%%%%%%%%%%%%%%%%%%%%%%%%%%%%%%%%%%
%%% Pädiatrie                                           %%%
%%%%%%%%%%%%%%%%%%%%%%%%%%%%%%%%%%%%%%%%%%%%%%%%%%%%%%%%%%%

\section{Prinzipieller Umgang mit Kindern}

\begin{itemize}
    \item schonend
    \item Aufregung vermeiden
    \item Bezugsperson wichtig
    \item Vertrauensperson bzw. Eltern in die Handlung einbeziehen
    \item Kind während Transport am Besten eng bei Vertrauensperson belassen
\end{itemize}

\paragraph{Altersgruppen bei Kinder}
%  TODO MED: PÄD: Altersgrenzen bei Kindern checken (#19)

\cite{LPN3}

\begin{center}
\begin{tabulary}{12cm}{LL}\toprule
\gEs{Neugeborenes} & bis zum 28. Lebenstag
\\\addlinespace
\gEs{Säugling} & bis zum Ende des 1. Lebensjahr
\\\addlinespace
\gEs{Kleinkind} & 1-5 Jahre
\\\addlinespace
\gEs{Schulkind} & 6-13 Jahre
\\\addlinespace
\gEs{Jugendlicher} & 14-18 Jahre
\\\bottomrule
\end{tabulary}
\end{center}

\gCave{In der Literatur kommt es gelegentlich zu
abweichenden Definitionen. So ist z. B. in den
ERC-Richtlinien zur Reanimation ein Neugeborenes nur bis
unmittelbar nach der Geburt als solches klassifiziert.}

\section{Anatomische und physiologische  Besonderheiten}

Kinder unterscheiden sich wesentlich von erwachsenen Patienten, nicht nur in
ihrer Fähigkeit ihre Umwelt und gesetzte Handlungen zu verstehen, sondern auch
hinsichtlich ihrer Anatomie und Physiologie. 



\gFazit{Kinder sind keine kleinen Erwachsenen!}

\begin{table}[H]
\gCaptionof{table}{Anatomische und physiologische Besonderheiten von Kindern.}

\begin{tabulary}{12cm}{LL}\addlinespace
\gTabSub{Atemfrequenz}
&
Neugeborenes: 40\,\nicefrac{x}{min}

Säugling: 20\,\nicefrac{x}{min}

Kleinkind: 15\,\nicefrac{x}{min}

Beatmung bis zum Säugling immer in  Neutralposition (Kopf nicht
überstrecken)lt. ERC. 

\\\addlinespace
\gTabSub{Zeichen der Atemnot}
&
Einziehungen zwischen Rippen

\index{Nasenflügeln}\gEs{Nasenflügeln}
beschleunigte Atmung,
Atemgeräusche
\\\addlinespace
\gTabSub{Pulstasten beim Säugling}
&
Oberarmarterie (Oberarm Innenseite), Achsel
\\\addlinespace
\gTabSub{Herzfrequenz}
&
Neugeborenes: 140\,--\,160\,/\,min.

Säugling: 120\,/\,min.

Kinder: 80-100\,/\,min.

Bradykardie  ist fast immer  hypoxiebedingt $\rightarrow$
lebensgefährlich!

\\\addlinespace
\gTabSub{Wasser-, Elektrolyt- und Wärmehaushalt}
&
Kinder haben im Verhältnis zu ihrem  Körpergewicht eine
größere Körperoberfläche  als Erwachsene!

höherer Wasseranteil

Austrocknung und Unterkühlung schneller  möglich!
\\\bottomrule
\end{tabulary}
\end{table}

\section{Allgemeines zu Erkrankungen im Kindes- und Jugendalter}

Grundsätzlich können Kinder viele der bereits erwähnten Erkrankungen auch bekommen. 

Bei Kindernotfällen können als zusätzlicher Stressfaktor für
den Rettungsdienst stark emotional betroffene Angehörige sein. Daher gilt es besonders ruhig und gezielt zu handeln. Sie sollten sich der Verantwortung bewusst sein, es geht um die Rettung eines jungen Menschen mit langer Lebenserwartung.

Wichtig für den Transport ins Krankenhaus ist, alle erkrankte
Kinder bis zum vollendeten 18. Lebensjahr auf eine
Kinderambulanz zu transportieren.


\section{Speziell die Kindheit und Jugend betreffende Erkrankungen}




\subsection{Fremdkörperaspiration}


\gERROR{GABS}{????-??-??}{Fremdkörperaspiration: Erwachsen -- Kind überarbeiten
und zusammenführen}
{
\gA{GABS: unklare Aufteilung -- Fremdkörperaspiration Erwachsener -- Kind

Lt. ERC besteht ein Unterschied in der Behandlung nur bei < 1a (keine
Abdomenkompression, stattdessen Thoraxkompression) 

Abschnitt zu respiratorischen Notfällen nehmen und nur Verweis für <1 a auf
Kinderteil?}
}
% TODO Fremdkörperaspiration: Erwachsen -- Kind überarbeiten und zusammenführen
% (# 39)





\gDefinition{Verlegung der Atemwege durch Fremdkörper wie
z.B. Nahrungsmittel, Spielzeug o.ä.}

\gMarginNo{\cite{Erc2005Sec6}} % p 102

\gLayout{2}{\gTxSymptome}
{
% TODO MED: Fließtext fehlt.
\begin{itemize}
    \item Atemnot
    \item plötzlicher starker Husten
    \item Stridor oder Keuchen
    \item Würgen
    \item evt. Zyanose
\end{itemize}
}{

}{}{}{}




\gMassnahmenNG{Spezielle Maßnahmen}{Atemwegsverlegung}{}{

\begin{itemize}
    \item \gTxMassVitBes
    \begin{itemize}
        \item \ce{O2}: 15\gLMin
        \item Lagerung: aufrecht, Oberkörper hoch,
        Atemhilfsmuskulatur unterstützen
    \end{itemize}
    \item Manöver (lt. ERC) unterschiedlich zwischen leichter und schwerer Atemwegsverlegung
    \begin{itemize}
        \item \gEs{leichte Atemwegsverlegung} Kind weint und kann noch auf Fragen verbal Antworten; deutliches Einatmen ist möglich.
        \begin{itemize}
			\item zum Husten auffordern; überwachen; hospitalisieren
		\end{itemize}
		\item \gEs{schwere Atemwegsverlegung} unfähig zu sprechen, atmen oder husten
		\begin{itemize}
			\item 5 Schläge auf den Rücken
			\item 5 Heimlich-Manöver (!beim Säugling jedoch Thoraxkompressionen!)
		\end{itemize}
    \end{itemize}
    \item Bei Atemstillstand Reanimation
\end{itemize}

}

\subsection{Laryngitis, Pseudokrupp}
\label{topic:laryngitis}
\label{topic:pseudokrupp}


\gLayout{2}{\gTxBeschreibung}
{ Bei der Laryngitis (Pseudokrupp\index{Pseudokrupp}\footnote{"`Pseudokrupp"'
bezieht sich auf die Heiserkeit (griech.: pseudo für „unecht“; schottisch: croup,
„Heiserkeit“\cite{wiki:pseudokrupp}). Der echte Krupp bezeichnet die
Kehlkopfentzündung bei Diphtherie.}) handelt es sich um eine Entzündung des Kehlkopfes. Es können auch die oberen Atemwege durch das Anschwellen der Schleimhäute behindert sein (Pseudokrupp/Kruppsyndrom). Diese
Infektion wird durch Viren verursacht. 
}{
% TODO SLIDE fehlt
}{}{}{}


\gLayout{0}{\gTxSymptome}
{
Leitsymptom ist ein pfeifendes Geräusch beim Einatmen
\gE{(inspiratorischer Stridor)} sowie ein trockener, bellender Husten. Es besteht Heiserkeit und
Dyspnoe. Fieber ist eher diskret und steht nicht im Vordergrund. Die Laryngitis
tritt vor allem im Alter zwischen 3 Monaten und 6 Jahren auf.
}{
\begin{itemize}
    \item Pfeifendes Geräusch beim Einatmen
    \gZ{(inspiratorischer Stridor)}.
    \item Trockener, bellender Husten.
    \item Heiserkeit.
    \item Dyspnoe.
    \item Kaum Fieber.
    \item Typisches Alter: 3 Monate -- 6 Jahre.
\end{itemize}
}{}{}{}




\gMassnahmenNG{Spezielle Maßnahmen}{Laryngitis}{}{

\begin{itemize}
    \item Eltern und Kind beruhigen.
    \item Feuchte Luft einatmen lassen (evt. Fenster
    öffnen; feuchte Tücher im Raum aufhängen; Dusche im Badezimmer aufdrehen).
\end{itemize}

}

\subsection{Epiglottitis}
\label{topic:epiglottitis}

\gLayout{2}{\gTxBeschreibung}
{
Bei der Epiglottitis handelt es sich um eine \gE{bakterielle Entzündung} und
daraus folgende \gEs{Schwellung des Kehldeckels}. Sie kommt eher selten vor, kann aber
eventuell lebensbedrohlich werden. Die Epiglottitis ist eine schwere Erkrankung, mit
richtiger  Behandlung besteht aber meist eine gute Prognose.
}{
% TODO SLIDE fehlt
}{}{}{}



\gLayout{0}{\gTxSymptome}
{
Die Epiglottitis geht mit hohem Fieber einher. Es ist ein Atemgeräusch bei der
Einatmung wahrnehmbar (\gE{inspiratorischer Stridor}). Der Patient hält den
Mund offen und es ist ein Speichelfluss bemerkbar. Der
Patient ist eher ruhig und spricht kaum, er hat Angst. 
}{
\begin{itemize}
    \item Hoch fieberhafte Erkrankung.
    \item Inspiratorischer Stridor.
    \item Mund offen, Speichelfluss, spricht kaum!
    \item Angst!
\end{itemize}
}{}{}{}

\gMassnahmenNG{Spezielle Maßnahmen}{Epiglottitis}{}{

\begin{itemize}
    \item \gTxMassVit
    \item Luftbefeuchtung
    \item Kinder sitzen lassen!
    \item Keine Manipulation im Mundraum (Spatel)!
\end{itemize}
}

 \subsubsection{Laryngitis vs. Epiglottitis}
 
Die Unterscheidung beider Krankheiten kann auch für den erfahrenen Sanitäter schwer sein. Darum hier eine Gegenüberstellung der Differenzialdiagnosen Laryngitis (Kruppsymdrom) und Epiglottitis

\begin{table}[H]
\gCaptionof{table}{Gegenüberstellung Laryngitis vs. Epiglottitis}

\begin{tabulary}{12cm}{LL}\addlinespace
\gTabHeading{Laryngitis}
&
\gTabHeading{Epiglottitis}
\\\midrule
Wenig bedrohlich
&
Lebensgefährlich
\\\addlinespace
Kind ist weinerlich
&
Ruhiges Kind, apathisch
\\\addlinespace
Kann schlucken
&
Speichelfluss
\\\addlinespace
Kaum Fieber
&
Hohes Fieber
\\\addlinespace
Kind schreit; weint evt.
&
Sehr leise, verfällt zusehends
\\\bottomrule
\end{tabulary}
\end{table}
 
 
\subsection{SIDS -- Sudden Infant Death Syndrome -- Plötzlicher Kindstod}
\index{plötzlicher Kindstod}
\index{SIDS}
\index{Sudden Infant Death Syndrome}
\label{topic:sids}

\gLayout{0}{\gTxBeschreibung}
{
Der Plötzliche Kindstod  ist eine plötzliche, einschneidende Katastrophe im
Familienleben. Aus \gEs{unbekannter Ursache} -- und aus oft völliger Gesundheit
-- hört das Kind plötzlich auf zu atmen und verstirbt. 

Dieses Schicksal ereilt Kinder im ersten bis vierten Lebensmonat, es ist eine
Häufung in den Monaten Jänner bis März zu beobachten.

}{
\begin{itemize}
    \item Plötzlich, meist im Schlaf.
    \item Vorwiegend 1.-4.Lebensmonat.
    \item Gehäuft Jänner bis März.
    \item Bei Frühgeburten gehäuft.
\end{itemize}
}{}{}{}


\gLayout{0}{Ursachen}
{
Es wird viel über die möglichen Ursachen spekuliert. \gZ{Als mögliche Ursache
wird oft ein unreifes Atemzentrum und eine zentrale Atemregulationsstörung
angeführt}

Letzten Endes ist die genaue \gEs{Ursache unbekannt}, eine
verlässliche Vorhersage, welche Kinder betroffen sein werden,
gibt es nicht. Es wurden lediglich Risikogruppen wie
Frühgeburten und Zwillingsgeschwister von betroffenen
Kindern, identifiziert. Auch scheint die Lagerung in
Bauchlage das Risiko zu erhöhen.

Es muss betont werden, dass eine verlässliche Vorhersage nicht möglich ist und
\gEs{die Eltern keine Schuld trifft}!

}{
\begin{itemize}
    \item Ursache unbekannt.
    \item Risikogruppen.
    \item Zentrale Atemregulationsstörung \gZ{(unreifes 
    Atemzentrum?)}.
    \item Eltern trifft keine Schuld!
\end{itemize}
}{}{}{}


\gMassnahmenNG{Spezielle Maßnahmen}{SIDS}{}{

\begin{itemize}
	\item Reanimation.
	\item Betreuung der Angehörigen. Den Eltern muss das Gefühl
	gegeben werden, dass sie alles für ihr Kind getan haben. Ein solcher
	Unglücksfall tritt schicksalshaft auf und wäre nicht abwendbar gewesen. 
	\item In jedem Fall nimmt die Kriminalpolizei Ermittlungen auf. Darauf müssen
	die Eltern vorbereitet werden. 
	\item Betreuung für die Angehörigen organisieren
	(Kriseninterventionsdienst\opt{REGVIE}{\footnote{In Wien ist die
	\gE{Akutbetreuung Wien} der Stadt Wien zuständig. Sie kann über die
	Leitstelle angefordert werden.}})
	\item Zwillingsgeschwister unbedingt hospitalisieren.
	\item Einsatznachbesprechung und eigene Psychohygiene. Derartige Einsätze
	können -- auch für erfahrenes Fachpersonal -- sehr belastend sein. 
	\opt{ORGASBFLO}{\item \ASBFLO{} Meldung an Gesamteinsatzleiter, Erkundigung
	nach Verfügbarkeit von Peers (auch wenn Peer-Angebot nicht angenommen wird!). }
\end{itemize}

}

\subsection{Ertrinkungsunfall}\index{Ertrinkungsunfall}

\gDefinition{\gEs{Ertrinken}: Tod in den ersten 24\,h nach
Unfall: Ertrinken}

ansonsten:
\index{Beinahe-Ertrinken}\gT{Beinahe-Ertrinken}

Kleinkinder sind besonders gefährdet, da sie meist nicht
schwimmen können und sich oft schnell und unbemerkt der Aufmerksamkeit der Aufsichtsperson
entziehen. Besonders ungesicherte Gartenteiche, Seen, Swimmingpools,
Badeanstalten oder auch Wannenbäder stellen eine Gefahrenquelle im Alltag dar. 

\gMassnahmenNG{Spezielle Maßnahmen}{Ertrinkungsfall}{}{
    \begin{itemize}
        \item Wenn erforderlich Reanimation mit fünf
        Initialbeatmungen. Die Reanimation kann auch nach
        länger andauerndem Kreislaufstillstand erfolgreich sein, da bei einer Unterkühlung der Stoffwechsel verlangsamt wird.
        \item \gTxBeiVit \gTxMassVitBes
        \item Bei Aspiration: Wenn der Verdacht auf eine
        Aspiration besteht, ist der Patient zur Beobachtung zu hospitalisieren.
    \end{itemize}
}


\gFazit{Nobody is dead until warm and dead!}


\subsection{Krampfanfälle im Kindesalter}

Grundsätzlich gelten auch bei Kindern die Ausführungen im Abschnitt
\ref{topic:zerebrale-krampfanfaelle}. Eine Besonderheit stellt jedoch der
\gT{Fieberkrampf} dar.


\subsubsection{Fieberkrampf}
\label{topic:fieberkrampf}

\gLayout{2}{\gTxBeschreibung}
{
Kleinkinder unter 5~Jahren können bei Fieber sogenannte Fieberkrämpfe
entwickeln. Sie entsprechen einem vom Gehirn ausgehenden Krampfanfall und sind grundsätzlich
genau so zu behandeln. 

Mitunter ist die Neigung des Kindes zu Fieberkrämpfen schon bekannt und es
wurden bereits spezielle Zäpfchen für den Bedarfsfall verschrieben (und bei
Eintreffen eventuell schon von den Eltern gegeben). 
}{

}{}{}{}



\gMassnahmenNG{Spezielle Maßnahmen}{Krampfanfälle im Kindesalter}{}{
    
    Die Behandlung erfolgt grundsätzlich wie beim Erwachsenen, siehe
    \ref{topic:zerebrale-krampfanfaelle}.
    
    Die Körpertemperatur ist immer zu messen!
}



\subsection{Kindesmisshandlung}
\label{topic:kindesmisshandlung}

\gDefinition{Absichtliche psychische oder physische Gewalt gegen Kinder}


Kindesmisshandlung kommt in \gEs{allen sozialen Schichten} vor!

\gLayout{2}{Hinweiszeichen}
{
\begin{itemize}
    \item Langes Intervall zwischen  Verletzungszeitpunkt und Verständigung der
    Rettung.
    \item Anamnese und Verletzungsmuster  widersprechen einander!
    \item Mehrere Verletzungen verschiedenen Alters.
    \item Erkennbarer, verletzungsverursachender Gegenstand (Schuhabdruck,
    \ldots).
    \item Zeichen der Vernachlässigung.
    \item Verletzungen bzw. Schmerzen in der Genitalregion (Achtung: Es gibt auch andere
    Ursachen als sexueller Missbrauch!).
    \item "`Untypische"' Verletzungen.
    \item Verhaltensstörungen.
\end{itemize}
}{
\begin{itemize}
    \item Langes Intervall Verletzung -- Hilfeersuchen.
    \item Widerspruch Anamnese -- Verletzungsmuster.
    \item Verletzungen verschiedenen Alters.
    \item Erkennbarer, verletzungsverursachender Gegenstand .
    \item Vernachlässigung.
    \item Verletzungen / Schmerzen in Genitalregion (Achtung: auch andere
    Ursachen!).
    \item "`Untypische"' Verletzungen.
    \item Verhaltensstörungen.
\end{itemize}
}{}{}{}



\gMassnahmenNG{Spezielle Maßnahmen}{Verdacht auf Kindesmisshandlung}{}{
    Die Aufgaben des Rettungsdienstes bei Verdacht auf Kindesmisshandlung
besteht in:

\begin{enumerate}
    \item Daran denken!
    \item Umstände und Umfeld abklären
    \item Unstimmigkeiten und Auffälligkeiten auf eine neutrale Weise
    \gEs{dokumentieren}
    \item Obiges verlässlich bei der Übergabe an die nachbetreuende Einrichtung
    weiterleiten. Evt. Name des Gesprächpartners bei der  Übergabe dokumentieren.
    \item \gEs{Patienten nicht im Stich lassen}: Gelindestes Mittel wählen,
    evt. Vorwand für die Hospitalisierung suchen, aber wenn ein begründeter
    Verdacht auf eine Misshandlung besteht und die
    Erziehungsberechtigten unkooperativ sind notfalls auch die Exekutive beiziehen.
\end{enumerate}

\paragraph{Vorsicht!}

Jedenfalls zu \underline{unterlassen} sind:

\begin{itemize}
    \item Voreilige Verdachtsäußerungen. \gZ{Nicht alles was stinkt ist ein
    Fisch!}
    \item Andeutungen gegenüber den Erziehungsberechtigten oder Dritten. 
\end{itemize}

}





\chapter{[MED-PSY] \emph{Spezielle Patientengruppe:} Psychologie und
Psychiatrie}
\minitoc
\section{Psychologie}

\gDefinition{}

\gLayout{0}{Einleitung}
{
%%%%%%%%%%%%%%%%%%%%%%%%%%%%%%%%%%%%%%%%%%%%%%%%%%%%%%%%%%%
% TODO Psychologie: Neuerstellung
\gTODO{GABS}{????-??-??}{Psychologie/Psychiatrie: Neuerstellung von Text
notwendig}{}
%%%%%%%%%%%%%%%%%%%%%%%%%%%%%%%%%%%%%%%%%%%%%%%%%%%%%%%%%%%
}{}{}{}{}

\subsection{Stress}
%%%%%%%%%%%%%%%%%%%%%%%%%%%%%%%%%%%%%%%%%%%%%%%%%%%%%%%%%%%
% TODO Psychologie: Neuerstellung
\gTODO{GABS}{????-??-??}{Psychologie/Psychiatrie: Stress: Neuerstellung von Text
notwendig}{}
%%%%%%%%%%%%%%%%%%%%%%%%%%%%%%%%%%%%%%%%%%%%%%%%%%%%%%%%%%%

\subsection{Coping}
%%%%%%%%%%%%%%%%%%%%%%%%%%%%%%%%%%%%%%%%%%%%%%%%%%%%%%%%%%%
% TODO Psychologie: Neuerstellung
\gTODO{GABS}{????-??-??}{Psychologie/Psychiatrie: Coping: Neuerstellung von Text
notwendig}{}
%%%%%%%%%%%%%%%%%%%%%%%%%%%%%%%%%%%%%%%%%%%%%%%%%%%%%%%%%%%

\subsection{PTS}
%%%%%%%%%%%%%%%%%%%%%%%%%%%%%%%%%%%%%%%%%%%%%%%%%%%%%%%%%%%
% TODO Psychologie: Neuerstellung
\gTODO{GABS}{????-??-??}{Psychologie/Psychiatrie: PTS: Neuerstellung von Text
notwendig}{}
%%%%%%%%%%%%%%%%%%%%%%%%%%%%%%%%%%%%%%%%%%%%%%%%%%%%%%%%%%%

\subsection{Belastungen im Dienst}
%%%%%%%%%%%%%%%%%%%%%%%%%%%%%%%%%%%%%%%%%%%%%%%%%%%%%%%%%%%
% TODO Psychologie: Neuerstellung
\gTODO{GABS}{????-??-??}{Psychologie/Psychiatrie: Neuerstellung von Text
notwendig}{}
%%%%%%%%%%%%%%%%%%%%%%%%%%%%%%%%%%%%%%%%%%%%%%%%%%%%%%%%%%%

Karutz, H.: Wenn die Belastungsgrenze erreicht ist: Psychologische Selbsthilfe
in Extremsituationen Rettungsdienst, 2009, 2009/12, 40-45 \cite{Karutz200912}

\section{Psychiatrie}
\nocite{WunnPsych1,OHCSP7}

\gMarriedLiii{Einleitung}
{
Die Psychiatrie ist ein sehr kompliziertes und
anspruchsvolles Fachgebiet der Medizin. Sie beschäftigt
sich einerseits mit handfesten, körperlich bedingten
Störungen, die Auswirkungen auf das Denken und die geistige
Leistungsfähigkeit haben. Andererseits hat man es in der
Psychiatrie sehr oft mit Störungen zu tun, von deren
Ursache die Wissenschaft kaum eine Ahnung hat. In der
Psychiatrie ist es auch besonders schwierig, zwischen
"`gesund"' und "`krank"' treffsicher zu unterscheiden. 
}{

}
{./images/wikipedia-pd/van_Gogh_Cafeterasse_bei_Nacht.jpg}{Vincent van Gogh,
Caféterasse bei Nacht.}


Grundsätzlich muss aber gesagt werden, dass psychiatrische
Erkrankungen\footnote{Sofern der Begriff
\emph{"`Erkrankung"'} überhaupt angebracht ist \ldots}
etwas alltägliches und nichts außergewöhnliches sind. Je nach Studie leiden bis zu 30\,\% der Bevölkerung an psychiatrischen Erkrankungen, allen voran an
Depressionen und der Alkoholkrankheit. 


\dictum[Otto Pötzl, nach \cite{Jolesch2}]{Der, der den
Schlüssel hat, ist der Arzt.}

\paragraph{Der "`Krankheitsbegriff"' in der Psychiatrie} Man
trifft oft auch nicht klar umrissene Krankheitsbilder, und auch deren jeweiliger Krankheitswert ist oft umstritten,
sehr vom einzelnen Patienten und seiner Situtation abhängig: Während der eine
Mensch an seiner Depression leidet, verarbeitet ein anderer seinen Weltschmerz
zu einem Lied und sichert sich dank der AKM\footnote{\gZ{Autoren, Komponisten
und Musikverleger. Die AKM ist die größte Urheberrechts- und
Verwertungsgesellschaft in Österreich.}} eine Einnahmequelle. (Subjektiv
betrachtet erscheint die Depression manchmal als ein treibender Motor der
Singer-/Songwriter-Szene.) 


\dictum[Redensart]{If there's no problem, don't fix it.}

Einer der bekanntesten psychiatrischen Patienten war der niederländische Maler
\gE{Vincent van Gogh} (\textborn\,1853\,--\,\textdied\,1890). Im Alter von 35
Jahren erkrankte er an einer psychiatrischen Erkrankung die von
Wahnvorstellungen, Albträumen, Depressionen und wiederholten
Selbstmordversuchen geprägt war. Es sollte nicht bei den Versuchen bleiben,
doch bis zu seiner Selbsttötung am 29. Juli 1890 schuf er hunderte Gemälde und
über tausend Zeichnungen. Heute gilt van Gogh als bedeutender Mitbegründer
und Impulsgeber der modernen Malerei. \cite{wikiVanGogh}

\gFazit{Die psychiatrische Diagnose verhält sich zum Patienten wie ein Barcode
zur Tafel Schokolade:

Niemand kauft eine Packung wegen des Codes.}

\begin{figure}
\begin{center}
\includegraphics[width=0.75\linewidth]{./images/wikipedia-pd/van_Gogh_Starry_Night_Over_the_Rhone.jpg}
\caption{Krank? \gSourcePicture{Vincent van Gogh: Sternennacht über der
Rhone (September 1888).}}
\end{center}
\end{figure}


Der Krankheitswert eines Zustandes beginnt dort, wo dieser Zustand für den
Patienten selbst oder für seine Umgebung ein Problem darstellt. Dabei gibt es
keine scharfe Grenze und die Grauzone ist beträchtlich.

\gLayout{0}{Psychiatrische Erkrankung in der Praxis}
{
In der Gesellschaft (repräsentiert durch die Gesetzgebung) gibt es die Übereinkunft, dass bestimmte
Krankheitsbilder bzw. Symptome als behandlungswürdig bzw.
sogar behandlungspflichtig angesehen werden. Als
behandlungspflichtig gelten zum Beispiel die \gE{Eigen-} und
\gE{Fremdgefährdung} aufgrund einer psychiatrischen Erkrankung, hier hat der
Gesetzgeber Maßnahmen nach dem \gT{Unterbringungsgesetz} (\gT{UbG}) vorgesehen. Auch der
geistige Verfall (z.B. im Rahmen einer (Alters-)Demenz) kann
rechtliche Schritte nach sich ziehen, im Sinne einer
Besachwaltung. 
}{←}{}{}{}



\gLayout{0}{Der psychiatrische Patient ist nicht entrechtet!}
{
Es ist hierbei jedoch festzuhalten, dass auch ein psychiatrischer Patient -- genauso wie jeder andere
Patient auch -- über die selben Selbstbestimmungsrechte
verfügt, d.h. es ihm grundsätzlich zusteht eine Behandlung
zu verweigern. Sollten im Rahmen der gesetzlichen
Bestimmungen (\gEs{Unterbringungsgesetz}) Maßnahmen
durchzuführen sein, so sind diese von der \gEs{Exekutive}
(Polizei) vorzunehmen. Das Ziel des Unterbringungsgesetz
ist es, das Maß der \gE{Einschränkung der Freiheit des
Patienten so gering wie möglich zu halten}. 

Auch im Rahmen der
\gEs{Besachwaltung} kann ein \gE{Gericht} das Selbstbestimmungsrecht des
Patienten einschränken. 
}{
\begin{itemize}
    \item Der psychiatrische Patient hat grundsätzlich die gleichen Rechte!
    \item Besonderheiten:
    \begin{itemize}
        \item Unterbringungsgesetz (UbG).
        \item Besachwaltung.
    \end{itemize}
\end{itemize}
}{}{}{}





\gFazit{Der Rettungsdienst führt von sich aus keine
Zwangsmaßnahmen bei psychiatrischen Patienten durch.}

Für die Arbeit im Rettungsdienst ist die Kenntnis einiger
spezieller Krankheitsbilder und Symptome notwendig, sowie
ein allgemeines Verständnis über den Umgang mit
psychiatrischen Patienten.

\subsection{Umgang mit psychiatrischen Patienten}

Wie in der Einleitung erwähnt, ist eine psychiatrische
Erkrankung etwas durchaus Alltägliches, oft stllt selbige
nicht einmal den Behandlungs- bzw. Transportgrund dar.
Ebenso wurde erwähnt, dass der psychiatrische Patient
grundsätzlich über die selben (Selbstbestimmungs-)Rechte
wie jeder andere Patient verfügt. 

\paragraph{Haltung gegenüber dem Patienten} Dementsprechend
gelten grundsätzlich auch die gleichen Regeln für die Gesprächsführung wie sie unter
\gSee{topic:kommunikationsregeln} besprochen wurden. Der
Patient verdient den gleichen Respekt und die gleiche
Wertschätzung wie ein körperlich Erkrankter.

\paragraph{Auftreten} 

\begin{itemize}
    \item Besonders wichtig ist ein \gEs{ruhiges,
    sachliches Auftreten}.
    \item  Eventuelle \gEs{Aufregung} (durch den
    Patienten; Angehörige; Feuerwehr; Polizei;
    Spezialeinsatzgruppe mit Rammbock, Tränengas,
    Maschinengewehr und Angst; \ldots) \gEs{soll nicht auf
    den Helfer abfärben}.
    \item Der Patient soll sich \gEs{nicht beengt},
    bedroht und -- wenn möglich -- nicht bevormundet
    \gEs{fühlen}.
    \item Es sollen \gEs{keine falschen oder nicht
    einlösbare Versprechen} gegeben werden.
   \item \gEs{Keine Lügen!}
\end{itemize}

\paragraph{Eigenschutz} Auch wenn vom durchschnittlichen
psychiatrischen Patienten keine höhere Gefahr als von der
Normalbevölkerung ausgeht, so ist es doch ratsam
grundlegende Punkte des Eigenschutzes zu beachten:

\begin{enumerate}
    \item Fluchtweg freihalten.
    \item Rücken freihalten. 
    \item Patienten nicht beengen/bedrängen.
    \item Patienten über Handlungen informieren und evt.
    extra um Einverständnis fragen (angurten, \ldots).
\end{enumerate}


Diese Punkte sind grundsätzlich für jeden Einsatz sinnvoll!



\subsection{Symptome}

\subsubsection{Der Patient mit \emph{Wahnvorstellungen}}

Menschen können -- aus unterschiedlichen Gründen --
Wahnvorstellungen (Halluzinationen) entwickeln. Diese
können sehr unterschiedlich sein und alle Sinne betreffen.
Man kann Dinge \gE{hören} (Stimmen hören, \ldots)  oder
Sachen \gE{sehen} ("`weiße Mäuse"', \ldots) die nicht
existieren. Der Betroffene nimmt dies jedoch als vollkommen
"`echt"' und real war. 

Der Wirklichkeitsbezug kann teilweise erhalten oder nahezu
komplett verloren gegangen sein. 



\minisec{Fallgeschichte}

\begin{quotation}
Eine 37-jährige Patientin kommt aufgrund eines
grippalen Infektes während der Wintermonate in die
Ordination eines Arztes für Allgemeinmedizin. Während des
Gespräches erzählt die Patienten, dass sie schon seit
mehreren Jahren arbeitsunfähig sei, da sie an
Halluzinationen leide. Jedes Mal wenn sie das Badezimmer
betritt, sieht sie eine ihr unbekannte Frau beim Fenster
hinein schauen, obwohl sie im zweiten Stock wohnt. Sie
erzählt, dass sie wisse, dass es nicht real sein kann,
trotzdem nimmt sie diese Frau so wahr als würde sie
tatsächlich vor dem Fenster stehen. Sie ist nun schon seit
Jahren in psychiatrischer Behandlung, allerdings stellt
sich keine Besserung ihrer Symptome ein. Die Patientin gibt
an einen massiven Leidensdruck zu haben.
\end{quotation}



\paragraph{Fazit}

Auch psychiatrische Patienten haben "`normale"'
Alltagsprobleme, wie z.B. einen grippalen Infekt. Die
Wahnvorstellungen werden von den Patienten komplett real
erlebt, auch wenn sie vielleicht sogar wissen, dass sie nicht
real sein können. 

Diese Patientin hat noch eine teilweise
intakte Realitätswahrnehmung: Psychiatrische Erkrankungen
oder Symptome funktionieren \gE{nicht} nach dem
"`Alles-oder-Nichts"'-Prinzip.


\paragraph{Wahnidee -- Wahnsysteme}

Auch die Wahrnehmung der Umgebung insgesamt kann gestört
sein, auch hier gibt es verschiedene Abstufungen. 

Die einfachste Form ist die \gT{Wahnidee}. Hierbei hat der
Patient eine klar abgegrenzte Vorstellung von etwas, was mit
der Realität nicht vereinbar ist. 

\begin{quotation}
\emph{"`Das
Gebühren-Info-Service des ORF kontrolliert jede Woche am
Verteilerschrank, ob das Antennenkabel korrekt vom
Verteiler abgesteckt wurde."'}
\end{quotation}

Diese Wahnideen können in weiterer Folge zu einem
\gT{Wahnsystem} ausgebaut werden 



\begin{quotation}
\emph{"'Das Gebühren-Info-Service
des ORF ist eine Deckmantelorganisation des israelischen
Geheimdienstes Mossad, welcher im Auftrag der CIA geheime
Abhöranlagen im ganzen Wohnhaus installiert hat. Bei
unliebsamen Äußerungen der Bewohner werden Strahlen, welche
das Denken beeinflussen können, in die Wohnungen geleitet.
Um diesen zu entgehen hat der Patient das Antennenkabel am
Verteilerkasten abgesteckt und zahlt keine Fernsehgebühren
mehr. Da ihn nun der Geheimdienst nicht mehr erreichen
kann, haben sie vom Gebühren-Info-Service einen Agenten
vorbei geschickt, welcher nun gewaltsam die Gedanken des
Patienten wegnehmen wollte."'}
\end{quotation}

\paragraph{Wahnthemen} \gZ{Ein Wahn kann unterschiedliche
Themen haben, die häufigsten sind z.B. Verfolgungswahn,
Größenwahn oder der Verarmungswahn.}

\gMassnahmenNG{Spezielle Maßnahmen}{Umgang mit Patienten mit
Wahnvorstellungen}{}{ Die -- für den Außenstehenden -- wirren Vorstellungen
stellen das Fachpersonal vor die Frage, wie mit diesen
Patienten umzugehen ist.
Ziel ist:
\begin{enumerate}
    \item Abwendung von unmittelbaren Gefahren für den
    Patienten und die Umgebung
    \item Korrekte und professionelle Durchführung eines
    Transportes in eine geeignete Einrichtung.
\end{enumerate}
Es gelten somit folgende Grundsätze:
\begin{enumerate}
    \item Sofern eine Gefährdung für den Patienten oder für
    Dritte besteht, ist diese (wenn möglich) zu beseitigen,
    dazu zählt u. U. auch der Rückzug und die Herbeiziehung
    der Exekutive.
    \item Wahnvorstellungen sollen \gEs{nicht ausgeredet}
    werden.
    \item  Wahnvorstellungen sollen aber auch \gEs{nicht
    bestätigt} werden.
    \item Am besten man nimmt die Schilderungen des
    Patienten \gEs{interessiert, aber neutral und möglichst
    kommentarlos} zur Kenntnis (\emph{"`Mhm!"', "`Aha."'}).
\end{enumerate}
}

% \subsubsection{Realitätsverlust}
% 
% \paragraph{Mögliche Ursachen}
% 
% Endogene Psychosen (Schizophrenie,  schizoaffektive Psychose)
% 
% Organische Psychosen (Alkoholparanoia,  Drogenpsychose (z.B. Coke bugs
% ),  Encephalitis)
% 
% Realitätsverlust - Wahn
% 
% Wahn = falsche, allerdings in sich  logische Auffassung der Realität,
% an der  unkorrigierbar festgehalten wird  (paranoid = wahnhaft) 
% 
% Themen: Verfolgungswahn, Größenwahn,  Verarmungswahn,...


% \subsubsection{Realitätsverlust - Halluzination}
% 
% 
% Halluzinationen: Trugwahrnehmungen,  für alle Sinne möglich
% (weiße Mäuse,  Stimmen,
% {\quotedblbase}Geisterberührungen{\textquotedblleft},...)
% 
% 
% Psychomotorische Unruhe /  Aggressivität
% 
% 
% Organische Ursachen ausschließen!!!
% 
% {\mdseries\color[rgb]{0.0,0.0,0.5019608}
% Intoxikation durch Alkohol, Drogen,  Medikamente}
% 
% 
% Hypoglykämie
% 
% St.p. SHT
% 
% Hirntumor
% 
% psychomotorische Unruhe /  Aggressivität
% 
% Schwere Erregungszustände:
% 
% Katatonie : Starre mit Sprechstereotypien, Grimassieren,
% Befehlsautomatismen, Patient imitiert Untersucher
% 
% Manie: Distanzlosigkeit, aggressiv gefärbte Agitiertheit
% 
% psychoreaktiv: abnorme Persönlichkeit
% 
% psychosoziale Krise: Reaktion auf Anlaß / Streßreaktion
% 
% psychomotorische Unruhe /  Aggressivität
% 
% Hysterie (funktionelle Erregungszustände):
% 
% anlaßgebunden, entsprechende Grundpersönlichkeit
% 
% psychosomatischer Funktionsverlust (Lähmungen, Erblindung) ohne echte
% organische Ursache
% 
% Entspricht psychischer Abwehrreaktion
% 
% \subsubsection{hysterischer Krampfanfall:}
% 
% {\mdseries
% keine Verletzungen, keine Apnoe, atypische Krämpfe, kein
% Stuhl-/Harnverlust}
% 
% {\mdseries
% unbewußt {\quotedblbase}gespielt{\textquotedblleft}}



\subsubsection{\gZ{Reserviert}\gA{Der \emph{stuporöse}
Patient}}

%\gDefinition{Stupor: Erschlaffung}
% 
% \paragraph{\gTxSymptome}
% 
% \begin{itemize}
%     \item ausdrucksloser Patient
%     \item unbeweglich
%     \item starre Mimik
%     \item spricht nicht, ißt nicht ={\textgreater} ev. bis Verhungern!
%     \item Suchterkrankungen
%     \item primär psychisches Problem
%     \item sekundär körperliche und soziale Folgen
%     \item eigengesetzlicher Ablauf
%     \item zunehmender Kontrollverlust über  eigenes Verhalten
% \end{itemize}

% \subsubsection{Manie}
% 
% {\mdseries
% rastlose Unruhe - {\quotedblbase}energiegeladen{\textquotedblleft}}
% 
% {\mdseries
% gereiztes Verhalten}
% 
% {\mdseries
% Überaktivität}
% 
% {\mdseries
% Distanzverlust}
% 
% {\mdseries
% Verwirrungszustände}
% 
% {\mdseries
% Aggression}
% 
% {\mdseries
% Selbstüberschätzung}
% 
% {\mdseries
% ev. als bipolare Störung abwechselnd mit  Depression}


\subsubsection{Psychomotorische Unruhe und Aggressivität}

Unruhe und Aggressivität kann ein Anzeichen für viele
Krankheitsbilder sein. Die Differentialdiagnosen umfassen
z.B. Blutzuckerstörungen, Demenz, Intoxikation,
Panikattacken, Manie, akute Psychosen u.v.a.m. Hier soll nur
auf die grundsätzlich zu bedenkenden Maßnahmen im Umgang mit diesen Patienten
eingegangen werden. 

\gMassnahmenNG{Spezielle Maßnahmen}{Unruhiger oder aggressiver Patient}{}{

\begin{itemize}
  \item Ruhe vermitteln und deeskalieren:
  \begin{itemize}
  \item  Ruhig bleiben! Achtung auf eigenen Tonfall!
  \item Nicht provozieren!
  \item Aufregung nicht abfärben lassen.
  \item \emph{"`Talk down"'}
  \end{itemize}
  \item Selbstschutz, Fluchtweg offen halten.
\end{itemize}

}

\subsubsection{Suizidalität}

\gDefinition{\gT{Suizid}: absichtliche
Selbsttötung. \gT{Erweiterter Suizid}: Tötung der eigenen und
fremder Personen.}

Selbstmorde und Selbstmordversuche (\gT{SMV}) gehören zu den
häufigen Einsatzgründen im Rettungsdienst. Häufig sind
Suizid- oder sonstige selbstschädliche Handlungen nicht auf
den ersten Blick als solche erkennbar. 

Wenn der Verdacht auf eine suizidale Handlung besteht, soll
dies vertraulich, aber \gEs{direkt angesprochen und geklärt
werden}. Es ist dabei wichtig, einen \gCa{neutralen Tonfall}
zu verwenden und \gCa{nicht anklagend zu klingen!}

\begin{quotation}
\emph{"`Wollten Sie sich verletzen?"'}

\emph{"`Wollten Sie sich umbringen?"'}
\end{quotation}

Häufig berichten Patienten über Gefühle der
Hoffnungslosigkeit, Einsamkeit, Verzweiflung, innerer Leere
oder lassen einen fehlenden Realitätsbezug erkennen.
\gEs{Die geäußerten Aussagen und Gefühle des Patienten sollen nicht
gewertet, verharmlost oder verniedlicht werden!}

\gCave{Jede Suizidäußerung des Patienten muss ernst
genommen und dokumentiert werden.}

Ein misslungener -- oft zögerlicher -- Selbstmordversuch
wird oft von Unwissenden im Sinne von \emph{"`Der hat es gar nicht so gemeint"'}
aufgefasst. Das ist falsch\footnote{Meistens zumindestens.}!
Ein Selbstmordversuch ist oft der letzte Hilferuf, den ein
Mensch -- sei es aus Vereinsamung oder aus anderen Gründen
-- aussenden kann. 

\gCave{Jeder Selbstmordversuch muss ernst genommen
werden, auch wenn er noch so lächerlich oder unglaubwürdig
aussieht.}


\subsubsection{Der \emph{unmittelbar eigen- oder
fremdgefährdende} Patient}

\gLayout{22}{2 Möglichkeiten:}
{


\begin{enumerate}
    \item Freiwillige Behandlung (Einwilligung des Patienten)
    \item Unterbringung (Behandlung gegen oder ohne Willen
    des Patienten) gem. Unterbringungsgesetz
\end{enumerate}



% {\mdseries
% früher: {\quotedblbase}Parere{\textquotedblleft}}


Die \index{Unterbringung}Unterbringung wird durch das
\index{Unterbringungsgesetz}Unterbringungsgesetz geregelt. Auszug:

{\S} 3. In einer Anstalt darf nur untergebracht werden, wer
\begin{itemize}
    \item 1. an einer \textbf{psychischen Krankheit} leidet und im Zusammenhang
damit \textbf{sein Leben} oder seine \textbf{Gesundheit} oder das Leben
\textbf{oder} die Gesundheit \textbf{anderer} ernstlich und erheblich
gefährdet und
\item 2. \textbf{nicht in anderer Weise}, insbesondere außerhalb einer
Anstalt, ausreichend ärztlich behandelt oder betreut werden kann.
\end{itemize}

}{}
{./images/shared/mann_mit_pistole.jpg}
{Sie kommen in eine Wohung und finden diese
Situation vor \ldots Was machen Sie?}
{Gabriel, AASS-CC-AT.at}





\paragraph{Unterbringung ohne Verlangen}~

{\S} 8. Eine Person darf gegen oder ohne ihren Willen nur dann in eine
Anstalt gebracht werden, wenn ein im \textbf{öffentlichen
Sanitätsdienst stehender Arzt} oder ein \textbf{Polizeiarzt} sie
untersucht und bescheinigt, dass die Voraussetzungen der
Unterbringung vorliegen. In der Bescheinigung sind im einzelnen die
Gründe anzuführen, aus denen der Arzt die Voraussetzungen der
Unterbringung für gegeben erachtet.

{\S} 9. (1) Die \textbf{Organe des öffentlichen Sicherheitsdienstes}
sind berechtigt und verpflichtet, eine Person, bei der sie aus
besonderen Gründen die Voraussetzungen der Unterbringung für
gegeben erachten, zur Untersuchung zum Arzt ({\S} 8) zu bringen oder
diesen beizuziehen. Bescheinigt der Arzt das Vorliegen der
Voraussetzungen der Unterbringung, so haben die Organe des
öffentlichen Sicherheitsdienstes die betroffene Person in eine
Anstalt zu bringen oder dies zu veranlassen. Wird eine solche
Bescheinigung nicht ausgestellt, so darf die betroffene Person nicht
länger angehalten werden.

(2) \textbf{Bei Gefahr im Verzug} können die \textbf{Organe des
öffentlichen Sicherheitsdienstes} die betroffene Person \textbf{auch
ohne Untersuchung und Bescheinigung in eine Anstalt bringen}.

(3) Der Arzt und die Organe des öffentlichen Sicherheitsdienstes haben
unter möglichster Schonung der betroffenen Person vorzugehen und die
notwendigen Vorkehrungen zur Abwehr von Gefahren zu treffen. Sie haben,
soweit das möglich ist, mit psychiatrischen Einrichtungen
außerhalb einer Anstalt zusammenzuarbeiten und
\textbf{erforderlichenfalls den }\textbf{örtlichen Rettungsdienst
beizuziehen}.

Daraus ergibt sich bei \gEs{fehlender}\gEs{Freiwilligkeit}:

\paragraph{Voraussetzungen für eine Unterbringung ({\S}3)}

\begin{itemize}
    \item \gEs{Psychiatrische Erkrankung}
    (ursächlich!), (${\rightarrow}$~{\S}3~Abs.~1)
    \item \gEs{Eigen-}  und/oder
    \gEs{Fremdgefährdung},
    (${\rightarrow}$~{\S}3~Abs.~1)
    \item \gEs{Keine weniger einschränkende
    Möglichkeit} zur Abwendung der Gefährdung,
    (${\rightarrow}$~{\S}3~Abs.~2)
\end{itemize}

\paragraph{Vorgehensweise ({\S}{\S}~8, 9)}

Ein {\quotedblbase}im öffentlichen Sanitätsdienst stehender Arzt
oder ein Polizeiarzt{\textquotedblleft} die Voraussetzungen zu
bescheinigen und zu Begründen. In Wien ist das der
{\quotedblbase}\index{Amtsarzt}\gEs{Amtsarzt}{\textquotedblleft},
welcher über die Polizei verständigt wird. 

\gEs{Der Notarzt, auch wenn er von der Gemeinde Wien
gestellt wird, hat nicht die Kompetenz eine Unterbringung anzuordnen}.
(${\rightarrow}$~{\S}8)

Für den Transport ist grundsätzlich die Polizei zuständig
(${\rightarrow}$~{\S}9~Abs.~1)! Die Polizei darf des Rettungsdienst
quasi als {\quotedblbase}Taxi{\textquotedblleft} beiziehen
(${\rightarrow}$~{\S}9~Abs.~3). 

Bei Gefahr im Verzug darf die Polizei auch ohne Bescheinigung eine
Person auf eine psychiatrische Abteilung zwecks Unterbringung
verbringen (${\rightarrow}$~Abs.~2). 

Im Zweifelsfall Polizei nachalarmieren...

\paragraph{Weiters: }~

\gEs{Eine Unterbringung kann nur in einer psychiatrischen Abteilung erfolgen!}

Eine Unterbringung \gEs{kann auch freiwillig}
erfolgen. Da dann auch der Transport freiwillig
stattfindet, sind dabei keine besonderen Bestimmungen zu
beachten.







\subsection{Krankheitsbilder}


\subsubsection{Demenz}

\gDefinition{Fortschreitende Einschränkung der  geistigen
Leistungsfähigkeit des Gehirns. Die häufigste Form ist die
Altersdemenz.}

\gParallel{
\paragraph{Betrifft}

\begin{itemize}
    \item Gedächtnis
    \item Denkvermögen
    \item Sprache
    \item Motorik
    \item Persönlichkeitsstruktur
\end{itemize}
}{
\paragraph{\gTxSymptome}

\begin{itemize}
    \item Zeitliche und örtliche Desorientierung
    \item psychomotorische Unruhe
    \item Personenverkennung
    \item Wechselnde Symptomatik!
\end{itemize}
}




 

% TODO: Depressionen 
% \subsubsection{\gZ{Reserviert}\gA{Depression}}
% 
% {\mdseries
% Traurigkeit}
% 
% {\mdseries
% Schuldgefühle}
% 
% {\mdseries
% fehlender Antrieb}
% 
% {\mdseries
% Schlafstörungen}
% 
% {\mdseries
% Suizidgedanken}
% 
% \paragraph{Depression / Verzweiflung / Suizidalität}~
% 
% 
% {\mdseries
% Grundstimmung entweder:}
% 
% {\mdseries
% antriebslos  oder}
% 
% {\mdseries
% agitiert (höhere Suizid gefahr!)}
% 
% \paragraph{Maßnahmen:}
% 
% {\mdseries
% auf Gegenwart konzentrieren, soziale  Umgebung miteinbeziehen, 
% Einfühlungsvermögen!!!}
% 
% {\mdseries
% Depression / Verzweiflung / Suizidalität}
% 
% {\mdseries
% Selbstmordversuch ={\textgreater} stationäre  Aufnahme, falls
% notwendig Parere!!!}
% 
% {\mdseries
% Selbstmordphantasien: je konkreter,  desto gefährlicher!}
% 
% {\mdseries
% Suizid = Selbstmord}
% 
% {\mdseries
% Achtung, evt. sogar Mitnahmeselbstmord  möglich!}
% 
% {\mdseries
% Angst}
% 
% {\mdseries
% Gefühl der Bedrohung mit vegetativer Begleitsymptomatik}
% 
% {\mdseries
% oft ohne erkennbaren Grund}
% 
% {\mdseries
% bei psychiatrischen und organischen  Krankheitsbildern}
% 
% {\mdseries
% Angst}
% 
% 
% 
% 



\subsubsection{Alkohol- und Drogenentzug}
\label{topic:alkohol-entzug}

\gParallel{
\paragraph{\gTxSymptome}

\begin{itemize}
    \item Schwitzen
    \item Unruhe
    \item Tremor  (Zittern der Hände)
    \item Delirium  (Bewusstseinsstörungen, Halluzinationen )
    \item Krampfanfälle
    \item RR-Entgleisung
    \item \gE{"`Cold turkey"'} bei Heroinentzug
\end{itemize}
}{
\paragraph{Anamnese}

\begin{itemize}
    \item Plötzlich kein Zugang mehr zu Alkohol:
    \item Kürzlicher Einzug in Altersheim
    \item kürzlich pflegebedürftig geworden
    \item \ldots
\end{itemize}
}









%%%%%%%%%%%%%%%%%%%%%%%%%%%%%%%%%%%%%%%%%%%%%%%%%%%%%%%%%%%

\paragraph{Entzugsdelir} Im Extremfall leidet der Patient
an einem Entzugsdelir (Delirium tremens). Dabei kann  es zu
Zittern, Unruhe, Angst, \gEs{Halluzinationen} ("`weiße Mäuse"', \ldots), vegetativer Entgleisung (Tachykardie,
Hypo-/Hypertonie) und zerebralen Krampfanfällen kommen. Es
handelt sich um einen \gCa{lebensbedrohlichen Zustand}!



\subsubsection{Panikattacke}

Eine Panikattacke kann sich sehr unterschiedlich äußern und
körperliche Beschwerden vortäuschen.

\begin{itemize}
    \item Plötzliches Angstgefühl
    \item Dauer: Sekunden bis Minuten
    \item Thoraxbeschwerden (Infarktangst)
    \item Hyperventilation \index{Hyperventilation!bei Panikattacke}
    \item Mehrfach durchuntersuchte  Patienten, wirken gesund
    \item Erwartungsangst
    \item Vermeidungsverhalten
    \item Antriebslosigkeit / Stupor
\end{itemize}



% TODO Schizophrenie
% \subsubsection{\gZ{Reserviert}\gA{Schizophrenie}}














% 
%  {\mdseries 10\% aller psychiatrischen Erkrankungen 
% liegt eine organische Ursache zugrunde.}
% 
% {\mdseries
% Bei Erstmanifestation einer  psychiatrischen Erkrankung MUSS  daher eine
% organische Ursache  ausgeschlossen werden!!!}
% 
% {\mdseries
% Der Umgang mit psychiatrischen Patienten ist heikel und erfordert viel
% Feingefühl sowie die Fähigkeit, auf das Gegenüber einzugehen. }




% 
% \subparagraph*{
% Situation:}
% 
% {\mdseries
% Belassung? (ist Patient überhaupt  reversfähig?)}
% 
% {\mdseries
% Unterbringung? (notwendig / rechtlich  zulässig?)}
% 
% {\mdseries
% Polizei / Amtsarzt bzw. Notarzt? (stabil?  Eskalation?)}
% 
% \subparagraph*{Leitsymptome}
% 
% \begin{itemize}
% \item Verwirrtheit, Bewußtseinseintrübung
% \item Realitätsverlust (Wahn, Halluzinationen )
% \item psychomotorische Unruhe, Aggressivität
% \item Depression , Verzweiflung, Suizidalität
% \item Angst
% \item Antriebslosigkeit, Stupor
% \item Verwirrtheit / Bewußtseinstrübung
% \item einfache Verwirrung
% \item Altersdemenz
% \item delirante Verwirrung
% \item Einfache Verwirrung
% 
% \end{itemize}
% {\mdseries
% Kann bei allen {\quotedblbase}normalen{\textquotedblleft},
% {\quotedblbase}körperlichen{\textquotedblleft} Notfällen, die mit
% Bewusstseins\-ein\-schränkungen ein\-her\-gehen, auftreten!}
% 
% {\mdseries
% Bsp.: Epilepsie , Commotio cerebri, Exsikkose etc.}
% 
% 
% delirante Verwirrung 
% organische Symptome:
% 
% Tachykardie
% 
% Mydriasis (Pupillen weit)
% 
% trockene rote Haut
% 
% Krämpfe
% 
% 
% z.B. Entzugsdelir bei Alkohol, Drogen
% 
% \subparagraph*{Maßnahmen}
% 
% 
% NAW, ad Interne Überwachungsstation (RR  instabil!)
% 
% 
% RR-Kontrollen

\chapter{[MED-THE] \emph{Spezielle Patientengruppe:} Thermische Schädigungen}
\minitoc
\section{Systemische Hitzeschäden}

Der menschliche Körper hält seine Kerntemperatur  im Normalfall
zwischen 36,5{\textdegree}  und 37,5{\textdegree} C  konstant.
Übersteigt  die Wärmezufuhr die effektive  Wärmeabgabe,
resultiert ein Hitzeschaden: \gT{Hitzekollaps} , 
\gT{Hitzeerschöpfung} , \gT{Hitzschlag}
 (bzw. bei  direkter Sonneneinstrahlung auf den Kopf ein  Sonnen\-stich
).



\subsection{Hitzekollaps}
\label{topic:hitzekollaps}

\gMarried{Bei einem Wärmestau, bei dem die zugeführte Wärme
größer ist, als die abgegebene, kann es zu einem so
genannten Hitzekollaps kommen. Die Blutgefäße in der Haut erweitern sich und das Blut versackt, der Blutdruck fällt und die
Durchblutung des Gehirns nimmt kurzfristig ab. Die Folge ist eine kurze Ohnmacht,
den Patient wird "`schwarz vor Augen"'. Lagert man den Patienten flach oder
eventuell mit erhöhten Beinen, kühlt ihn beziehungsweise
bringt ihn an einen kühlen Ort, so bessert sich sein Zustand recht rasch. } 
{\begin{itemize} \item
\textbf{Wärmestau}  (zugeführte Wärme  kurzfristig {\textgreater} abgegebene)
\item Erweiterung der Blutgefäße in der Haut  (maximal)
\item RR\,${\downarrow}$ , Durchblutung des Gehirns
kurzzeitig ${\downarrow}$
\item kurze Ohnmacht, {\quotedblbase}schwarz vor Augen{\textquotedblleft}
\end{itemize}}



\gMassnahmenNG{Spezielle Maßnahmen}{Hitzekollaps}{}{

\begin{itemize}
\item Flachlagerung: ${\rightarrow}$ Gehirndurchblutung
wieder ${\uparrow}$, bewusstseinsklar
\item nach Sekunden gebessert, {\quotedblbase}regelt sich von
selbst{\textquotedblleft}
\item Patient sitzen lassen, Flüssigkeit
\end{itemize}
}


\subsection{Hitzeerschöpfung}
\label{topic:hitzeerschoepfung}

\gMarried{Bei der Hitzeerschöpfung ist die Wärmeaufnahme
dauerhaft größer als die Wärmeabgabe. Der Körper versucht
dies durch starkes Schwitzen zu kompensieren, dadurch kommt 
es zu einem starken Wasser- und Elektrolytverlust. Der
Blutdruck fällt, es kommt zu einer Tachykardie, die
Körpertemperatur bleibt jedoch noch normal. 
Hier ist es wichtig den Patienten zu
kühlen.}{\begin{itemize}
\item Wärmeaufnahme {\textgreater}  Wärmeabgabe
\item durch starkes Schwitzen Wasser- und  Elektrolytverlust
\item starkes Durstgefühl
\item RR ${\downarrow}$, Tachykardie
\item Körpertemperatur normal
\end{itemize}}




\gMassnahmenNG{Spezielle Maßnahmen}{Hitzeerschöpfung}{}{

\begin{itemize}
\item
Kühlung, Schatten,  Flüssigkeitszufuhr, Vitalparameter,
evt. BZ messen
\end{itemize}
}



\subsection{Hitzschlag}
\label{topic:hitzschlag}

\gMarried{Beim Hitzeschlag ist nun auch die Regulation der Körpertemperatur gestört, es
sind auch sehr hohe Werte (bis 43\textcelsius ) möglich. Dies bedeutet akute
Lebensgefahr. Der Patient hat eine zunächst trockene, rote und weiße,
später grau-bläuliche Haut, der Kopf ist oft hochrot.
Weiters kommt es zu neurologischen Symptomen wie massiven
Kopfschmerzen, Schwindel, Erbrechen und
Bewusstseinsstörungen; bis hin zur Bewusstlosigkeit. Der Patient 
kann Schockzeichen zeigen, evtl.. auch eine schnelle, flache
Atmung.}{
\begin{itemize}
    \item Kopfschmerz, Schwindel, Erbrechen
    \item Haut zunächst trocken, rot und heiß,  später grau-blau, Kopf
    hochrot
    \item stark  erhöhte  Körpertemperatur (bis  43{\textdegree}C möglich
    ={\textgreater} LEBENSGEFAHR !!!)
    \item Bewusstseinsstörung, Bewusstlosigkeit
    \item Schockzeichen
    \item schnelle, flache Atmung
\end{itemize}}

\gMassnahmenNG{Spezielle Maßnahmen}{Hitzschlag}{}{

\begin{itemize}
    \item Vitale Bedrohung einschätzen. \gTxBeiVit \gTxMassVit
\item Flachlagerung an kühlem Ort (Schatten), Beine und Kopf hochlagern
\item beengende Kleidungsstücke lockern
\item Kühlung von außen (Fächer, mit Eiswürfeln abreiben etc.)
\item Kühlung von innen: für ANSPRECHBARE Patienten viel kühle
(alkoholfreie) Flüssigkeit, ev. Elektrolyt-Getränk
\item Schocklagerung u. -bekämpfung
\item Vitalparameter, \ce{O2}, BZ messen
\end{itemize}
}


\subsection{Sonnenstich}
\label{topic:sonnenstich}

\gDefinition{Hitzeschaden durch Hirnerwärmung infolge von
direkter Sonneneinstrahlung auf den \textbf{ungeschützten
Kopf}. }

\gMarried{Der Sonnenstich ist eine Sonderform des Hitzeschadens: 
er entsteht durch direkte Sonneneinstrahlung auf den
ungeschützten Kopf. Dadurch steigt die Temperatur im Kopf,
es kommt zu einer Reizung der Hirnhäute: 

Symptome wie bei
einer Hirnhautentzündung (Meningitis) beziehungsweise andere
neurologische Symptome wie zum Beispiel Kopfschmerz sind die
Folge. Die Gesichts- und Kopfhaut ist hochrot und heiß, der 
restliche Körper bleibt dagegen eher kühl. Der Patient
wirkt abgeschlagen und klagt über heftige 
Kopfschmerzen, Schwindel und Übelkeit, auch Unruhe, 
Verwirrtheitszustände und Nackensteifigkeit sind zu
beobachten. In schweren Fällen kann es sogar bis zu 
Krampfanfällen und Bewusstlosigkeit kommen.}
{
\begin{itemize}
    \item direkte  Sonneneinstrahlung auf den \textbf{ungeschützten
    Kopf}${\rightarrow}$ Temperatur  im Gehirn  steigt  ${\rightarrow}$
    Reizung der Hirnhäute  ${\rightarrow}$ Symptome wie bei
    Hirnhautentzündung (Meningitis) bzw.  neurologische Symptome durch
    Schwellung des ganzen Gehirnes
    \item Kann kombiniert mit anderen Hitzeschäden auftreten
    \item Gesichts- u. Kopfhaut hochrot u. heiß, restlicher Körper kühl
    \item Abgeschlagenheit, heftiger Kopfschmerz, Schwindel
    \item Unruhe,Verwirrtheitszustände, Brechreiz
    \item Nackensteifigkeit
    \item Schwere Fälle: Krampfanfälle, Bewusstlosigkeit
\end{itemize}
}





\gMassnahmenNG{Spezielle Maßnahmen}{Sonnenstich}{}{

\begin{itemize}
\item Vitale Bedrohung einschätzen.
    \gTxBeiVit \gTxMassVit
    \item Flachlagerung an kühlem Ort (Schatten), Kopf hochlagern
    \item beengende Kleidungsstücke lockern
    \item Kühlung von außen (Fächer, mit Eiswürfeln abreiben, kalte
    Umschläge auf die Stirn etc.)
    \item Kühlung von innen: für ANSPRECHBARE Patienten viel kühle
    (alkoholfreie) Flüssigkeit, ev. Elektrolyt-Getränk
    \item Schocklagerung u. -bekämpfung
    \item Vitalparameter, \ce{O2}, BZ messen
\end{itemize}
}



%\clearpage
\section{Unterkühlung}
\label{topic:unterkuehlung}

\gDefinition{Bei einer Körperkerntemperatur unter
36\textcelsius spricht man von Unterkühlung.}

\begin{wrapfigure}{r}{4cm}
\includegraphics[width=4cm]{./images/teaser/huette_winter_mariazell.jpg}
\end{wrapfigure}
Ähnlich
wie beim Schock kommt es auch  bei der Unterkühlung zur Zentralisation  des Blutes, um die  Körperkerntemperatur aufrecht
erhalten  zu können. Je niedriger diese  Kerntemperatur jedoch sinkt,
desto  schwerwiegender sind die Folgen!

Es besteht die Gefahr, dass kaltes Blut aus der Peripherie
zum Körper hin fließt und sich mit dem noch relativ warmen zentralisierten Blut vermischt. Es kommt
dann zu einer weiteren Unterkühlung des Körperstammes, im schlimmsstem Fall zu
einem \gT{Bergungstod} (Blutfluß in Folge von Bewegungen während des
Bergungsversuches).

\paragraph{Phasen}

\begin{itemize}
    \item \gT{Kampfphase}:  36,5{\textdegree}-34{\textcelsius},
    Kältezittern, Herzfrequenz steigt ({\textgreater}\,100\,/\,min), schnelle
    Atmung, Schmerzen
    
    \item \gT{Erschöpfungsphase}:
    34{\textdegree}-30{\textcelsius}, Bewusstseinstrübung, Kältestarre,
    Herzfrequenz sinkt wieder ({\textless}\,60\,/\,min), zu langsame Atmung,
    keine Schmerzen mehr
    
    \item \gT{Kältenarkose}: 30{\textdegree}-27{\textcelsius},
    Bewusstlosigkeit, unregelmäßiger Herzschlag, Atempausen
    
    \item {\textless}\,27\,{\textcelsius} ${\rightarrow}$ klinisch tot

\end{itemize}


\gMassnahmenNG{Spezielle Maßnahmen}{Schwere Unterkühlung
(\textless\,34\textcelsius)}{}{

\begin{itemize}
    \item \gTxMassVit
    \item Den Patient so wenig als möglich (nur soviel wie unbedingt nötig) 
   bewegen (drohender Bergungstod).
    \item Schutz vor weiterer Abkühlung, \gCa{keine \underbar{aktive}
    Erwärmung!}
    \begin{itemize}
    \item Woll- + Aludecken! Mütze!
    \end{itemize}
\end{itemize}
}



Bei Unterkühlung gilt:

\gFazit{{\quotedblbase}Nobody is dead until warm and dead.{\textquotedblleft}

Schwer unterkühlte Patienten können auch noch nach Stunden (!) ohne
Kreislauftätigkeit erfolgreich reanimiert werden!}


\chapter{[MED-TOX] \emph{Spezielle Patientengruppe:} Vergiftungen}
\minitoc
% %%%%%%%%%%%%%%%%%%%%%%%%%%%%%%%%%%%%%%%%%%%%%%%%%%%%%%%%%%%
% \section{Vergiftungen -- Toxikologie}
\label{topic:vergiftung}

\dictum[Hohenheim von Theoprast]{Die Dosis macht das Gift.}

\gDefinition{\gT{Vergiftung} bedeutet eine (meist  dosisabhängige)
organische  Funktionsbeeinträchtigung bzw. Schädigung des
Organismus durch eine aufgenommene Substanz.}

Vergiftungen sind eine interessante Sache: Ein an sich
nützlicher Stoff kann sich in eine todbringende Waffe
verwandeln, wenn er im Übermaß zugeführt wird. So ist zum
Beispiel ein Überleben ohne Salz nicht möglich. Führt man
jedoch 20\,dag zu, endet das meist tödlich. Sonst nützliche
Medikamente können schon bei einer Dosisabweichung im
Mikrogramm-Bereich eine fatale Wirkung haben.

Andererseits können Vergiftung in Abhängigkeit vom ursächlichen Stoff so
ziemlich alle Erscheinungsformen darbieten. Oft findet man
nur \gEs{unspezifische Symptome} vor, wie z.B. Erbrechen,
Durchfall, Schmerzen im Magen-Darm-Trakt,
Bewusstseinsstörungen, Krampfanfälle oder Entgleisung der
Vitalparameter.

Es ist weder möglich, noch sinnvoll, auf sämtliche mögliche
Vergiftungen im Detail einzugehen. Im Folgenden werden
einige häufig anzutreffende Vergiftungen sowie
Herangehensweisen für andere, weniger häufige, vorgestellt
und besprochen. 






\paragraph{Der Weg in den Körper -- Aufnahme des Giftes}~

Die Aufnahme der Substanz kann über verschiedene Wege
erfolgen:

\begin{description}
    \item[Oral] über den Mund 
    \item[Inhalativ] Einatmen
    \item[Stechen, spritzen] in die Vene (intravenös, zB.
    Drogenkonsum), in den Muskel, in die %(Unter-)
    Haut (z.B. Insektenstich)
    \item[Über die Haut] transkutan : fettlösliche Substanz kann  Hautbarriere durchdringen
(z.B.:  organische Lösungsmittel,  Insektenschutzmittel etc.)
\end{description}




\subparagraph*{Vorgehensweise}

{\mdseries
5-Fingerregel:}

\begin{enumerate}
\item Vitalparameter
\item Giftentfernung
\item Antidot-Therapie (Notarzt)
\item Gift-Asservierung
\item Transport
\end{enumerate}

\paragraph{Vergiftungsinformationszentrale:}

\gFazit{\gEs{{\ttfamily 01 / 406 43 43}}}

Bei der Vergiftungsinformationszentrale kann Expertenrat
eingeholt werden, sofern der verantwortliche Stoff bekannt
ist. Unter Umständen kann wird hier auch bei der
Identifizierung des Stoffes geholfen.


Gegen die meisten Vergiftungen kann man präklinisch,
ursächlich wenig unternehmen. Das Hauptaugenmerk liegt
daher auf: 

\gMassnahmenNG{Allgemeine Maßnahmen und
Herangehensweise}{Vergiftungen}{\label{m:am-vergiftung}}{
\begin{enumerate}
    \item Selbstschutz
    \item Beendigung des Giftkontaktes
    \item Evt. Dekontamination bei
    Gefahrstoffen\footnote{z.B. Insektizide, Lösungsmittel
    in der chemischen Industrie, \ldots}
    \item Sicherung der Vitalfunktionen
    \item Entscheidung, ob es sich um einen akut vital
    bedrohten Patienten handelt
    \item Quellenforschung (Mistkübel
    durchsuchen)
 \item Nüchtern lassen
    \item NICHT absichtlich Erbrechen
    herbeiführen
    \item Asservierung (Packungen, Erbrochenes in
    Nierentasse mitnehmen!)
    \item Anamnese, Besonderheiten:
    \begin{multicols}{2}
\begin{enumerate}\raggedbottom
    \item \gEs{Wer}? (Alter, KG, Vorerkrankungen, \ldots
    $\rightarrow$ Risikoabschätzung!)
    \item \gEs{Was}? (Art des Giftes $\rightarrow$ ev.
    zu erwartende Symptome)
    \item \gEs{Wann}? (Zeit seit Einnahme)
    \item \gEs{Wie}?
    (inhalativ, oral, perkutan, intravenös etc.)
    \item \gEs{Wieviel}? (Dosisabschätzung)
    \item \gEs{Warum}? (Ursachenforschung:
    Selbst\-mord\-ver\-such, Unfall, Fremdverschulden,
    Vorsatz?)
\end{enumerate}
\end{multicols}
\item ggfs. \gEs{Expertenrat einholen} und Entscheidung
bezüglich des weiteren Vorgehens

\begin{multicols}{2}
\begin{enumerate}\raggedbottom
    \item Gegenmittel verfügbar?
    \item 'normales' Spitalsbett?
    \item Intensivüberwachung?
    \item Intensivüberwachung mit Entgiftungseinheit?
    \item Dialyse?
    \item Druckkammer?
    \item Sonstige Spezialbehandlungseinheit?
\end{enumerate}
\end{multicols}
\end{enumerate}
}



\section{Vergiftungen mit Alkohol}

\gMargin{\includegraphics[width=\linewidth]{./images/TAE/dif-w2-zelt-1.jpg}}Die
Vergiftung mit Alkohol ist häufig anzutreffen, weil "`sozial akzeptiert"'. Wenn der Betroffene "`über's Ziel
hinaus schießt"', kann es schnell zu lebensbedrohlichen
Zuständen kommen. 

\gZ{Die akute Alkoholvergiftung ist von der chronischen
Alkoholabhängigkeit (Alkoholkrankheit) zu unterscheiden.}

\gLayout{0}{Wirkung}
{
Alkohol  wirkt enthemmend, z.T. stimulierend und
stimmungsmodulatorisch. Weiters kommt es zu einer Sprach- 
und Gangunsicherheit. 
}{
\begin{itemize}
    \item Enthemmend, stimmungsmodulatorisch
    \item Sprach- und Gangunsicherheit
    \item Bewusstseinstrübung
\end{itemize}
}{}{}{}

\gLayout{0}{\gTxGefahren}
{
Mit zunehmender Menge kommt es zu einem totalen
Kontrollverlust, Bewusstseinstrübung,
Schmerzunempfindlichkeit, Temperaturregulationsstörung,
Inkontinenz\footnote{\emph{Inkontinenz:} Unfähigkeit Harn
und/oder Stuhl zu halten.} Im Extremfall kann es zu einer 
vegetativen Entgleisung und einem tief komatösen Zustand 
kommen. Spätestens dann besteht akute Lebensgefahr. Weitere 
Gefahren sind Atemdepression, Aspiration, Ausfall der 
Schutzreflexe und -- besonders in der kalten Jahreszeit --
Unterkühlung.
}{
\gSlideHeading{}
\begin{itemize}
    \item Koma
    \item Atemdepression
    \item  Ausfall der
    Schutzreflexe
    \item Aspiration
    \item Schmerzunempfindlichkeit
    \item Unterkühlung
\end{itemize}
}{}{}{}


Es ist unbedingt zu erheben, ob noch andere Substanzen
(Medikamente, Drogen, \ldots) eingenommen wurden
(\gE{"`Mischintoxikation"'}). Aufgrund der oft reduzierten
Schmerzempfindlichkeit ist besonders auf Verletzungen zu
achten.

Die \gEs{vitale Bedrohung} ist vor allem anhand der
Vitalparameter, der körperlichen Untersuchung und des
Bewusstseinszustandes zu beurteilen.

\gMassnahmenNG{Spezielle Maßnahmen}{Vergiftungen mit Alkohol}{}{
\begin{itemize}
    \item Allgemeine Maßnahmen bei Vergiftungen
    (\prettyref{m:am-vergiftung})
    \item \gTxBeiVit \gTxMassVit
\end{itemize}
}


Der Alkoholentzug und das Entzugsdelir wird unter
\gSee{topic:alkohol-entzug} abgehandelt.

\gLayout{2}{Chronische Schäden}
{
\begin{itemize}
    \item Leberschäden, Zirrhose
    \item Aszites (Wasserbauch)
    \item Krampfadern in der unteren Speiseröhre  (Blutungsgefahr)
    \item Verminderte Blutgerinnung
    \item Vitaminmangel
    \item Geistiger Abbau, Gedächtnisschwäche
    \gZ{(Korsakow-Syndrom, Wernicke-Enzephalopathie)}
    \item Anämie
    \item Pankreatitis
\end{itemize}
}{

}{}{}{}













\section{Vergiftungen durch Medikamente}

\gLayout{28}{}
{
Überdosierungen von Medikamenten können komplett
unterschiedliche Symptome erzeugen, je nach eingenommenen
Medikament. Manchmal können die Symptome sogar erst Stunden
nach der Einnahme auftreten und tödlich enden. Vergiftungen
mit Medikamenten gehören daher immer abgeklärt und die
Patienten hospitalisiert. 

Neben unabsichtlichen Fehleinnahmen oder Überdosierungen
werden Arzneimittel oft auch in Selbstmordabsicht
eingenommen. Die Umstände der Einnahme müssen daher
abgeklärt werden.
}{

}{./images/teaser/imgp0482.jpg}{}{Gabriel}

\gZ{
\paragraph{Paracetamol (Mexalen\texttrademark,
APA\texttrademark, \ldots)} Paracetamol kann in hoher Dosis
zu Leberversagen führen. Symptome treten erst nach Stunden
auf. Es gibt ein Gegenmittel, dieses muss jedoch
rechtzeitig gegeben werden. Ohne Behandlung kommt es zum
Leberversagen und in Folge zum Tod. 

\paragraph{Digitalis} Digitalis ist ein auf das Herz
wirkendes Medikament und wird seit langer Zeit zur
Behandlung der Herzinsuffizienz und von
Herzrhythmusstörungen eingesetzt. Es ist jedoch sehr
schwierig die richtige Dosis für einen Patienten zu finden,
daher kommt es -- trotz korrekter Einnahme -- häufig zu
Überdosierungen. Die Folge sind Herzrhythmusstörungen und
andere EKG-Veränderungen und "`Farbensehen"'.
}

\section{Vergiftungen mit Suchtmitteln und 'Drogen'}

\subsection{'Fixer' -- Heroin \& Co.}

Heroin und Morphium gehören zu der Gruppe der Opiate und
sind im Drogenmilieu weit verbreitet, die Anwendung erfolgt
i.d.R. intravenös. 

\gLayout{0}{\gTxSymptome{} Opiatvergiftung}
{
Die gefährlichste Wirkung der Opiate ist die
Atemdepressionen, d.h. die Atemfrequenz wird dosisabhängig
immer weiter reduziert bis es schließlich zu einem
Atemstillstand kommt. Die Patienten zeigen
dementsprechend Zeichen einer Hypoxie
(Zyanose). Charakteristisch ist die starke Verengung der
Pupillen \gE{("`stecknadelkopfgroße
Pupillen"')}\index{Pupillen!stecknadelkopfgroße}. Oft
findet man alte Einstichstellen in der Armbeuge. Da in der
"`Szene"' oft unsauberes oder getauschtes Spritzbesteck
verwendet wird, besteht eine erhöhte Erkrankungsrate an HIV
und Hepatitis C. Eine weitere Gefahr geht vom Spritzbesteck
aus, wenn es nicht sicher verwahrt wird (\gCa{Vorsicht vor
der Nadel!})
}{
\begin{itemize}
    \item Atemdepression
    \item Hypoxie
    \item Stecknadelkopfgroße Pupillen
    \item Oft Begleiterkrankungen (HIV, Hep. C) 
    \item \gCa{Vorsicht Spritzbesteck!}
    \item Gegenmittel verfügbar, wirkt aber nicht so lang
\end{itemize}
}{}{}{}

\gLayout{2}{Gegenmittel}
{
Jedes Notarztmittel führt ein Opiat-Gegenmittel mit sich,
dieses ist sehr effektiv\footnote{\gZ{Naloxon. Im Handel
als: \emph{Narcanti}\texttrademark,
\emph{Naloxon-ratiopharm}\texttrademark, u.a.}}. Es ist
jedoch wichtig zu wissen, dass dieses Gegenmittel nicht so
lange wirkt wie das Opiat, d. h. es kann wieder zu einer
verminderten Atmung kommen, wenn die Wirkung des Gegenmittels nachlässt. Daher soll der Transport nur mit Arztbegleitung erfolgen.

}{

}{}{}{}


\gMassnahmenNG{Spezielle Maßnahmen}{Opiatvergiftung}{}{
\begin{itemize}
    \item Allgemeine Maßnahmen bei Vergiftungen
    (\prettyref{m:am-vergiftung})und
    symptomatische Therapie
    \item \gTxBeiVit \gTxMassVit
\end{itemize}

%Allgemeine Maßnahmen bei Vergiftungen
%(\ref{vergiftung-allgemein-massnahmen})und symptomatische
%Therapie. 
% Gegenmittel verfügbar.
}

\subsubsection{'Clubber' -- Ecstacy, Schwammerl,
Special K \ldots}

\gLayout{0}{\gTxBeschreibung}
{
\gZ{MDMA (3,4-Methylen\-dioxy\-meth\-amphet\-amin),} bekannt
unter dem Namen \gE{"`Ecstasy"' (XTC, E, X)}, ist eine der
bekanntesten \gE{"`Partydrogen"'}. Auf den ersten Blick
scheint die Substanz auch perfekt für diesen Zweck zu sein:
Der Patient erlebt ein gesteigertes Glücksgefühl, ein
Gefühl von unerschöpflicher Energie und ein gesteigertes
Erleben der Umwelt. 
}{
\begin{itemize}
    \item Ecstasy"' (XTC, E, X).
        \item Gesteigertes "`Glücksgefühl"'.
        \item Gefühl von unerschöpflicher Energie.
        \item Gesteigertes Erleben der Umwelt.
\end{itemize}
}{}{}{}

\gLayout{0}{\gTxSymptome}
{
Die Nebenwirkungen sind jedoch
gravierend und können lebensgefährlich sein. Man kann sie
unterteilen in körperliche und psychische Nebenwirkungen. 
Psychisch kommt es dosisabhängig zu Halluzinationen
und Panik-Attacken. 

Körperlich können Patienten tachykarde
Herzrhythmusstörungen zeigen, Krampfanfälle, hypertensive
Krisen, massiver Körpertemperaturanstieg über
40\textdegree{} und Muskelzerfall.

Bei einer ausgeprägten Vergiftung kann eine vitale
Bedrohung bestehen. }{
\begin{itemize}
    \item Nebenwirkungen Psychisch
    \begin{itemize}
        \item Halluzinationen
        \item Panik-Attacken
    \end{itemize}
    \item Nebenwirkungen Körperlich
    \begin{itemize}
        \item Tachykarde Herzrhythmusstörungen.
        \item Hypertensive Krisen.
        \item Krampfanfälle.
        \item Massiver Temperaturanstieg über
        {\textgreater}\,40\textdegree{}
    \end{itemize}
\end{itemize}
}{}{}{}


\gMassnahmenNG{Spezielle Maßnahmen}{Vergiftung mit
Szenedrogen}{}{
\begin{itemize}
    \item Allgemeine Maßnahmen bei Vergiftungen
    (\prettyref{m:am-vergiftung}) und
    symptomatische Therapie
    \item \gTxBeiVit \gTxMassVit
\end{itemize}

%Allgemeine Maßnahmen bei Vergiftungen
%(\ref{vergiftung-allgemein-massnahmen})und symptomatische Therapie 
}


%„Liquid Ecstasy”: Gamma-Butyrolacton-Entzugsdelir mit
%Rhabdomyolyse und dialysepflichtiger Niereninsuffizienz
\cite{Supady2009}









\section{Gase}

\subsection{Vergiftungen mit Stickgasen}

\gMassnahmenNG{Allgemeine Maßnahmen}{Stickgasvergiftungen}{\label{m:am-stickgase}}{


\begin{enumerate}
    \item Auf Selbsschutz achten!
    \item Patienten retten, aber nur wenn es der
    Eigenschutz zulässt.
    \item \ce{O2}-Berieselung mit Maske ist in keinem Fall
    ein ausreichender Atemschutz!

\item Spezialkräfte anfordern (Schwerer Atemschutz). I.d.R.
Feuerwehr
\item Weitere Opfer ausschließen (lassen):
Nachbarwohnungen, etc.
\item Sicherung der Vitalfunktionen, Beurteilung ob der
Patient vital bedroht ist. 
\item \gTxBeiVit \gTxMassVitBes
\begin{itemize}
    \item \ce{O2} hochdosiert, 15\gLMin
    \item Lagerung: situationsgerecht. 
\end{itemize}
\item Traumacheck (Folgeverletzungen)
\item Neurocheck inkl. Blutzuckerbestimmung
(Differentialdiagnosen, Beurteilung der Beeinträchtigung im Verlauf) 
\end{enumerate}}


\subsubsection{Kohlenmonoxid (\ce{CO}) -- oder: Von defekten
Heizlüftern, Gasthermen und Autoabgasen}

Kohlenmonoxid (\ce{CO}) entsteht bei unvollständiger
Verbrennung (\ce{O2}-Zufuhr ${\downarrow}$), dies passiert
vor allem bei schlecht gewarteten Gasheizungen und
Durchlauferhitzern. \ce{CO} bindet sich 300 mal besser an den
roten Blutfarbstoff Hämoglobin als $O2$, daher verdrängen
schon kleine Mengen den Sauerstoff aus dem Blut! Die Folge
ist logischerweise eine Unterversorgung mit Sauerstoff.


\gLayout{2}{Anamnese und Szene}
{
\begin{itemize}
    \item Mehrere, unerklärliche Opfer, leblose Haustiere
    \item Gasofen, Durchlauferhitzer
    \item \emph{"`Immer beim duschen!"'}
\end{itemize}
}{

}{}{}{}

\gLayout{0}{\gTxSymptome}
{
Kohlenmonoxid-beladenes Blut ist -- wie
Sauerstoff-beladenes -- hellrot. Der Patient ist
daher eher nicht bleich oder zyanotisch, sondern
die Haut ist oft rosig. Pulsoxymeter älterer Bauart
verwechseln Kohlenmonoxid-beladenes und Sauerstoff-beladenes Blut, daher wird ein falsch hoher Sauerstoffsättigungswert angezeigt!

Es stehen die Symptome des Sauerstoffmangels im Vordergrund
(Hypoxie): Der Patient leidet Atemnot, allerdings ohne
Zyanose, begleitet von Allgemeinsymptomen wie Schwindel,
Kopfschmerz, Übelkeit und Erbrechen, evt. besteht auch eine
Tachykardie. 

Es können sich rasch Bewusstseinsstörungen bis hin zur
Bewusstlosigkeit entwickeln, auch Krämpfe sind gut möglich.
}{
\begin{itemize}
    \item Dyspnoe  \gE{ohne}  Zyanose.
    \item Tachykardie.
    \item Schwindel, Kopfschmerz, Übelkeit, Erbrechen
    \item Bewusstseinsstörungen, Koma. 
    \item Krämpfe.
    \item Pulsoxymeter zeigt falsche Werte an!
\end{itemize}
}{}{}{}


\gMassnahmenNG{Spezielle Maßnahmen}{Kohlenmonoxid-Vergiftung}{}{

\begin{itemize}
    \item Allgemeine Maßnahmen bei
Stickgas-Vergiftungen
(\prettyref{m:am-stickgase})
\item \gZ{Hyperbare $O2$-Therapie in der Druckkammer erwägen}
\end{itemize}
}


% \begin{figure}[H]
%     \begin{gWideParEnv}
% \gParallel{
% \includegraphics[width=\linewidth]{./images/2006-01-31-gaswienheute2.png}
% }{
% \gCaptionof{figure}{\gAuthNotesPicLegalOk{Gas Wien Heute
% Homepage}{OK, wenn richtig titriert}\ce{CO}-Vergiftung.
% Regelmäßig kommt es zu Unfällen durch defekte Gasthermen
% oder Heizlüfter. Bei diesem Einsatz war ein RTW des ASB
% Floridsdorf-Donaustadt beteiligt. Die dritte Patientin
% wurde erst \gEs{nachdem die Feuerwehr \gE{sämtliche} (!)
% Wohnungen des Wohnhauses aufgebrochen hatte} in der
% \gE{darüberliegenden Wohnung} gefunden.
% \gSourcePicture{Wien heute \url{http://wien.orf.at}}}
% }
% \end{gWideParEnv}
% 
% \end{figure}

\begin{figure}[H]
\begin{gWideParEnv}
\begin{minipage}{0.6\linewidth}
\hspace{0.1\linewidth}
\includegraphics[width=\linewidth]{./images/2006-01-31-gaswienheute2.png}
\end{minipage}
\begin{minipage}{0.3\linewidth}
\gCaptionof{figure}{\gAuthNotesPicLegalOk{Gas Wien Heute
Homepage}{OK, wenn richtig titriert}\ce{CO}-Vergiftung.
Regelmäßig kommt es zu Unfällen durch defekte Gasthermen
oder Heizlüfter. Bei diesem Einsatz war ein RTW des ASB
Floridsdorf-Donaustadt beteiligt. Die dritte Patientin
wurde erst \gEs{nachdem die Feuerwehr \gE{sämtliche} (!)
Wohnungen des Wohnhauses aufgebrochen hatte} in der
\gE{darüberliegenden Wohnung} gefunden.
\gSourcePicture{Wien heute \url{http://wien.orf.at}}}
\end{minipage}
\end{gWideParEnv}
\end{figure}

\subsubsection{Kohlendioxid (\ce{CO2}) -- Tod im Weinkeller}


Kohlendioxid fällt als Gär- oder Verbrennungsgas gehäuft in
der Landwirtschaft und der Industrie an. Es ist schwerer als
Luft und sammelt sich in Mulden, Weinkellern, Silos, etc.


\gLayout{0}{\gTxSymptome}
{
\begin{itemize}
    \item Kopfschmerz, Schwindel, Übelkeit
    \item Krämpfe
    \item Zyanose
    \item Hyperventilation, da der hohe
\ce{CO2}-Spiegel im Blut das Atemzentrum stimuliert 
\end{itemize}
}{

}{}{}{}


\gMassnahmenNG{Spezielle Maßnahmen}{Kohlendioxid-Vergiftung}{}{
Allgemeine Maßnahmen bei
Stickgas-Vergiftungen
(\prettyref{m:am-stickgase})
}

\gCave{
Bei Gasunfällen hat der \gEs{Selbstschutz} Vorrang.
Oft ist die Bergung nur mit \gEs{schwerem
Atemschutz} möglich ${\rightarrow}$ FEUERWEHR!}

\gFazit{Ein toter Helfer ist keine Hilfe.}

\subsection{Reizgase}

\gLayout{2}{Typen}
{
\begin{description}
    \item[Soforttyp] (z.B. Tränengas,...) Augenrötung, starker Hustenreiz
    \item[Verzögerungstyp] \gZ{(z.B. Nitrosegase)} kaum
    Hustenreiz, Latenzphase (kaum Symptome), \gEs{3-24\,h
    später: toxisches Lungenödem}; bei Nichtbehandlung evt.
    irreversible Lungenschäden
\end{description}
}{

}{}{}{}

\gLayout{46}{Vorkommen}
{

}{

}{
\begin{tabulary}{\linewidth}{LL}
\gTabHeading{Vorkommen} & \gTabHeading{Art}
\\\midrule
Chemische Industrie 
&
Nitrosegase,  Schwefelwasserstoffe, Ammoniak
\\\addlinespace
Kunststoffindustrie 
&
Chlorwasserstoffe,  Formaldehyd
\\\addlinespace
Düngemittelherstellung 
&
Ammoniak,  Nitrosegase, Schwefelwasserstoffe
\\\addlinespace
Haushalt 
&
Chlorgas (Kombination von best.  Haushaltsreinigern mit
Essig), Lacke,  Imprägniersprays, ...
\\\bottomrule
\end{tabulary}
}{Vorkommen von Reizgasen.}{}





\gMassnahmenNG{Allgemeine
Maßnahmen}{Reizgasvergiftungen}{\label{m:am-reizgase}}{


\begin{enumerate}
    \item Auf Selbstschutz achten!
    \begin{itemize}
        \item Patienten retten, aber nur wenn es der
        Eigenschutz zulässt.
        \item \ce{O2}-Berieselung mit Maske ist in keinem Fall
        ein ausreichender Atemschutz!
    \end{itemize}
    \item Gefahr erkennen: Welcher Stoff? Expertenrat einholen!
    \item Wenn erforderlich: Absperrmaßnahmen durchführen
    \item Wenn erforderlich:  Spezialkräfte anfordern.
    \item Vorgehen analog zu Abschnitt "`Gefahrengutunfall"'
    (\gSee{topic:gefahrengutunfall})
    \item Weitere Opfer ausschließen (lassen):
    Nachbarwohnungen, etc.
    \item Sicherung der Vitalfunktionen, Beurteilung ob der
    Patient vital bedroht ist.
    \item \gTxBeiVit \gTxMassVitBes
    \begin{itemize}
        \item \ce{O2} hochdosiert, 15\gLMin
        \item Lagerung: situationsgerecht.
    \end{itemize}
    \item Traumacheck (Folgeverletzungen)
    \item Neurocheck inkl. Blutzuckerbestimmung
    (Differentialdiagnosen, Beurteilung der Beeinträchtigung im Verlauf)
\end{enumerate}
}




\subsubsection{Kampfgase}

\gLayout{29}{}
{
Kurz und bündig: Hier gehören Spezialisten her.
Selbstschutz und Expertenrat einholen.
}{
}{./images/gabriel-sebastian-aass/schild-atemschutzmaske-abnehmen.jpg}{}{}


\section{Wovon man nicht trinken sollte \ldots}

\subsection{Einnahme von Säuren und Laugen}
\label{topic:veraetzung-intoxikation}

Verletzungen der Haut und Augen durch Säuren und
Laugen werden Kapitel \emph{"`Traumatologische Notfälle"'}
unter \gSee{topic:veraetzung-trauma} besprochen.

\gDefinition{\gT{Verätzung}: Lokale Schädigung der
Oberfläche aufgrund einer irreversiblen Zerstörung (Denaturierung) von
Eiweißstoffen.}

% 
% Betroffene Organe
% 
% \begin{itemize}
%     \item Magen-Darm-Trakt
%     \item Haut
%     \item Augen
% \end{itemize}

\gLayout{2}{Symptome bei Einnahme}
{
\begin{itemize}
    \item Ätzspuren (Mund/Rachen), Abrinnspuren
    \item Würgen
    \item Erbrechen
    \item Schmerzen
    \item Vorsicht: Ruptur/Perforation/Blutungsgefahr!
\end{itemize}
}{

}{}{}{}




\gMassnahmenNG{Spezielle Maßnahmen}{Einnahme von Säuren oder Laugen}{}{
\begin{itemize}
    \item Allgemeine Maßnahmen bei Vergiftungen
    (\prettyref{m:am-vergiftung})
\item Schluckweise Wasser zur Verdünnung trinken lassen
\item \gEs{Niemals absichtlich Erbrechen  herbeiführen!}
\item \gZ{Keine Neutralisationsversuche! (Milch 
obsolet!)}
\end{itemize}
}


% Verätzung der Haut \gSee{topic:veraetzung}



\subsection{Schaumbildner (Wasch-/Putzmittel)}


\gLayout{2}{\gTxBeschreibung}
{
Schaumbildner vermindern die Oberflächenspannung, dadurch
werden entspannende Schaumbäder mit Tannenwaldduft möglich,
genauso wie eine effektive Geschirrwäsche. }{

}{}{}{}

\gLayout{2}{\gTxGefahren{} und \gTxKomplikationen{}}
{
\begin{itemize}
    \item Verminderte Oberflächenspannung: 
    Aspirationsgefahr ${\uparrow}$.
    \item \gEs{Verlegung der Atemwege}.
\end{itemize}
}{

}{}{}{}

\gLayout{2}{\gTxSymptome}
{
\begin{itemize}
    \item Schaum
    \item Übelkeit/Erbrechen.
\end{itemize}
}{

}{}{}{}







\gMassnahmenNG{Spezielle Maßnahmen}{Einnahme von
schaumbildenden Substanzen}{}{
\begin{itemize}
    \item Allgemeine Maßnahmen bei Vergiftungen
    (\prettyref{m:am-vergiftung})
    \item Patienten nüchtern lassen
    \item \gEs{Niemals absichtlich Erbrechen 
    herbeiführen!}
\end{itemize}
}





% 
% 
% \subsection{\gZ{Reserviert}\gA{Vergiftungen mit
% Insektiziden}}
% % TODO Vergiftungen mit Insektiziden
% 
% \subsection{\gZ{Reserviert}\gA{Vergiftungen durch giftige
% Tiere}}
% % TODO Vergiftungen durch giftige Tiere
% 
% 
% 
















































%\putbib
% Rea während Transport
\nocite{Eisenburger2008,Havel2007}
% General
\nocite{SiewertChirurgie8}
\nocite{AMLS3}
\nocite{Erc2005DeSec5}
\nocite{Erc2005Sec5}
\nocite{Erc2005DeSec6}
\nocite{Erc2005DeSec3}
\nocite{Erc2005Sec3}
\nocite{Juurlink2005}
\nocite{Kalla2006}

\nocite{Vlcek2008}
\nocite{Wells2000}
\nocite{BoehmerTaschRett6}
\nocite{Fauci2008}
\nocite{Lissauer2007}
\nocite{HaagGyn2007}
\nocite{LPN3,Boecker2}
\nocite{Helfen2008}
\nocite{LippertAnatomie7}
\nocite{Nilsson1997}
\nocite{Patzelt2005}
\nocite{Sachsenweger2003}
\nocite{RedelsteinerHRNS1}
\nocite{ForMed2}
\nocite{StriebelAn7}
\nocite{Zeiler2001}
\nocite{KlinkePhysio3}
\nocite{DRAn3}
\nocite{ForthPharma8}
\nocite{LuellmannPharma16}
\nocite{ThierbachInterhospitaltransfer1}
\nocite{AnCompact3}
\nocite{RueckerBildatlas1}


\clearpage
\gChapterBibliography 

%\bibliographystyle{unsrtdin2}{\scriptsize \bibliography{Literatur-AAS}}
